\documentclass[11pt,a4paper]{article}
\usepackage[spanish]{babel}
\usepackage[utf8]{inputenc}
\usepackage[T1]{fontenc}
\usepackage{amsmath,amssymb}
\usepackage{booktabs}
\usepackage{longtable}
\usepackage{tabularx}
\usepackage{array}
\usepackage[numbers]{natbib}
\usepackage{hyperref}
\usepackage{enumitem}
\usepackage{algorithm}
\usepackage{algpseudocode}
\usepackage{geometry}
\usepackage{tikz}
\usepackage{graphicx}
\usepackage{adjustbox}
%Includes "References" in the table of contents
\usepackage[nottoc]{tocbibind}

\usetikzlibrary{arrows.meta,positioning,fit,calc,backgrounds,shapes.geometric,shapes.multipart}
\geometry{margin=1in}
\hypersetup{colorlinks=true,linkcolor=blue,citecolor=blue,urlcolor=blue}

\tikzset{
box/.style={draw, rounded corners, align=center, minimum height=1cm, inner sep=6pt},
line/.style={-Latex, thick},
group/.style={draw, dashed, rounded corners, inner sep=8pt}
}

\newcommand{\fitfigure}[1]{\adjustbox{max width=0.98\linewidth,max totalheight=0.42\textheight,keepaspectratio}{#1}}

\title{Frankestein Transformer:\\Biblioteca Unificada de Codificador-Decodificador, CLI y Notas de Diseño Basadas en Investigación}
\author{
Erick F. Merino M. \\
Este trabajo no está afiliado \\
\texttt{erickfmm@gmail.com}
}
\date{Febrero 2026}

\begin{document}

\maketitle

\begin{abstract}
Frankestein Transformer presenta un conjunto de herramientas unificado y basado en configuración para la experimentación sistemática con arquitecturas modernas de transformers, abarcando diecisiete variantes de mezcladores de secuencia y veintitrés familias de optimizadores. El sistema admite tanto el modelado de lenguaje enmascarado (MLM) tipo codificador como la predicción autorregresiva (AR) del siguiente token tipo decodificador mediante una configuración flexible de clase de modelo y modo, con flujos de trabajo de ajuste fino especializados para ambas arquitecturas. Las contribuciones de investigación son tres: (i) un contrato de configuración estricto basado en esquemas que permite la experimentación reproducible a través de diversos mecanismos de atención, incluyendo atención softmax estándar, atención sigmoide, redes retentivas, modelos selectivos de espacio de estados, transformers de profundidad continua, enrutamiento adaptativo de profundidad (Mixture-of-Depths) \citep{zhu2026mixtureofdepthsattention}, aumento de memoria condicional (Engram) \citep{cheng2026conditionalmemoryscalablelookup}, patrones de atención dispersa y mecanismos con compuerta; (ii) un marco integral de enrutamiento de optimizadores que admite métodos de reducción de varianza (MARS, Adan, AdEMAMix), variantes eficientes en memoria (Adafactor, GaLore, Lion), enfoques sin programación de tasa de aprendizaje, precondicionadores de segundo orden (Shampoo, SOAP, Sophia) y optimizadores de bajo rango de la familia APOLLO (Apollo, Apollo-Mini, Q-Apollo) \citep{zhu2025apollosgdlikememoryadamwlevel}; y (iii) flujos de trabajo de extremo a extremo que abarcan el despliegue cuantizado mediante empaquetado ternario de pesos y el entrenamiento de incrustaciones de oraciones inspirado en SBERT, respaldados por una pila de aseguramiento de calidad ampliada con amplia cobertura de pruebas unitarias, validación de ejemplos YAML y ejecución de integración continua. El conjunto de herramientas implementa una interfaz de configuración basada en web que proporciona renderizado de formularios impulsado por esquemas con documentación en línea y validación en tiempo real. Este documento de referencia técnica incluye diagramas arquitectónicos, visualizaciones de flujo de ejecución, tablas de decisión y apéndices completos que sintetizan la literatura sobre arquitecturas de transformers, mecanismos de atención dispersa, arquitecturas con compuerta y familias de optimizadores.
\end{abstract}

\tableofcontents
\newpage

\section{Introducción}

\subsection{Motivación y Planteamiento del Problema}

Las arquitecturas de transformers han transformado fundamentalmente el panorama del modelado de secuencias en el procesamiento del lenguaje natural \citep{vaswani_attention_2017}, la visión por computadora y la biología computacional. El éxito de modelos como BERT \citep{devlin_bert_2019}, GPT y sus variantes ha establecido al transformer como el estándar de facto para el aprendizaje de representaciones en el aprendizaje profundo. Sin embargo, la rápida proliferación de innovaciones arquitectónicas presenta desafíos prácticos significativos para investigadores y profesionales.

Los últimos años han sido testigos de una explosión de mecanismos alternativos de atención y mezcla de secuencias, cada uno abordando limitaciones específicas de la atención softmax estándar: complejidad computacional cuadrática \citep{child_sparse_transformer_2019}, requisitos de inferencia eficiente en memoria \citep{sun_retentive_2023,gu_mamba_2023}, gestión selectiva de estados basada en contenido \citep{behrouz_titans_2025} y optimizaciones conscientes del hardware \citep{ramapuram_theory_2024}. Simultáneamente, la literatura sobre optimización se ha diversificado más allá del AdamW clásico, introduciendo técnicas de reducción de varianza \citep{xie_adan_2022,pagliardini_ademamix_2024,yuan_mars_2024}, variantes eficientes en memoria \citep{shazeer_adafactor_2018,zhao_galore_2024}, enfoques sin programación de tasa de aprendizaje \citep{defazio_road_2024} y métodos de precondicionamiento de segundo orden \citep{gupta_shampoo_2018,vyas_soap_2024}.

La consecuencia práctica es un ecosistema de investigación fragmentado donde la comparación experimental entre arquitecturas y optimizadores requiere un esfuerzo de ingeniería significativo. Los investigadores deben implementar y depurar múltiples variantes desde cero, garantizar tuberías de entrenamiento consistentes y gestionar espacios de hiperparámetros complejos. Esta fragmentación dificulta la reproducibilidad, ralentiza el progreso científico y aumenta la barrera de entrada para nuevos investigadores.

Este trabajo aborda estos desafíos mediante un conjunto de herramientas de experimentación unificado y basado en configuración, disponible en \url{https://github.com/erickfmm/frankestein-transformer}, que proporciona:
\begin{enumerate}[leftmargin=*]
    \item \textbf{Diseño Basado en Esquemas}: Un contrato de configuración estricto y validado que garantiza la reproducibilidad mientras admite diecisiete arquitecturas de mezcladores de secuencia distintas y veintitrés familias de optimizadores.
    \item \textbf{Entrenamiento Agnóstico a la Arquitectura}: Infraestructura de entrenamiento común que soporta líneas base de atención densa (estándar, sigmoide), alternativas recurrentes (RetNet, Mamba, bloques tipo EDO), enrutamiento adaptativo de profundidad (Mixture-of-Depths) \citep{zhu2026mixtureofdepthsattention}, bloques de memoria condicional (Engram) \citep{cheng2026conditionalmemoryscalablelookup}, patrones de atención dispersa (Sparse Transformer, Longformer, BigBird, SparseK, NSA, SpargeAttn, FASA) y mecanismos con compuerta (GLA, DeltaNet, Gated DeltaNet, Gated DeltaNet-2, HGRN2, FoX, Gated Softmax).
    \item \textbf{Marco de Enrutamiento de Optimizadores}: Grupos de hiperparámetros con prefijos que permiten un control detallado sobre incrustaciones, capas de normalización, bloques recurrentes, bloques de atención y otros subconjuntos de parámetros a través de diversos optimizadores, incluyendo la familia APOLLO (Apollo, Apollo-Mini, Q-Apollo) \citep{zhu2025apollosgdlikememoryadamwlevel}.
    \item \textbf{Flujos de Trabajo de Extremo a Extremo}: Despliegue integrado mediante cuantización (empaquetado ternario de pesos, activaciones INT8) y capacidades de incrustación de oraciones inspiradas en SBERT \citep{reimers_sentence-bert_2019}.
    \item \textbf{Configuración Interactiva}: Interfaz basada en web que proporciona renderizado de formularios impulsado por esquemas, validación en tiempo real y generación de comandos CLI.
    \item \textbf{Herramientas de Confiabilidad}: Pruebas automatizadas integrales con suites de pruebas unitarias, validación de presets YAML y automatización de CI en versiones de Python compatibles.
\end{enumerate}

La superficie de comandos del proyecto es:
\begin{center}
\texttt{frankestein-transformer}
\end{center}
con subcomandos para flujos de trabajo de codificador y decodificador:
\texttt{train}, \texttt{finetune}, \texttt{deploy}, \texttt{quantize}, \texttt{infer}, \texttt{sbert-train}, \texttt{sbert-infer}, \texttt{web-server}.

El comando \texttt{web-server} lanza un constructor de configuración basado en Streamlit que proporciona:
\begin{itemize}[leftmargin=*]
    \item Campos de formulario impulsados por esquemas con títulos de parámetros y descripciones detalladas
    \item Información sobre herramientas en tiempo real y texto de ayuda para cada opción de configuración
    \item Vista previa YAML en vivo y funcionalidad de descarga
    \item Comandos CLI generados para entrenamiento, despliegue, inferencia y flujos de trabajo SBERT
\end{itemize}

Esta interfaz interactiva sirve como alternativa a la edición manual de YAML, mejorando la usabilidad para usuarios que exploran las opciones de configuración disponibles y comprenden su impacto en el comportamiento del modelo y la dinámica de entrenamiento.

\subsection{Contribuciones}

Este trabajo realiza las siguientes contribuciones principales:
\begin{enumerate}[leftmargin=*]
    \item \textbf{Esquema de Configuración Unificado}: Un contrato de esquema basado en YAML con validación estricta que soporta diecisiete arquitecturas de mezcladores de secuencia distintas en cuatro categorías (líneas base densas, bloques recurrentes/retentivos, patrones de atención dispersa y mecanismos con compuerta), junto con controles de profundidad adaptativa y memoria condicional, mientras garantiza la reproducibilidad mediante restricciones \texttt{additionalProperties: false}.
    \item \textbf{Soporte Integral de Arquitecturas}: Implementación de variantes modernas de transformers incluyendo atención softmax estándar \citep{vaswani_attention_2017}, atención sigmoide \citep{ramapuram_theory_2024}, RetNet \citep{sun_retentive_2023}, Mamba \citep{gu_mamba_2023}, transformers continuos tipo EDO \citep{zhang_continuous_2021}, atención con aumento de memoria Titans \citep{behrouz_titans_2025}, capas de memoria condicional Engram \citep{cheng2026conditionalmemoryscalablelookup}, enrutamiento de tokens Mixture-of-Depths \citep{zhu2026mixtureofdepthsattention}, mecanismos de atención dispersa \citep{child_sparse_transformer_2019,beltagy_longformer_2020,zaheer_bigbird_2020,lou_sparsek_2024,yuan_nsa_2025,zhang_spargeattn_2025,wang_fasa_2026} y arquitecturas de atención con compuerta \citep{yang_gla_2023,yang_deltanet_2024,yang_gated_deltanet_2024,hatamizadeh_gated_deltanet2_2026,qin_hgrn2_2024,lin_forgetting_transformer_2025,qiu_gated_attention_2025}. El sistema soporta tanto entrenamiento en modo codificador con atención bidireccional para modelado de lenguaje enmascarado (MLM) como entrenamiento en modo decodificador con enmascaramiento causal para predicción autorregresiva del siguiente token.
    \item \textbf{Marco de Enrutamiento de Optimizadores}: Sistema de hiperparámetros con prefijos que permite control por grupo de parámetros a través de veintitrés optimizadores, incluyendo Anon con adaptabilidad ajustable \citep{anon2026anon}, métodos de reducción de varianza (MARS, Adan, AdEMAMix, Cautious AdamW), variantes eficientes en memoria (Adafactor, GaLore, Lion), enfoques sin programación de tasa de aprendizaje (Schedule-Free AdamW), métodos conscientes de curvatura (Sophia), precondicionadores de segundo orden (Shampoo, SOAP), optimizadores orientados a ortogonalidad (Muon, Turbo-Muon), escalado de lotes grandes (LAMB) y la familia APOLLO (Apollo, Apollo-Mini, Q-Apollo) \citep{zhu2025apollosgdlikememoryadamwlevel}.
    \item \textbf{Cuantización y Despliegue}: Tubería de despliegue integrada que soporta empaquetado ternario de pesos y cuantización de activaciones INT8 con estimaciones de tamaño siguiendo la aproximación de almacenamiento de $1.58$ bits.
    \item \textbf{Flujos de Trabajo de Incrustación de Oraciones}: Tuberías de entrenamiento e inferencia inspiradas en SBERT que soportan puntuación de similitud, recuperación, agrupamiento y exportación persistente de incrustaciones.
    \item \textbf{Interfaz de Configuración Interactiva}: Servidor web basado en Streamlit que proporciona generación de formularios impulsada por esquemas, validación en tiempo real, documentación en línea y generación de comandos CLI.
    \item \textbf{Tubería de Validación Automatizada}: Integración continua y pruebas unitarias ampliadas que verifican las rutas de entrenamiento del modelo, la integración optimizador/esquema y la compatibilidad de ejemplos YAML.
\end{enumerate}

\subsection{Guía de Lectura}

Este documento está organizado como una referencia técnica que aborda cuatro preocupaciones operativas:
\begin{enumerate}[leftmargin=*]
    \item \textbf{Requisitos Previos}: El Apéndice~\ref{sec:annex-tutorial} proporciona una introducción accesible a los transformers y mecanismos de atención para lectores nuevos en el campo.
    \item \textbf{Contrato de Configuración}: La Sección~\ref{sec:config-centric-architecture} describe el esquema YAML que garantiza experimentos válidos y la Sección~\ref{sec:schema-validation} explica las reglas de validación.
    \item \textbf{Selección de Arquitectura}: La Sección~\ref{sec:attention-families} proporciona una comparación exhaustiva de las familias de mezcladores de secuencia; los Apéndices~\ref{sec:annex-transformer}, \ref{sec:annex-latent}, \ref{sec:annex-sparse} y \ref{sec:annex-gated} sintetizan la literatura de respaldo; el Apéndice~\ref{sec:annex-norm-bitnet-mod} detalla las variantes de normalización, la cuantización BitNet, el enrutamiento Mixture-of-Depths y la memoria condicional Engram.
    \item \textbf{Dinámica de Optimización}: La Sección~\ref{sec:optimizer-families} detalla el enrutamiento de optimizadores y la dinámica de entrenamiento; el Apéndice~\ref{sec:annex-optimizers} proporciona un análisis exhaustivo de las familias de optimizadores.
    \item \textbf{Despliegue e Inferencia}: Las Secciones~\ref{sec:quantization} y~\ref{sec:sbert} describen el despliegue cuantizado y los flujos de trabajo de incrustación de oraciones.
\end{enumerate}

\subsection{Constructor de Configuración Basado en Web}

Además de la edición directa de YAML, este proyecto proporciona una interfaz web basada en Streamlit (accesible mediante el comando \texttt{web-server}) que mejora la accesibilidad y descubribilidad de la configuración. La interfaz presenta campos de esquema con:

\begin{itemize}[leftmargin=*]
    \item \textbf{Renderizado de formularios impulsado por esquemas} --- Todos los campos se generan dinámicamente a partir del esquema autoritativo, garantizando consistencia y validación.
    \item \textbf{Documentación de parámetros en línea} --- Cada campo del formulario muestra un \emph{título} del esquema como su etiqueta, con la \emph{descripción} mostrada como información sobre herramientas al pasar el cursor.
    \item \textbf{Vista previa de configuración en tiempo real} --- Los usuarios ven la salida YAML en vivo mientras modifican los campos del formulario, permitiendo retroalimentación de validación inmediata.
    \item \textbf{Generación de comandos CLI} --- La interfaz genera comandos CLI completos para entrenamiento, despliegue, inferencia y flujos de trabajo SBERT basados en la configuración actual.
    \item \textbf{Mejoras de accesibilidad} --- La información sobre herramientas y los formularios estructurados facilitan la comprensión de las opciones de configuración, especialmente para usuarios nuevos en el proyecto o que exploran arquitecturas novedosas.
\end{itemize}

Este enfoque basado en web aborda barreras de usabilidad comunes en la experimentación impulsada por configuración:
\begin{itemize}[leftmargin=*]
    \item Reduce la necesidad de memorizar la estructura YAML y los nombres de los campos
    \item Previene errores tipográficos mediante la validación del esquema
    \item Proporciona contexto educativo a través de documentación en línea
    \item Permite la experimentación rápida con ajuste guiado de parámetros
    \item Sirve tanto como herramienta de configuración como recurso de aprendizaje
\end{itemize}

La implementación de la interfaz web utiliza widgets de formulario de Streamlit con:
\begin{itemize}[leftmargin=*]
    \item \texttt{st.checkbox()} para alternativas binarias con texto de ayuda
    \item \texttt{st.number\_input()} para campos numéricos con tamaño de paso y formato
    \item \texttt{st.selectbox()} para opciones de enumeración con visualización de opciones
    \item \texttt{st.multiselect()} para selecciones de arreglo a partir de opciones definidas
    \item \texttt{st.text\_input()} para campos de cadena
    \item \texttt{st.info()}, \texttt{st.caption()} para información complementaria
\end{itemize}

Los metadatos del esquema (campos de título y descripción) se extraen y renderizan sistemáticamente en todas las secciones del formulario, incluyendo:
\begin{itemize}[leftmargin=*]
    \item Parámetros de arquitectura del modelo (tamaño oculto, capas, cabezas de atención, etc.)
    \item Configuraciones de ejecución de entrenamiento (tamaño de lote, acumulación, programador)
    \item Configuración del optimizador con hiperparámetros por grupo de parámetros
    \item Opciones de despliegue y cuantización
    \item Parámetros específicos de entrenamiento e inferencia SBERT
\end{itemize}


\section{Trabajo Relacionado}

Este trabajo se sitúa en la intersección de frameworks de aprendizaje profundo, herramientas de entrenamiento y experimentación basada en configuración. Mientras que las Secciones~\ref{sec:attention-families} y~\ref{sec:optimizer-families} examinan en profundidad la literatura arquitectónica y de optimización, esta sección se centra en el ecosistema de software para construir, entrenar y configurar modelos transformer.

\subsection{Frameworks de Aprendizaje Profundo Base}

El ecosistema moderno de investigación en transformers se apoya en varios frameworks fundamentales que proporcionan diferenciación automática, aceleración GPU y APIs de alto nivel para construcción de modelos.

\textbf{PyTorch} \citep{paszke_pytorch_2019} se ha convertido en el framework dominante para la investigación en transformers. Su modelo de programación imperativa y grafos de cómputo dinámicos permiten depuración flexible y prototipado rápido de arquitecturas novedosas. La abstracción \texttt{torch.nn.Module}, el entrenamiento automático con precisión mixta mediante \texttt{torch.cuda.amp} y el compilador JIT \texttt{torch.compile} forman la columna vertebral de la mayoría de las implementaciones contemporáneas de transformers. Frankestein Transformer está construido completamente sobre PyTorch, aprovechando su sistema de módulos para definiciones de arquitectura componibles.

\textbf{TensorFlow} \citep{abadi_tensorflow_2016} fue pionero en aprendizaje profundo orientado a producción con su paradigma de grafo de cómputo estático y la API de alto nivel Keras. Aunque su influencia en la investigación de transformers ha disminuido en relación con PyTorch, la infraestructura de serving de TensorFlow (\texttt{TF Serving}, \texttt{TF Lite}) y la biblioteca \texttt{tensorflow-text} siguen siendo relevantes para escenarios de despliegue.

\textbf{JAX} \citep{bradbury_jax_2018} ofrece un modelo de programación funcional construido alrededor de transformaciones componibles (\texttt{jit}, \texttt{vmap}, \texttt{grad}) y compilación XLA. Su diseño permite implementaciones limpias de algoritmos complejos y ha sido adoptado por varias bibliotecas recientes de transformers (por ejemplo, Flax, Haiku). El diseño sin estado de JAX y su vectorización automática son particularmente adecuados para la investigación en mecanismos de atención novedosos y algoritmos de optimización.

\textbf{scikit-learn} \citep{pedregosa_scikit-learn_2011} proporciona algoritmos clásicos de aprendizaje automático y utilidades (preprocesamiento, métricas, selección de modelos) que sirven como líneas base y herramientas complementarias en flujos de trabajo con transformers. Aunque no está diseñado para aprendizaje profundo, su API consistente y documentación extensa lo convierten en un punto de referencia estándar para comparar enfoques basados en transformers con métodos tradicionales.

\subsection{Frameworks de Entrenamiento Listos para Usar}

Varias bibliotecas proporcionan pipelines de entrenamiento completos que abstraen la complejidad del entrenamiento distribuido, la carga de datos y la gestión de hiperparámetros.

\textbf{Hugging Face Transformers} \citep{wolf_huggingface_2020} es el estándar de facto para acceder, ajustar y compartir modelos transformer preentrenados. Su API \texttt{Trainer}, el hub de modelos con más de 500,000 modelos y su estrecha integración con las bibliotecas \texttt{datasets} y \texttt{tokenizers} lo han convertido en el punto de entrada principal para NLP basado en transformers. Sin embargo, su arquitectura está optimizada para el ajuste fino de modelos preentrenados existentes en lugar de experimentar con diseños novedosos de mezcladores de secuencia o entrenar desde cero con estrategias de optimización personalizadas.

\textbf{Oumi} (\texttt{\url{https://github.com/oumi-ai/oumi}}, 2025) es una plataforma de código abierto de extremo a extremo para entrenar, ajustar, evaluar y desplegar modelos fundacionales. Proporciona una interfaz unificada a través de múltiples proveedores de nube y soporta tanto ajuste fino supervisado como optimización de preferencias (DPO, RLHF). Oumi enfatiza la preparación para producción con monitoreo integrado, checkpointing y herramientas de despliegue.

\textbf{LLaMA Factory} (\texttt{\url{https://github.com/hiyouga/LLaMA-Factory}}, 2024) ofrece un framework unificado para el ajuste fino eficiente de más de 100 modelos de lenguaje grandes. Proporciona tanto una interfaz web como una CLI, soportando LoRA, QLoRA, ajuste fino de parámetros completos y diversas técnicas de alineación. Su enfoque está en adaptar LLMs preentrenados existentes a tareas posteriores en lugar de experimentación arquitectónica.

\textbf{Axolotl} (\texttt{\url{https://github.com/axolotl-ai-cloud/axolotl}}, 2024) es una herramienta simplificada de ajuste fino que soporta LoRA, QLoRA, DeepSpeed y FSDP. Enfatiza flujos de trabajo basados en configuración mediante archivos YAML y se ha vuelto popular en la comunidad de ajuste fino de LLMs de código abierto. Al igual que LLaMA Factory, se enfoca en el ajuste fino de modelos preentrenados en lugar de entrenar arquitecturas novedosas desde cero.

\subsection{Herramientas de Bajo Código y Eficiencia}

Una categoría creciente de herramientas prioriza la accesibilidad y la eficiencia computacional, reduciendo la barrera para el entrenamiento de transformers.

\textbf{Unsloth} (\texttt{\url{https://github.com/unslothai/unsloth}}, 2024) proporciona ajuste fino de bajo código con aceleración de 2--5$\times$ y requisitos de VRAM significativamente reducidos mediante kernels CUDA escritos a mano y operaciones Triton optimizadas. Soporta exportación GGUF para despliegue cuantizado y se integra con el ecosistema Hugging Face. El enfoque de Unsloth está en hacer el ajuste fino más rápido y eficiente en memoria, no en exploración arquitectónica.

\textbf{dashAI} (\texttt{\url{https://www.dash-ai.com}}, 2024) es una aplicación de escritorio para entrenamiento de aprendizaje automático sin código con una interfaz basada en esquemas. Permite a los usuarios configurar conjuntos de datos, modelos y parámetros de entrenamiento a través de una interfaz gráfica sin escribir código. Aunque dashAI demuestra el valor de la configuración basada en esquemas para la accesibilidad, se enfoca en flujos de trabajo de ML clásico en lugar de experimentación con arquitecturas transformer.

\subsection{Experimentación Basada en Configuración}

Un hilo común entre las herramientas examinadas anteriormente es su enfoque en el ajuste fino de modelos preentrenados existentes. Hugging Face Transformers, LLaMA Factory, Axolotl y Unsloth asumen un checkpoint preentrenado como punto de partida. Oumi soporta entrenamiento desde cero pero se enfoca en arquitecturas estándar. Ninguna de estas herramientas proporciona un marco sistemático para experimentar con arquitecturas novedosas de mezcladores de secuencia, comparar diversos algoritmos de optimización bajo condiciones controladas o explorar la interacción entre elecciones arquitectónicas y dinámicas de entrenamiento.

\textbf{Frankestein Transformer} llena este vacío proporcionando un sistema de configuración basado en esquemas diseñado específicamente para la experimentación arquitectónica. Sus diferenciadores clave incluyen:

\begin{itemize}[leftmargin=*]
    \item \textbf{Más de 35 mezcladores de secuencia} que abarcan atención densa, bloques recurrentes/retentivos, patrones de atención dispersa, mecanismos con compuerta, enrutamiento adaptativo por profundidad y capas de memoria condicional, todos seleccionables mediante un único campo YAML.
    \item \textbf{23 familias de optimizadores} con un sistema de enrutamiento de hiperparámetros con prefijos que permite control por grupo de parámetros a través de embeddings, capas de normalización, bloques de atención y redes feed-forward.
    \item \textbf{Modos de entrenamiento duales} que soportan tanto modelado de lenguaje enmascarado (MLM) estilo encoder como predicción del siguiente token autorregresiva (AR) estilo decoder, con flujos de trabajo de ajuste fino especializados para cada uno.
    \item \textbf{Validación estricta de esquemas} mediante restricciones \texttt{additionalProperties: false} que previenen errores de configuración antes de que comience el entrenamiento, asegurando reproducibilidad entre experimentos.
    \item \textbf{Flujos de trabajo de extremo a extremo} que abarcan despliegue cuantizado (empaquetado de pesos ternarios, activaciones INT8) y entrenamiento de sentence embeddings inspirado en SBERT \citep{reimers_sentence-bert_2019}.
    \item \textbf{Interfaz de configuración web} que proporciona renderizado de formularios basado en esquemas con documentación en línea, validación en tiempo real y generación de comandos CLI.
\end{itemize}

Al desacoplar la definición de arquitectura de la infraestructura de entrenamiento mediante un contrato de configuración validado, Frankestein Transformer permite la comparación sistemática de familias de mezcladores de secuencia y estrategias de optimización bajo condiciones de entrenamiento idénticas---una capacidad no proporcionada por ninguna herramienta existente en el ecosistema.

\section{Diseño y Arquitectura del Sistema}
\label{sec:system-design}

\subsection{Arquitectura Centrada en Configuración}
\label{sec:config-centric-architecture}

La decisión de diseño principal en este repositorio es que la experimentación es \emph{esquema primero}. En lugar de exponer una gran cantidad de indicadores laxamente verificados, el proyecto obliga a la topología del modelo, la familia de optimizadores, los límites de entrenamiento y las opciones de telemetría a través de un único documento de configuración validado. Esto reduce la ambigüedad al reproducir resultados y permite comparar muchas arquitecturas bajo una interfaz operativa consistente.

El contrato autoritativo es \texttt{configs/schema.yaml}. Este impone tres objetos de primer nivel:
\begin{itemize}[leftmargin=*]
    \item \texttt{model\_class}
    \item \texttt{model}
    \item \texttt{training}
\end{itemize}

\paragraph{Selección de Clase de Modelo.}
El campo \texttt{model\_class} determina la variante arquitectónica instanciada por el pipeline de entrenamiento. Se soportan tres opciones:
\begin{itemize}[leftmargin=*]
    \item \textbf{frankenstein}: Modelos encoder de arquitectura mixta que soportan diversos mecanismos de atención (estándar, sigmoide, retentiva, espacio de estados, dispersa y mezcladores con compuerta) con MoE (Mezcla de Expertos) y características avanzadas. Optimizado para entrenamiento bidireccional tipo encoder con objetivos de modelado de lenguaje enmascarado (MLM).
    \item \textbf{mini}: Variante de encoder simplificada diseñada para escenarios de entrenamiento a menor escala. Proporciona una sobrecarga de parámetros reducida e iteración más rápida para experimentación y creación de prototipos.
    \item \textbf{frankesteindecoder}: Decoder causal autorregresivo para generación de siguiente token estilo LLM. Permite el enmascaramiento de atención causal para tareas de generación secuencial de texto. Cuando se selecciona esta clase, en tiempo de ejecución se fuerza \texttt{mode='decoder'}.
\end{itemize}

\paragraph{Selección de Modo de Entrenamiento.}
El campo \texttt{model.mode} controla el comportamiento de enmascaramiento de atención en todo el modelo:
\begin{itemize}[leftmargin=*]
    \item \textbf{encoder}: Utiliza atención bidireccional donde todos los tokens atienden a todos los demás tokens en la secuencia. Adecuado para tareas de preentrenamiento de modelado de lenguaje enmascarado (MLM) donde el modelo aprende a predecir tokens enmascarados aleatoriamente basándose en el contexto completo.
    \item \textbf{decoder}: Utiliza enmascaramiento causal donde cada token solo puede atender a tokens anteriores en la secuencia. Requerido para tareas autorregresivas (AR) de predicción del siguiente token, como modelado de lenguaje y generación de texto. Cuando \texttt{model\_class='frankesteindecoder'}, el sistema fuerza automáticamente \texttt{mode='decoder'} en tiempo de ejecución.
\end{itemize}

Este soporte de arquitectura dual permite al sistema manejar tanto preentrenamiento tipo encoder (MLM en contextos bidireccionales) como generación tipo decoder (predicción causal autorregresiva) a través de una interfaz de configuración unificada.

\textbf{Selección de Patrones de Mezcladores de Secuencia}: El campo \texttt{layer\_pattern} acepta una lista ordenada de identificadores de mezcladores que abarcan seis familias: líneas base densas (\texttt{standard\_attn}, \texttt{sigmoid\_attn}, \texttt{gqa}), recurrentes/retentivas (\texttt{retnet}, \texttt{retnet\_attn}, \texttt{mamba}, \texttt{ode}, \texttt{titan\_attn}), atención dispersa (\texttt{sparse\_transformer\_attn}, \texttt{longformer\_attn}, \texttt{bigbird\_attn}, \texttt{sparsek\_attn}, \texttt{nsa\_attn}, \texttt{sparge\_attn}, \texttt{fasa\_attn}, \texttt{msa\_attn}, \texttt{sparda\_attn}), mecanismos con compuerta (\texttt{gla\_attn}, \texttt{deltanet\_attn}, \texttt{gated\_deltanet\_attn}, \texttt{gated\_deltanet2\_attn}, \texttt{hgrn2\_attn}, \texttt{fox\_attn}, \texttt{gated\_softmax\_attn}, \texttt{kda\_attn}), compresión latente (\texttt{mla\_attn}, \texttt{gqla\_attn}, \texttt{mlra\_attn}, \texttt{tucker\_attn}, \texttt{iha\_attn}, \texttt{gta\_attn}, \texttt{mtla\_attn}, \texttt{cca\_attn}, \texttt{ccgqa\_attn}) y memoria condicional (\texttt{engram\_attn}). La taxonomía completa con referencias se proporciona en la Sección~\ref{sec:architecture-taxonomy} y el Apéndice~\ref{sec:annex-summary-tables}.

El \texttt{training.optimizer.optimizer\_class} soporta una amplia familia de optimizadores:
\texttt{sgd\_momentum}, \texttt{adamw}, \texttt{adafactor}, \texttt{galore\_adamw}, \texttt{prodigy}, \texttt{lion}, \texttt{sophia}, \texttt{muon}, \texttt{turbo\_muon}, \texttt{radam}, \texttt{adan}, \texttt{adopt}, \texttt{ademamix}, \texttt{mars\_adamw}, \texttt{cautious\_adamw}, \texttt{lamb}, \texttt{schedulefree\_adamw}, \texttt{shampoo}, \texttt{soap}, \texttt{anon}, \texttt{apollo}, \texttt{apollo\_mini}, y \texttt{q\_apollo}.

\begin{figure}[t]
\centering
\fitfigure{\begin{tikzpicture}[
font=\scriptsize,
box/.append style={inner sep=4pt, minimum height=0.62cm}
]
% Top level
\node[box, fill=green!10, minimum width=1.9cm] (yaml) at (0,0) {\textbf{YAML Config}};
\node[box, fill=green!8, minimum width=2.0cm] (validation) at (0,-1.55) {\textbf{Validation}\\\texttt{schema + rules}};
\node[box, fill=blue!10, minimum width=1.9cm] (web) at (-3.9,-2.95) {\textbf{Web}\\\texttt{Streamlit}};
\node[box, fill=blue!8, minimum width=2.1cm] (cli) at (3.9,-2.95) {\textbf{CLI}\\\texttt{train / deploy / quantize / infer / sbert-train / sbert-infer / transformers-export / bitnet-gguf / web-server}};
% Main blocks
\node[box, fill=orange!12, minimum width=1.9cm] (model) at (-2.6,-4.3) {\textbf{Model}\\35+ mixers};
\node[box, fill=yellow!15, minimum width=2.0cm] (train) at (0,-4.3) {\textbf{Training}\\AMP + scheduler};
\node[box, fill=red!12, minimum width=1.9cm] (opt) at (2.6,-4.3) {\textbf{Optimizer}\\23 families};
% Outcomes
\node[box, fill=purple!12, minimum width=2.0cm] (modelout) at (-2.6,-5.95) {\textbf{Model Build}\\pattern + loops};
\node[box, fill=purple!10, minimum width=2.0cm] (trainout) at (0,-5.95) {\textbf{Train Loop}\\checks + logs};
\node[box, fill=purple!10, minimum width=2.0cm] (optout) at (2.6,-5.95) {\textbf{Opt Step}\\group routing};
% Deployment
\node[box, fill=cyan!12, minimum width=2.0cm] (deploy) at (0,-7.75) {\textbf{Deploy}\\quantization};
\node[box, fill=cyan!10, minimum width=2.0cm] (sbert) at (2.6,-7.75) {\textbf{SBERT}\\search/cluster};
\node[box, fill=cyan!8, minimum width=2.0cm] (export) at (-2.6,-7.75) {\textbf{Export}\\transformers-export / bitnet-gguf};
% Arrows
\draw[line] (yaml) -- (validation);
\draw[line] (validation) -- (web);
\draw[line] (validation) -- (cli);
\draw[line] (validation) -- (model);
\draw[line] (validation) -- (train);
\draw[line] (validation) -- (opt);
\draw[line] (model) -- (modelout);
\draw[line] (train) -- (trainout);
\draw[line] (opt) -- (optout);
\draw[line] (trainout) -- (deploy);
\draw[line] (optout) -- (sbert);
\draw[line, dashed] (modelout) -- (trainout);
\draw[line, dashed] (optout.south) |- (deploy.east);
\draw[line] (cli) -- (export);
\end{tikzpicture}}
\caption{Arquitectura del sistema: la configuración fluye a través de la validación, se divide en modelo/entrenamiento/optimizador, ejecuta en tiempo real y produce artefactos de despliegue/SBERT/exportación.}
\label{fig:system-architecture}
\end{figure}

\subsection{Alcance del Esquema y Reglas de Validación}
\label{sec:schema-validation}

El esquema es estricto: los objetos de primer nivel y anidados establecen \texttt{additionalProperties: false}. Esto garantiza que las claves desconocidas fallen rápidamente en lugar de ser ignoradas silenciosamente. El objeto \texttt{training.optimizer.parameters} está adicionalmente restringido por reglas de prefijo específicas del optimizador mediante verificaciones de patrón \texttt{allOf}+\texttt{if/then}.

Los valores de normalización actualmente aceptados por el esquema son:
\[
\texttt{norm\_type} \in \{\texttt{layer\_norm}, \texttt{dynamic\_tanh}, \texttt{derf}, \texttt{rms\_norm}, \texttt{prms\_norm}\}
\]

\subsection{Esquema de Configuración}
El esquema completo de configuración del modelo y entrenamiento, incluyendo todos los campos, tipos, rangos y reglas de validación, se documenta en el Apéndice~\ref{sec:annex-schema}. El esquema impone \texttt{additionalProperties: false} en todos los niveles, garantizando que las claves desconocidas fallen rápidamente.

La profundidad en bucle inducida por el esquema es:
\[
L_{\text{logical}} = \texttt{num\_layers} \times \texttt{num\_loops}
\]
que es la definición a nivel de configuración de los bloques en bucle.

\subsection{Tipos de Tarea de Entrenamiento}

El campo \texttt{training.task} determina el objetivo de entrenamiento, trabajando en conjunto con \texttt{model.mode} para definir cómo aprende el modelo:

\paragraph{Modelado de Lenguaje Enmascarado (MLM).}
\begin{itemize}[leftmargin=*]
    \item \textbf{Modo}: Requiere \texttt{mode='encoder'} para atención bidireccional.
    \item \textbf{Objetivo}: Enmascarar aleatoriamente tokens en la secuencia de entrada (típicamente 15\%) y entrenar al modelo para predecir los tokens enmascarados basándose en el contexto bidireccional completo.
    \item \textbf{Caso de Uso}: Preentrenamiento de encoders para aprendizaje de representaciones, siguiendo la metodología BERT \citep{devlin_bert_2019}. El modelo aprende representaciones bidireccionales que capturan contexto tanto de izquierda como de derecha.
    \item \textbf{Configuración}: Utiliza el parámetro \texttt{mlm\_probability} para controlar la fracción de enmascaramiento.
\end{itemize}

\paragraph{Predicción Autorregresiva (AR) del Siguiente Token.}
\begin{itemize}[leftmargin=*]
    \item \textbf{Modo}: Requiere \texttt{mode='decoder'} para enmascaramiento causal.
    \item \textbf{Objetivo}: Entrenar al modelo para predecir el siguiente token en la secuencia dados solo los tokens anteriores.
    \item \textbf{Caso de Uso}: Generación de lenguaje y tareas estilo LLM siguiendo la metodología GPT. El modelo aprende dependencias secuenciales con atención causal donde cada token solo puede atender a tokens precedentes.
    \item \textbf{Clase de Modelo}: Típicamente usa \texttt{model\_class='frankesteindecoder'} para arquitecturas de decoder autorregresivo.
\end{itemize}

Este soporte de tareas dual permite la experimentación unificada tanto en preentrenamiento tipo encoder (MLM para comprensión bidireccional) como en generación tipo decoder (AR para producción secuencial de texto) dentro del mismo código base.

\subsection{Configuración del Optimizador}
El objeto \texttt{training.optimizer} selecciona entre 23 familias de optimizadores mediante \texttt{optimizer\_class} y proporciona control de hiperparámetros por grupo de parámetros a través de una convención de nombres con prefijos. El contrato completo de prefijos del optimizador y las familias soportadas se documentan en el Apéndice~\ref{sec:annex-schema}.

\subsection{Seguridad de Entrenamiento y Semántica de Ejecución}

Las características de seguridad a nivel de esquema incluyen acumulación, recorte, verificación de explosión post-recorte y reintentos por NaN:
\[
g_{\text{acc}}=\frac{1}{K}\sum_{i=1}^{K} g_i,\quad K=\texttt{gradient\_accumulation\_steps}
\]
\[
g_{\text{clip}} = g_{\text{acc}}\cdot
\min\left(1,\frac{\tau}{\|g_{\text{acc}}\|_2+\epsilon}\right),\quad
\tau=\texttt{grad\_clip\_max\_norm}
\]
luego los guardas de desbordamiento usan \texttt{inf\_post\_clip\_threshold} y lógica de reintentos limitada por \texttt{max\_nan\_retries}.

El pseudocódigo completo del paso de entrenamiento guiado por esquema---incluyendo acumulación de gradientes, recorte de norma global, guardas de explosión post-recorte, reintentos acotados de NaN/Inf, despacho del scheduler y rotación de checkpoints rodantes/mejores---se difiere al Apéndice~\ref{sec:annex-norm-bitnet-mod}.

\subsection{Variantes de Normalización}
\label{sec:normalization}

La normalización controla la escala de activación a través de la profundidad y es también una cuestión de compatibilidad con el esquema, ya que los valores aceptados de \texttt{norm\_type} son \texttt{layer\_norm}, \texttt{dynamic\_tanh}, \texttt{derf}, \texttt{rms\_norm} y \texttt{prms\_norm} (RMSNorm parcial). Las formulaciones matemáticas de LayerNorm, RMSNorm, pRMSNorm, Tanh Dinámico (DyT) \citep{zhu_transformers_2025} y Erf Dinámico (Derf) \citep{chen_stronger_2025}, junto con una tabla comparativa, se detallan en el Apéndice~\ref{sec:annex-norm-bitnet-mod}.


\section{Taxonomía de Arquitectura e Implementación}
\label{sec:architecture-taxonomy}

\subsection{Familias de Atención y Mezcladores de Secuencias}
\label{sec:attention-families}

Este sistema implementa diecisiete arquitecturas de mezcladores de secuencias distintas organizadas en cinco categorías funcionales que reflejan las tendencias de investigación en el diseño de modelado de secuencias. La organización taxonómica refleja la comprensión evolutiva de cómo equilibrar expresividad, eficiencia computacional y restricciones de memoria.

\begin{enumerate}[leftmargin=*]
    \item \textbf{Líneas Base de Atención Densas}: La atención softmax estándar y la atención sigmoide proporcionan contextualización global completa a costo computacional cuadrático, sirviendo como líneas base de referencia para comparación con alternativas más eficientes.
    \item \textbf{Arquitecturas Recurrentes y Retentivas}: RetNet, Mamba, bloques estilo ODE y Titans mantienen representaciones de estado que permiten un costo de inferencia $\mathcal{O}(1)$ mientras preservan la expresividad mediante dinámicas recurrentes, parámetros selectivos o adaptación de memoria en tiempo de prueba.
    \item \textbf{Patrones de Atención Dispersa}: Siete variantes dispersas (Sparse Transformer, Longformer, BigBird, SparseK, NSA, SpargeAttn, FASA) reducen la complejidad cuadrática mediante dispersión estructurada, selección de tokens o estrategias de poda sin entrenamiento.
    \item \textbf{Mecanismos de Memoria con Compuerta}: Siete arquitecturas con compuerta (GLA, DeltaNet, Gated DeltaNet, Gated DeltaNet-2, HGRN2, FoX, Gated Softmax) introducen control dependiente de datos sobre la retención de memoria, el olvido y la fuerza de actualización.
\end{enumerate}

\begin{figure}[t]
\centering
\fitfigure{\begin{tikzpicture}[
    font=\scriptsize,
    node distance=0.38cm and 1.0cm,
    box/.append style={inner sep=4pt, minimum height=0.55cm}
]
% Root
\node[box, fill=blue!15, minimum width=3.2cm] (root) {\textbf{Sequence Mixer Registry (17 variants)}};

% Category headers
\node[box, fill=red!10, below left=1.0cm and 5.4cm of root, minimum width=2.2cm] (dense) {\textbf{Dense (2)}};
\node[box, fill=green!10, below left=1.0cm and 1.8cm of root, minimum width=2.2cm] (recur) {\textbf{Recurrent (4)}};
\node[box, fill=yellow!12, below right=1.0cm and 1.8cm of root, minimum width=2.2cm] (sparse) {\textbf{Sparse (7)}};
\node[box, fill=purple!12, below right=1.0cm and 5.4cm of root, minimum width=2.2cm] (gated) {\textbf{Gated (7)}};

% Dense column
\node[box, fill=red!5, below=of dense] (d1) {\texttt{standard\_attn}};
\node[box, fill=red!5, below=of d1] (d2) {\texttt{sigmoid\_attn}};

% Recurrent column
\node[box, fill=green!5, below=of recur] (r1) {\texttt{retnet / retnet\_attn}};
\node[box, fill=green!5, below=of r1] (r2) {\texttt{mamba}};
\node[box, fill=green!5, below=of r2] (r3) {\texttt{ode}};
\node[box, fill=green!5, below=of r3] (r4) {\texttt{titan\_attn}};

% Sparse column
\node[box, fill=yellow!6, below=of sparse] (s1) {\texttt{sparse\_transformer\_attn}};
\node[box, fill=yellow!6, below=of s1] (s2) {\texttt{longformer\_attn}};
\node[box, fill=yellow!6, below=of s2] (s3) {\texttt{bigbird\_attn}};
\node[box, fill=yellow!6, below=of s3] (s4) {\texttt{sparsek\_attn}};
\node[box, fill=yellow!6, below=of s4] (s5) {\texttt{nsa\_attn}};
\node[box, fill=yellow!6, below=of s5] (s6) {\texttt{sparge\_attn}};
\node[box, fill=yellow!6, below=of s6] (s7) {\texttt{fasa\_attn}};

% Gated column
\node[box, fill=purple!6, below=of gated] (g1) {\texttt{gla\_attn}};
\node[box, fill=purple!6, below=of g1] (g2) {\texttt{deltanet\_attn}};
\node[box, fill=purple!6, below=of g2] (g3) {\texttt{gated\_deltanet\_attn}};
\node[box, fill=purple!6, below=of g3] (g4) {\texttt{gated\_deltanet2\_attn}};
\node[box, fill=purple!6, below=of g4] (g5) {\texttt{hgrn2\_attn}};
\node[box, fill=purple!6, below=of g5] (g6) {\texttt{fox\_attn}};
\node[box, fill=purple!6, below=of g6] (g7) {\texttt{gated\_softmax\_attn}};

% Root-to-category links
\draw[line] (root) -- (dense);
\draw[line] (root) -- (recur);
\draw[line] (root) -- (sparse);
\draw[line] (root) -- (gated);

% Category-to-list links
\draw[line] (dense) -- (d1);
\draw[line] (d1) -- (d2);

\draw[line] (recur) -- (r1);
\draw[line] (r1) -- (r2);
\draw[line] (r2) -- (r3);
\draw[line] (r3) -- (r4);

\draw[line] (sparse) -- (s1);
\draw[line] (s1) -- (s2);
\draw[line] (s2) -- (s3);
\draw[line] (s3) -- (s4);
\draw[line] (s4) -- (s5);
\draw[line] (s5) -- (s6);
\draw[line] (s6) -- (s7);

\draw[line] (gated) -- (g1);
\draw[line] (g1) -- (g2);
\draw[line] (g2) -- (g3);
\draw[line] (g3) -- (g4);
\draw[line] (g4) -- (g5);
\draw[line] (g5) -- (g6);
\draw[line] (g6) -- (g7);
\end{tikzpicture}}
\caption{Taxonomía completa de las diecisiete arquitecturas de mezcladores de secuencias soportadas. Las líneas base densas proporcionan enrutamiento global completo a costo cuadrático. Las arquitecturas recurrentes permiten inferencia en tiempo constante mediante compresión de estado. Las variantes dispersas reducen la complejidad mediante patrones estructurados o selección de tokens. Los mecanismos con compuerta introducen control dependiente de datos sobre la retención y el olvido de memoria.}
\label{fig:comprehensive-taxonomy}
\end{figure}

\subsection{Atención Estándar}
Dadas las matrices proyectadas $(Q,K,V)$:
\[
\text{Attn}(Q,K,V)=\text{softmax}\left(\frac{QK^\top}{\sqrt{d_k}}\right)V
\]
Este es el mecanismo base para el enrutamiento de contenido \citep{vaswani_attention_2017}.

\subsection{Atención Sigmoide}
La atención sigmoide elimina la normalización de probabilidad por fila:
\[
\text{SigmoidAttn}(Q,K,V)=\sigma\left(\frac{QK^\top}{\sqrt{d_k}}+b\right)V
\]
y tiene diferentes requisitos de estabilidad de entrenamiento, a menudo con normalización adicional \citep{ramapuram_theory_2024}.

\subsection{Formulación Retentiva}
RetNet usa retención con matriz de decaimiento $D$:
\[
\text{Retention}(Q,K,V)=\left(QK^\top \odot D\right)V
\]
con forma recurrente:
\[
S_n=\gamma S_{n-1}+k_n^\top v_n,\qquad o_n=q_n S_n
\]
permitiendo inferencia recurrente de bajo costo \citep{sun_retentive_2023,yang_survey_2025}.

\subsection{SSM Selectivo (Mamba)}
La recurrencia discreta selectiva de espacio de estados es:
\[
h_t=\bar{A}_t h_{t-1}+\bar{B}_t x_t,\qquad y_t=C_t h_t
\]
donde $(\bar{A}_t,\bar{B}_t,C_t)$ dependen de la entrada, preservando escalado en tiempo lineal con escaneo consciente del hardware \citep{gu_mamba_2023,hwang_hydra_2024}.

\subsection{Actualizaciones Continuas estilo ODE}
El marco de profundidad continua:
\[
\frac{dh(t)}{dt}=f_\theta(h(t),t)
\]
con integradores RK prácticos para ejecución discreta \citep{zhang_continuous_2021}.

\subsection{Memoria en Tiempo de Prueba (Titans)}
Una actualización aumentada con memoria puede escribirse:
\[
M_t=(1-\alpha_t)M_{t-1}+S_t,\qquad
S_t=\eta_t S_{t-1}-\theta_t\nabla \ell(M_{t-1};x_t)
\]
para adaptar la memoria en tiempo de inferencia \citep{behrouz_titans_2025,behrouz_titans_2024}.

\subsection{Extensiones Dispersas y con Compuerta en el Código Base Actual}
El registro de mezcladores ahora incluye bloques dispersos:
\texttt{sparse\_transformer\_attn}, \texttt{longformer\_attn}, \texttt{bigbird\_attn},
\texttt{sparsek\_attn}, \texttt{nsa\_attn}, \texttt{sparge\_attn}, y \texttt{fasa\_attn};
y bloques con compuerta:
\texttt{gla\_attn}, \texttt{deltanet\_attn}, \texttt{gated\_deltanet\_attn},
\texttt{gated\_deltanet2\_attn}, \texttt{hgrn2\_attn}, \texttt{fox\_attn}, y \texttt{gated\_softmax\_attn}.

La implementación impone una política de ejecución explícita para métodos dispersos sin entrenamiento:
\texttt{fasa\_attn} y \texttt{sparge\_attn} son solo para evaluación/inferencia y generan errores en tiempo de ejecución si se usan mientras el modelo está en modo de entrenamiento.

\begin{figure}[t]
\centering
\fitfigure{\begin{tikzpicture}[node distance=0.9cm and 0.9cm]
\node[box, fill=gray!10] (tokens) {Token sequence};
\node[box, fill=blue!8, below left=1.1cm and 2.5cm of tokens] (local) {Local / windowed\\Longformer};
\node[box, fill=yellow!10, below=1.1cm of tokens] (hybrid) {Hybrid sparse graph\\BigBird / Sparse Transformer};
\node[box, fill=orange!12, below right=1.1cm and 2.5cm of tokens] (select) {Selective pruning\\SparseK / NSA / FASA / SpargeAttn};
\draw[line] (tokens) -- (local);
\draw[line] (tokens) -- (hybrid);
\draw[line] (tokens) -- (select);
\node[group, fit=(local)(hybrid)(select), label=below:{\small Three sparse design strategies: locality, structured graph sparsity, and learned or predicted selection}] {};
\end{tikzpicture}}
\caption{Mapa conceptual de las estrategias de diseño de atención dispersa utilizadas en el código base. Diferentes métodos reducen el costo restringiendo vecindarios, construyendo grafos dispersos o seleccionando solo tokens/bloques de alto valor.}
\label{fig:sparse-map}
\end{figure}

\subsection{Bloques de Atención Dispersa Implementados (Detallado)}
Este código base incluye siete familias de atención dispersa \citep{child_sparse_transformer_2019,beltagy_longformer_2020,zaheer_bigbird_2020,lou_sparsek_2024,yuan_nsa_2025,zhang_spargeattn_2025,wang_fasa_2026}.

\paragraph{Sparse Transformer (\texttt{sparse\_transformer\_attn}).}
Utiliza máscaras dispersas factorizadas (con zancada + fijas) para aproximar la conectividad densa a menor costo que la atención completa $\mathcal{O}(n^2)$:
\[
\text{Attn}_i=\text{softmax}\!\left(\frac{q_iK_{A_i}^\top}{\sqrt{d_k}}\right)V_{A_i}
\]
donde $A_i$ es la vecindad dispersa inducida por reglas de zancada/fijas.
\citep{child_sparse_transformer_2019}

\paragraph{Longformer (\texttt{longformer\_attn}).}
Utiliza localidad de ventana deslizante con tokens globales opcionales:
\[
A_i=\{j:\lvert i-j\rvert \le w/2\}\cup\mathcal{G}
\]
produciendo escalado lineal en longitud de secuencia para ventana fija $w$.
\citep{beltagy_longformer_2020}

\paragraph{BigBird (\texttt{bigbird\_attn}).}
Combina ventanas locales, enlaces aleatorios y tokens globales:
\[
A_i=A_i^{\text{window}}\cup A_i^{\text{random}}\cup A_i^{\text{global}}
\]
para preservar una fuerte conectividad de contexto largo con cómputo disperso.
\citep{zaheer_bigbird_2020}

\paragraph{SparseK (\texttt{sparsek\_attn}).}
Utiliza una proyección diferenciable tipo top-$k$ sobre puntuaciones de importancia antes de la atención, de modo que solo los pares KV seleccionados participan en la ruta costosa de producto punto.
\citep{lou_sparsek_2024}

\paragraph{NSA (\texttt{nsa\_attn}).}
Implementa un diseño disperso de tres ramas: rama comprimida, rama seleccionada y rama de ventana local, y luego las combina con compuertas aprendidas:
\[
o_t=\sum_{c\in\{\text{cmp},\text{sel},\text{win}\}} g_t^c\,\text{Attn}(q_t,\tilde K_t^c,\tilde V_t^c)
\]
\citep{yuan_nsa_2025}

\paragraph{SpargeAttn (\texttt{sparge\_attn}).}
Filtrado de bloques en dos etapas sin entrenamiento: primero predice interacciones de bloques insignificantes, luego aplica poda consciente de softmax para eliminar bloques de baja contribución.
\citep{zhang_spargeattn_2025}

\paragraph{FASA (\texttt{fasa\_attn}).}
Atención sin entrenamiento consciente de frecuencia: utiliza fragmentos de frecuencia RoPE dominantes para la predicción de importancia de tokens, luego aplica atención completa solo en los tokens seleccionados.
\citep{wang_fasa_2026}

\small
\begin{longtable}{p{2.5cm}p{1.6cm}p{1.9cm}p{2.4cm}p{4.1cm}p{1.5cm}}
\toprule
\textbf{Bloque} & \textbf{Entrenable} & \textbf{Tendencia Asintótica} & \textbf{Unidad de Dispersión Principal} & \textbf{Notas de Integración Actual} & \textbf{Ref} \\
\midrule
\endhead
Sparse Transformer & Sí & sub-cuadrática & patrón de máscara (a nivel de token) & Máscaras factorizadas con zancada/fijas dentro del pipeline SDPA. & \citep{child_sparse_transformer_2019} \\
Longformer & Sí & lineal en $n$ ($w$ fijo) & ventana deslizante + tokens globales & Máscara de ventana con índices globales opcionales. & \citep{beltagy_longformer_2020} \\
BigBird & Sí & casi lineal & bordes ventana + aleatorios + globales & Máscara dispersa aleatorizada más rutas local/global. & \citep{zaheer_bigbird_2020} \\
SparseK & Sí & similar a lineal (KV seleccionados) & selección diferenciable top-$k$ de KV & Red de puntuación aprendida + proyección SparseK + atención con KV agrupados. & \citep{lou_sparsek_2024} \\
NSA & Sí & multi-rama con tokens reducidos & bloques comprimidos + bloques seleccionados + ventana local & Tres ramas dispersas combinadas con compuerta en un tensor de salida. & \citep{yuan_nsa_2025} \\
SpargeAttn & No (sin entrenamiento) & dependiente del bloque disperso & dispersión de bloque predecida a nivel de bloque & Solo evaluación en este repositorio; genera error en modo entrenamiento. & \citep{zhang_spargeattn_2025} \\
FASA & No (sin entrenamiento) & dependiente del token seleccionado & fragmentos de frecuencia dominantes + tokens seleccionados & Solo evaluación en este repositorio; genera error en modo entrenamiento. & \citep{wang_fasa_2026} \\
\bottomrule
\end{longtable}
\normalsize

\subsection{Bloques de Atención con Compuerta Implementados (Detallado)}
Este código base incluye ocho bloques con compuerta \citep{yang_gla_2023,yang_deltanet_2024,yang_gated_deltanet_2024,hatamizadeh_gated_deltanet2_2026,sun_retentive_2023,qin_hgrn2_2024,lin_forgetting_transformer_2025,qiu_gated_attention_2025}.

La idea unificadora es que las compuertas controlan \emph{qué información sobrevive}. Algunas compuertas actúan sobre actualizaciones de estado recurrente (GLA, variantes de DeltaNet, HGRN2), mientras que otras modifican la ruta de atención completa (FoX y Gated Softmax). Esto hace que las compuertas sean especialmente útiles cuando el modelo debe equilibrar recuerdo, actualidad y memoria limitada.

\begin{figure}[t]
\centering
\fitfigure{\begin{tikzpicture}[node distance=1.2cm and 1.0cm]
\node[box, fill=green!8] (state) {Previous state\\$S_{t-1}$};
\node[box, fill=red!8, right=of state] (gate) {Gate(s)\\$\alpha_t,\beta_t,G_t,f_t$};
\node[box, fill=blue!8, right=of gate] (write) {New key/value\\or SDPA output};
\node[box, fill=yellow!12, below=of gate] (next) {Updated memory / output\\$S_t$ or $O_t$};
\draw[line] (state) -- (gate);
\draw[line] (gate) -- (write);
\draw[line] (state) |- (next);
\draw[line] (write) -- (next);
\end{tikzpicture}}
\caption{Plantilla de compuerta genérica. Una compuerta puede atenuar la memoria existente, regular la fuerza de escritura o modular las salidas de atención densa, dependiendo de la familia del bloque.}
\label{fig:gating-template}
\end{figure}

\paragraph{GLA (\texttt{gla\_attn}).}
La Atención Lineal con Compuerta aplica decaimiento multiplicativo dependiente de datos en las actualizaciones de estado recurrente:
\[
S_t=G_t\odot S_{t-1}+v_tk_t^\top,\qquad o_t=S_tq_t
\]
para controlar la acumulación de memoria.
\citep{yang_gla_2023}

\paragraph{DeltaNet (\texttt{deltanet\_attn}).}
Utiliza una escritura correctora de error tipo delta-regla con fuerza de escritura aprendida $\beta_t$:
\[
S_t=S_{t-1}(I-\beta_tk_tk_t^\top)+\beta_tv_tk_t^\top
\]
lo que mejora el reemplazo de memoria dirigido.
\citep{yang_deltanet_2024}

\paragraph{Gated DeltaNet (\texttt{gated\_deltanet\_attn}).}
Añade compuerta de decaimiento $\alpha_t$ sobre las escrituras de delta-regla:
\[
S_t=\alpha_t S_{t-1}(I-\beta_tk_tk_t^\top)+\beta_tv_tk_t^\top
\]
tanto para olvido global como para actualizaciones correctivas locales.
\citep{yang_gated_deltanet_2024}

\paragraph{Gated DeltaNet-2 (\texttt{gated\_deltanet2\_attn}).}
Desacopla la compuerta delta escalar $\beta_t$ en una compuerta de borrado canalizada $\mathbf{b}_t$ (eje de claves) y una compuerta de escritura canalizada $\mathbf{w}_t$ (eje de valores), con decaimiento canalizado $\boldsymbol{\alpha}_t$ (estilo KDA):
\[
S_t=(I-k_t(\mathbf{b}_t\odot k_t)^\top)\,\mathrm{Diag}(\boldsymbol{\alpha}_t)\,S_{t-1}+k_t(\mathbf{w}_t\odot v_t)^\top
\]
para control de memoria más fino; se reduce a KDA y Gated DeltaNet como casos especiales.
\citep{hatamizadeh_gated_deltanet2_2026}

\paragraph{Alias RetNet Attn (\texttt{retnet\_attn}).}
Proporciona un alias explícito con compuerta que envuelve el comportamiento de retención multiescala para consistencia de nomenclatura en los registros de capas.
\citep{sun_retentive_2023}

\paragraph{HGRN2 (\texttt{hgrn2\_attn}).}
Utiliza compuertas de olvido con cota inferior y expansión de estado con producto exterior:
\[
S_t=\text{diag}(g_t)S_{t-1}+v_tk_t^\top
\]
para aumentar la expresividad del estado recurrente mientras se mantiene la eficiencia.
\citep{qin_hgrn2_2024}

\paragraph{FoX (\texttt{fox\_attn}).}
Inyecta sesgo de olvido a nivel de token directamente en los logits de softmax:
\[
O=\text{softmax}(QK^\top + D)V
\]
donde $D$ se deriva de compuertas acumulativas de log-olvido.
\citep{lin_forgetting_transformer_2025}

\paragraph{Gated Softmax (\texttt{gated\_softmax\_attn}).}
Aplica una compuerta sigmoide post-SDPA:
\[
Y'=\text{SDPA}(Q,K,V)\odot \sigma(XW_g)
\]
lo que añade compuerta de canal multiplicativa sin reemplazar la atención softmax.
\citep{qiu_gated_attention_2025}

\small
\begin{longtable}{p{2.4cm}p{2.0cm}p{2.2cm}p{1.8cm}p{4.0cm}p{1.5cm}}
\toprule
\textbf{Bloque} & \textbf{Tipo de Estado} & \textbf{Mecanismo de Compuerta} & \textbf{Ruta Softmax} & \textbf{Notas de Integración Actual} & \textbf{Ref} \\
\midrule
\endhead
GLA & estado recurrente matricial & decaimiento multiplicativo dependiente de datos & No (recurrente lineal) & Actualización recurrente con proyección de compuerta de bajo rango. & \citep{yang_gla_2023} \\
DeltaNet & estado recurrente matricial & compuerta de escritura ($\beta$) & No (recurrente lineal) & Corrección delta-regla con Q/K normalizados. & \citep{yang_deltanet_2024} \\
Gated DeltaNet & estado recurrente matricial & compuertas de decaimiento + escritura ($\alpha,\beta$) & No (recurrente lineal) & Olvido combinado y escritura dirigida. & \citep{yang_gated_deltanet_2024} \\
Gated DeltaNet-2 & estado recurrente matricial & compuertas de borrado + escritura decoupladas (canalizadas) & No (recurrente lineal) & Borrado canalizado (claves) + escritura (valores); decaimiento estilo KDA. & \citep{hatamizadeh_gated_deltanet2_2026} \\
RetNet Attn & estado recurrente matricial & decaimiento multiescala fijo & No (retención) & Envoltorio alias sobre el mezclador RetNet existente. & \citep{sun_retentive_2023} \\
HGRN2 & estado recurrente matricial & compuerta de olvido con cota inferior & No (recurrente lineal) & Compuerta recurrente jerárquica con productos exteriores. & \citep{qin_hgrn2_2024} \\
FoX & matriz de atención completa & compuerta de olvido en espacio de logits & Sí & Sesgo de olvido añadido antes de softmax. & \citep{lin_forgetting_transformer_2025} \\
Gated Softmax & matriz de atención completa & compuerta sigmoide post-atención & Sí & Compuerta sigmoide aplicada después de la salida SDPA. & \citep{qiu_gated_attention_2025} \\
\bottomrule
\end{longtable}
\normalsize

\begin{algorithm}[t]
\caption{Paso Adelante del Mezclador Guiado por Patrón (Conceptual)}
\begin{algorithmic}[1]
\Require Estados ocultos $H$, patrón $P$, índice de capa $\ell$
\State $m \gets P[\ell \bmod |P|]$
\If{$m \in \{\texttt{fasa\_attn}, \texttt{sparge\_attn}\}$ y el modelo está en modo entrenamiento}
    \State lanzar error de configuración/tiempo de ejecución (bloque sin entrenamiento en modo entrenar)
\ElsIf{$m=\texttt{standard\_attn}$}
    \State $H \gets \text{softmax-attention}(H)$
\ElsIf{$m=\texttt{sigmoid\_attn}$}
    \State $H \gets \text{sigmoid-attention}(H)$
\ElsIf{$m\in\{\texttt{retnet},\texttt{retnet\_attn}\}$}
    \State $H \gets \text{retention}(H)$
\ElsIf{$m=\texttt{mamba}$}
    \State $H \gets \text{selective-ssm}(H)$
\ElsIf{$m=\texttt{ode}$}
    \State $H \gets \text{rk-step}(H)$
\ElsIf{$m$ es una clave de atención dispersa}
    \State $H \gets \text{sparse-attention-family}(H)$
\ElsIf{$m$ es una clave de atención con compuerta}
    \State $H \gets \text{gated-attention-family}(H)$
\Else
    \State $H \gets \text{memory-augmented-attn}(H)$
\EndIf
\State \Return $H$
\end{algorithmic}
\end{algorithm}


\section{Familias de Optimizadores y Dinámicas de Entrenamiento}
\label{sec:optimizer-families}

La optimización de arquitecturas transformer altamente parametrizadas presenta desafíos significativos debido a paisajes de pérdida no convexos, puntos de silla y heterogeneidad de bloques entre grupos de parámetros. Este sistema aborda estos desafíos mediante un marco unificado que soporta veintitrés familias de optimizadores que abarcan seis categorías algorítmicas: (1) líneas base clásicas (SGD, AdamW), (2) momento avanzado y reducción de varianza (Adan, ADOPT, AdEMAMix, MARS, Cautious), (3) variantes eficientes en memoria (Adafactor, GaLore, Lion, APOLLO, APOLLO-Mini, Q-APOLLO), (4) métodos sin programación y sin parámetros (Schedule-Free AdamW, Prodigy) más escalado de lotes grandes mediante LAMB, (5) conscientes de curvatura y segundo orden (Sophia, Shampoo, SOAP), (6) orientados a geometría (Muon, Turbo-Muon), y (7) optimizadores con adaptabilidad ajustable (Anon). Las descripciones algorítmicas detalladas, pseudocódigo y una tabla comparativa de memoria/complejidad para todos los optimizadores se proporcionan en el Apéndice~\ref{sec:annex-optimizers}.

\begin{figure}[t]
\centering
\fitfigure{\begin{tikzpicture}[node distance=1.1cm and 1.0cm]
\node[box, fill=blue!8] (start) {Elegir objetivo de optimizador};
\node[box, fill=green!8, below left=of start, xshift=-1.5cm] (baseline) {Línea base confiable\\AdamW / RAdam};
\node[box, fill=yellow!12, below=of start] (memory) {Menor memoria de optimizador\\Adafactor / GaLore / Lion};
\node[box, fill=orange!12, below right=of start, xshift=1.5cm] (advanced) {Agresivo o estructurado\\Adan, ADOPT, Sophia,\\Shampoo, SOAP, Muon};
\node[box, fill=red!8, below=of memory] (route) {Grupos con prefijo de esquema\\LR, decay de peso, betas, eps};
\draw[line] (start) -- (baseline);
\draw[line] (start) -- (memory);
\draw[line] (start) -- (advanced);
\draw[line] (baseline) -- (route);
\draw[line] (memory) -- (route);
\draw[line] (advanced) -- (route);
\end{tikzpicture}}
\caption{Selección de optimizador en la práctica: la clase determina la regla de actualización, mientras que el esquema controla cómo se aplican los hiperparámetros entre embeddings, normas, bloques recurrentes, bloques de atención y otros parámetros.}
\label{fig:optimizer-routing}
\end{figure}


\section{Cuantización y Despliegue}
\label{sec:quantization}

El stack de despliegue utiliza empaquetado ternario de pesos más cuantización de activaciones INT8 para producir artefactos eficientes.

\begin{figure}[t]
\centering
\fitfigure{\begin{tikzpicture}[node distance=1.3cm and 0.9cm]
\node[box, fill=blue!8] (ckpt) {Checkpoint entrenado};
\node[box, fill=yellow!10, right=of ckpt] (quant) {Empaquetado de pesos\\ternario / bajo bit};
\node[box, fill=orange!12, right=of quant] (act) {Escalado de activaciones\\ruta INT8};
\node[box, fill=green!8, right=of act] (artifact) {Artefacto desplegable\\huella de almacenamiento reducida};
\draw[line] (ckpt) -- (quant);
\draw[line] (quant) -- (act);
\draw[line] (act) -- (artifact);
\end{tikzpicture}}
\caption{Ruta de despliegue desde un checkpoint entrenado hasta un artefacto compacto. El código base trata la cuantización como una transformación de la etapa de despliegue en lugar de una familia de modelos separada.}
\label{fig:deploy-flow}
\end{figure}

\subsection{Cuantización Ternaria y Escalado de Activaciones}
La etapa de despliegue transforma un checkpoint entrenado en un artefacto compacto mediante empaquetado ternario de pesos estilo BitNet \citep{wang_bitnet_2023}, escalado de activaciones INT8 y aproximación de almacenamiento de 1.58 bits alineada con objetivos de despliegue ligero \citep{samson_lightweight_2026,bravin_embbert_2026}. Las formulaciones matemáticas del escalado ternario de pesos, la cuantización/des-cuantización de activaciones INT8 y las estimaciones de tamaño FP32/FP16/1.58 bits se difieren al Apéndice~\ref{sec:annex-norm-bitnet-mod}.
\section{Tareas Posteriores de SBERT}
\label{sec:sbert}

El embedding de oraciones se construye sobre entrenamiento estilo Siamese \citep{reimers_sentence-bert_2019}. Para el par de oraciones $(s_1,s_2)$ con embeddings $(e_1,e_2)$:
\[
\text{cos}(e_1,e_2)=\frac{e_1^\top e_2}{\|e_1\|\|e_2\|}
\]
y pérdida coseno estilo regresión:
\[
\mathcal{L}_{\text{cos}}=\left(\text{cos}(e_1,e_2)-y\right)^2
\]
con $y\in[-1,1]$ en este pipeline.

Modos posteriores soportados:
\begin{itemize}[leftmargin=*]
    \item \textbf{Similitud}: puntuación por pares entre dos oraciones.
    \item \textbf{Búsqueda}: $k$ vecinos más cercanos sobre un corpus.
    \item \textbf{Clustering}: agrupación de embeddings (ej., k-means).
    \item \textbf{Codificación}: exportación persistente de embeddings para recuperación posterior.
\end{itemize}

\begin{figure}[t]
\centering
\fitfigure{\begin{tikzpicture}[node distance=1.1cm and 1.0cm]
\node[box, fill=blue!8] (encoder) {Codificador compartido};
\node[box, fill=yellow!10, below left=of encoder, xshift=-1.2cm] (sim) {Similitud\\puntuación coseno};
\node[box, fill=green!8, below=of encoder] (search) {Búsqueda\\recuperación top-$k$};
\node[box, fill=orange!12, below right=of encoder, xshift=1.2cm] (cluster) {Cluster / codificar\\análisis offline};
\draw[line] (encoder) -- (sim);
\draw[line] (encoder) -- (search);
\draw[line] (encoder) -- (cluster);
\end{tikzpicture}}
\caption{Reutilización del flujo de trabajo SBERT. Un único codificador soporta puntuación de pares en línea, recuperación sobre corpus, clustering y exportación persistente de embeddings.}
\label{fig:sbert-modes}
\end{figure}

\begin{algorithm}[t]
\caption{Enrutador de Modo de Inferencia SBERT}
\begin{algorithmic}[1]
\Require modo $m$, modelo $E$, entradas $X$
\If{$m=\texttt{similarity}$}
    \State retornar $\cos(E(x_1),E(x_2))$
\ElsIf{$m=\texttt{search}$}
    \State retornar top-$k$ por producto-punto/coseno contra embeddings del corpus
\ElsIf{$m=\texttt{cluster}$}
    \State retornar etiquetas de clustering sobre $E(X)$
\Else
    \State retornar embeddings serializados $E(X)$
\EndIf
\end{algorithmic}
\end{algorithm}


\section{Tablas Resumen}
\label{sec:summary-tables}

\subsection{Resumen de Atención y Mezcladores de Secuencia}
\small
\begin{longtable}{p{2.4cm}p{3.2cm}p{2.0cm}p{2.0cm}p{4.8cm}}
\toprule
\textbf{Tipo} & \textbf{Ecuación Principal} & \textbf{Entrenar} & \textbf{Inferir} & \textbf{Notas} \\
\midrule
\endhead
Atención Estándar & $\text{softmax}(QK^\top/\sqrt{d_k})V$ & $\mathcal{O}(n^2d)$ & $\mathcal{O}(n)$/paso & Enrutamiento global expresivo de línea base \citep{vaswani_attention_2017}. \\
Atención Sigmoide & $\sigma(QK^\top/\sqrt{d_k}+b)V$ & $\mathcal{O}(n^2d)$ & $\mathcal{O}(n)$/paso & Compuerta elemento a elemento; a menudo necesita norma de estabilización \citep{ramapuram_theory_2024}. \\
RetNet & $(QK^\top\odot D)V$ & $\mathcal{O}(n^2d)$ o por fragmentos & $\mathcal{O}(1)$/paso & Forma dual paralela/recurrente con retención de decaimiento \citep{sun_retentive_2023}. \\
Mamba & $h_t=\bar{A}_t h_{t-1}+\bar{B}_t x_t$ & $\mathcal{O}(nd)$ & $\mathcal{O}(1)$/paso & Espacio de estados selectivo con barrido consciente de hardware \citep{gu_mamba_2023}. \\
Bloque estilo EDO & $\frac{dh}{dt}=f_\theta(h,t)$ & depende del solver & depende del solver & Interpretación de profundidad continua; integración RK \citep{zhang_continuous_2021}. \\
Memoria Titans & $M_t=(1-\alpha_t)M_{t-1}+S_t$ & aprox.\ $\mathcal{O}(nd)$ & centrado en recuperación & Actualizaciones de memoria en tiempo de prueba con dinámica impulsada por sorpresa \citep{behrouz_titans_2025}. \\
\bottomrule
\end{longtable}
\normalsize

\subsection{Resumen de Optimizadores}
\small
\begin{longtable}{p{2.2cm}p{2.8cm}p{2.0cm}p{5.4cm}p{1.8cm}}
\toprule
\textbf{Optimizador} & \textbf{Familia} & \textbf{Costo de Estado} & \textbf{Idea Clave} & \textbf{Ref} \\
\midrule
\endhead
SGD+Momentum & Primer orden clásico & Bajo & Línea base acelerada por momento con estado mínimo & \citep{polyak_1964} \\
AdamW & Adaptativo primer/segundo momento & Alto & Línea base con decaimiento de peso desacoplado & \citep{loshchilov_decoupled_2017} \\
RAdam & Adaptativo con corrección de varianza & Alto & Rectifica la varianza adaptativa temprana & \citep{liu_variance_2019} \\
Adan & Momento + reducción de varianza & Alto & Actualización adaptativa estilo Nesterov & \citep{xie_adan_2022} \\
ADOPT & Variante de Adam & Alto & Actualizaciones reordenadas con mejores garantías de convergencia & \citep{taniguchi_adopt_2024} \\
AdEMAMix & Adaptativo multi-EMA & Alto & Mezcla EMAs de horizonte corto y largo & \citep{pagliardini_ademamix_2024} \\
MARS & Precondicionado con reducción de varianza & Alto & Corrección de momento recursiva & \citep{yuan_mars_2024} \\
Cautious AdamW & Momento enmascarado & Alto & Aplica actualizaciones solo en direcciones consistentes con signo & \citep{liang_cautious_2024} \\
LAMB & Momentos adaptativos por capa & Alto & Escalado de razón de confianza para entrenamiento con lotes muy grandes & \citep{you_largebatch_2020} \\
Schedule-free AdamW & Adaptativo sin programador & Alto & Elimina la dependencia de programación LR explícita & \citep{defazio_road_2024} \\
Adafactor & Adaptativo eficiente en memoria & Medio & Segundos momentos factorizados para tensores matriciales & \citep{shazeer_adafactor_2018} \\
GaLore AdamW & Proyección de gradiente de bajo rango & Medio & Optimiza en espacio de gradiente de bajo rango proyectado & \citep{zhao_galore_2024} \\
APOLLO & Proyección adaptativa de bajo rango & Medio & Escalado estructurado a partir de momentos estilo Adam proyectados & \citep{zhu2025apollosgdlikememoryadamwlevel} \\
APOLLO-Mini & Proyección adaptativa rango-1 & Bajo & Variante APOLLO escalada por tensor para ahorro extremo de memoria & \citep{zhu2025apollosgdlikememoryadamwlevel} \\
Q-APOLLO & Proyección de bajo rango cuantizada & Bajo & Estado APOLLO cuantizado para entrenamiento de memoria ultra-baja & \citep{zhu2025apollosgdlikememoryadamwlevel} \\
Anon & Similar a Adam con adaptabilidad ajustable & Alto & Adaptabilidad $\gamma \in \mathbb{R}$ ajustable continuamente con actualización de retardo incremental (IDU) para convergencia & \citep{anon2026anon} \\
Prodigy & Adaptación sin parámetros & Medio & Calibración de paso adaptativa a distancia & \citep{mishchenko_prodigy_2023} \\
Lion & Momento de signo & Bajo & Actualización de signo de momento, estado reducido & \citep{chen_symbolic_2023} \\
Sophia & Segundo orden aproximado & Medio & Precondicionamiento diagonal de Hessiana con recorte & \citep{liu_sophia_2023} \\
Shampoo & Precondicionador matricial & Alto & Estadísticas de segundo orden estructuradas Kronecker & \citep{gupta_shampoo_2018} \\
SOAP & Shampoo + base Adam & Alto & Seguimiento similar a Adam en base propia del precondicionador & \citep{vyas_soap_2024} \\
Muon & Basado en ortogonalidad & Medio & Actualizaciones matriciales ortogonalizadas & \citep{shen_convergence_2025} \\
Turbo-Muon & Ortogonalización acelerada & Medio & Aceleración Newton-Schulz precondicionada & \citep{boissin_turbo-muon_2025} \\
\bottomrule
\end{longtable}
\normalsize


\section{Discusión}
\label{sec:discussion}

Las decisiones de diseño en Transformer Encoder Frankenstein reflejan varias tensiones de ingeniería e investigación en las herramientas modernas de aprendizaje profundo.

\subsection{Compromisos de Diseño Basado en Esquemas}

El enfoque basado en esquemas proporciona beneficios significativos de reproducibilidad al imponer contratos explícitos y fallar rápidamente ante configuraciones inválidas. Sin embargo, este enfoque también introduce rigidez: agregar nuevas arquitecturas u optimizadores requiere extensiones del esquema en lugar de argumentos de línea de comandos flexibles. El sistema de hiperparámetros prefijados permite un control detallado pero aumenta la complejidad de configuración para usuarios acostumbrados a interfaces más simples.

La decisión de imponer \texttt{additionalProperties: false} en todos los niveles del esquema elimina la absorción silenciosa de parámetros que ha afectado a sistemas de configuración anteriores, pero esta rigurosidad requiere un mantenimiento cuidadoso del esquema al extender las capacidades del sistema. Cada nuevo mecanismo de atención o variante de optimizador debe integrarse adecuadamente en el marco de validación, incluyendo definiciones de campos del esquema con tipos y restricciones apropiados, mapeo de hiperparámetros prefijados para grupos específicos de optimizadores, valores predeterminados alineados con las mejores prácticas de investigación y cadenas de documentación para la representación en la interfaz web.

\subsection{Cobertura Arquitectónica y Vacíos}

Las diecisiete arquitecturas de mezclador implementadas abarcan las principales direcciones de investigación en modelado de secuencias, pero ciertos vacíos permanecen. El sistema carece de arquitecturas híbridas recientes como Griffin \citep{soares_griffin_2024} y Jamba \citep{gu_jamba_2024}, que combinan compuertas con modelos de espacio de estados. El enrutamiento MoE (Mezcla de Expertos) está implementado para capas FFN pero no para el cómputo de atención, donde trabajos recientes han mostrado beneficios \citep{lewis_moe_2021}.

La cobertura de atención dispersa es completa, pero la implementación de métodos sin entrenamiento (FASA, SpargeAttn) genera errores en tiempo de ejecución durante el entrenamiento, reflejando restricciones arquitectónicas: estos métodos requieren checkpoints preentrenados de modelos de atención completa o procedimientos específicos de ajuste fino que actualmente no están automatizados.

La cobertura de mecanismos de compuerta es sólida en todas las categorías principales (GLA, DeltaNet, Gated DeltaNet, Gated DeltaNet-2, HGRN2, FoX, Gated Softmax).

\subsection{Fragmentación del Panorama de Optimizadores}

El soporte para veintidós familias de optimizadores en seis categorías algorítmicas demuestra exhaustividad pero también resalta el estado fragmentado de la investigación en optimización. Los usuarios enfrentan una complejidad significativa de decisión al elegir entre métodos de reducción de varianza (Adan, MARS), variantes eficientes en memoria (GaLore, Adafactor, APOLLO) y enfoques conscientes de curvatura (Sophia, Shampoo). El sistema de hiperparámetros prefijados, aunque poderoso, requiere comprender qué parámetros son relevantes para cada clase de optimizador.

La calidad de implementación varía entre optimizadores: los métodos clásicos (AdamW, SGD con momento) están altamente optimizados en PyTorch, mientras que los métodos más nuevos (Muon, Turbo-Muon, SOAP) pueden requerir implementaciones personalizadas que afectan la estabilidad numérica y las características de rendimiento.

\subsection{Consideraciones de Despliegue y Producción}

El pipeline de cuantización demuestra preocupaciones prácticas de despliegue pero realiza compromisos de ingeniería específicos. El empaquetado ternario de pesos reduce el almacenamiento a aproximadamente $1.58$ bits por parámetro, pero esta compresión agresiva puede degradar el rendimiento, especialmente para modelos más pequeños donde el error de cuantización es más significativo. La implementación actual aplica cuantización uniforme entre todos los tipos de parámetros.

Los flujos de trabajo SBERT proporcionan utilidad práctica para tareas de similitud semántica y recuperación, pero la implementación asume estrategias de pooling estándar (token CLS, pooling promedio). Avances recientes como los embeddings Matryoshka \citep{muennighoff_matryoshka_2021} y los refinamientos de aprendizaje contrastivo \citep{chen_supervised_2020} aún no están incorporados.

\subsection{Desafíos de Integración y Extensibilidad}

La estructura actual del código base, aunque funcional, presenta desafíos de mantenimiento a medida que se expanden las familias de arquitecturas y optimizadores. El patrón de despachador para la selección de mezcladores y el enrutamiento de optimizadores maneja la extensibilidad pero corre el riesgo de convertirse en un ``cajón de sastre'' de lógica condicional. Las versiones futuras se beneficiarían de arquitecturas basadas en plugins donde nuevos mezcladores y optimizadores puedan registrarse de forma declarativa en lugar de modificar la lógica central de despacho.

La interfaz de configuración web proporciona mejoras significativas de usabilidad pero introduce complejidad de despliegue: ejecutar Streamlit junto con trabajos de entrenamiento requiere recursos adicionales y consideraciones de infraestructura que pueden no ser apropiadas para todos los entornos, particularmente clústeres HPC sin acceso web.


\section{Conclusión}
\label{sec:conclusion}

Transformer Encoder Frankenstein presenta una plataforma de experimentación unificada, impulsada por configuración, que aborda desafíos críticos en la investigación moderna de aprendizaje profundo: fragmentación arquitectónica entre atención densa, modelos recurrentes, patrones dispersos y mecanismos de compuerta; complejidad del panorama de optimizadores que abarca líneas base clásicas, métodos de reducción de varianza, variantes eficientes en memoria, enfoques sin programación, algoritmos conscientes de curvatura y métodos orientados a geometría; y flujos de trabajo de despliegue de extremo a extremo que abarcan cuantización y aplicaciones de embeddings de oraciones.

Las contribuciones principales del sistema son:
\begin{enumerate}[leftmargin=*]
    \item \textbf{Diseño Basado en Esquemas}: Un contrato de configuración estricto basado en YAML con validación y enrutamiento de hiperparámetros prefijados que permite experimentos reproducibles en diecisiete arquitecturas de mezclador y veintitrés familias de optimizadores.
    \item \textbf{Soporte Arquitectónico Integral}: Implementación que abarca las principales categorías de investigación incluyendo líneas base densas (atención estándar, sigmoide), alternativas recurrentes (RetNet, Mamba, estilo EDO, Titans), atención dispersa (Sparse Transformer, Longformer, BigBird, SparseK, NSA, SpargeAttn, FASA) y mecanismos de compuerta (GLA, DeltaNet, Gated DeltaNet, Gated DeltaNet-2, HGRN2, FoX, Gated Softmax).
    \item \textbf{Marco Unificado de Optimizadores}: Grupos de hiperparámetros prefijados que permiten control detallado sobre embeddings, capas de normalización, bloques recurrentes, pesos de atención y parámetros FFN en líneas base clásicas (SGD+Momentum, AdamW), reducción de varianza (Adan, ADOPT, AdEMAMix, MARS, Cautious), eficientes en memoria (Adafactor, GaLore, Lion, APOLLO, APOLLO-Mini, Q-APOLLO), lotes grandes y simplificación de programación (LAMB, Schedule-Free AdamW, Prodigy), conscientes de curvatura (Sophia), segundo orden (Shampoo, SOAP) y orientados a geometría (Muon, Turbo-Muon).
    \item \textbf{Flujos de Trabajo de Extremo a Extremo}: Pipeline de despliegue integrado que soporta empaquetado ternario de pesos y cuantización de activaciones INT8; entrenamiento e inferencia inspirados en SBERT para similitud semántica, recuperación y tareas de clustering.
    \item \textbf{Configuración Interactiva}: Interfaz web basada en Streamlit que proporciona generación de formularios impulsada por esquemas, validación en tiempo real, documentación en línea y síntesis de comandos CLI.
\end{enumerate}

Este sistema permite iteración experimental rápida mientras mantiene reproducibilidad mediante contratos de configuración estrictos. Al consolidar diversas contribuciones de investigación en un conjunto de herramientas unificado, reduce las barreras para explorar arquitecturas y estrategias de optimización novedosas, particularmente para investigadores que pueden carecer de recursos para implementar y validar cada variante de forma independiente.

\subsection{Limitaciones y Direcciones Futuras}

Varias limitaciones y direcciones prometedoras para trabajo futuro surgen del diseño e implementación de este sistema:

\begin{enumerate}[leftmargin=*]
    \item \textbf{Integración Arquitectónica}: Arquitecturas híbridas recientes (Griffin, Jamba, Mamba-X) demuestran beneficios de combinar múltiples mecanismos en bloques unificados. Versiones futuras deberían integrar estas arquitecturas y explorar patrones de composición sistemática.
    
    \item \textbf{Cuantización Avanzada}: La implementación actual utiliza empaquetado ternario uniforme en todos los parámetros. La investigación sobre cuantización por capas, por canales y consciente de importancia sugiere que estrategias más sofisticadas podrían mejorar los compromisos calidad-eficiencia.
    
    \item \textbf{Extensibilidad Basada en Plugins}: El patrón de despacho actual se vuelve cada vez más complejo con cada nueva adición. Una arquitectura de plugins que permita el registro declarativo de nuevos mezcladores, optimizadores y métodos de normalización mejoraría el mantenimiento y reduciría el riesgo de errores en la lógica central de despacho.
    
    \item \textbf{Optimización Automatizada de Hiperparámetros}: El esquema soporta espacios extensos de hiperparámetros, pero los usuarios deben explorar estos espacios manualmente. La integración con optimización bayesiana, estrategias de bandidos multi-brazo o ajuste de hiperparámetros basado en gradientes podría automatizar el descubrimiento de configuraciones efectivas.
    
    \item \textbf{Despliegue en Producción}: La interfaz web mejora la usabilidad pero puede no ser apropiada para todos los entornos de despliegue. Modos de configuración sin interfaz gráfica, gestión de configuración basada en API o ergonomía CLI mejorada podrían servir a flujos de trabajo HPC y de producción.
    
    \item \textbf{Evaluación Comparativa}: Si bien el sistema permite entrenamiento con diversas arquitecturas, una evaluación comparativa exhaustiva que compare el rendimiento entre mezcladores y optimizadores en tareas estandarizadas proporcionaría orientación valiosa para la selección de configuraciones.
    
    \item \textbf{Garantías de Estabilidad de Entrenamiento}: La implementación actual incluye guardas contra NaN/Inf y recorte de gradiente, pero el análisis formal de condiciones de estabilidad para diferentes combinaciones mezclador-optimizador, particularmente con bloques en bucle y cuantización agresiva, sigue abierto.
    
    \item \textbf{Extensiones Multimodales y Específicas de Tarea}: El diseño actual se centra en modelado de secuencias. Extensiones para modelos de visión-lenguaje, arquitecturas multimodales y flujos de trabajo de ajuste fino específicos de tarea (ej., ajuste por instrucciones, RLHF) ampliarían la aplicabilidad.
\end{enumerate}

La trayectoria de investigación del modelado de secuencias continúa hacia enfoques híbridos que combinan fortalezas de múltiples paradigmas---compresión de la recurrencia, selectividad de la atención, compuertas para gestión de memoria y dispersión para eficiencia. Una plataforma de experimentación unificada como Transformer Encoder Frankenstein es cada vez más valiosa a medida que esta convergencia se acelera, permitiendo a los investigadores explorar sistemáticamente este creciente espacio de diseño con infraestructura reproducible y bien diseñada.



\bibliographystyle{plainnat}
\bibliography{../bibliography/other,../bibliography/optimizers,../bibliography/attention_types}

\appendix

\input{appendices/annex-a-optimizadores}
\section{Anexo B: Transformadores Densos, Recurrentes y Aumentados con Memoria}
\label{sec:annex-transformer}

Este anexo exhaustivo sintetiza las arquitecturas modernas de transformers más allá de la línea base estándar de softmax. El campo ha evolucionado hacia seis paradigmas principales: atención densa global con variantes, compresión de cabezas clave-valor mediante atención por consultas agrupadas, compresión de estado recurrente con decaimiento, modelos de espacio de estado selectivos, integración numérica de profundidad continua, arquitecturas aumentadas con memoria y enfoques híbridos. Comprender esta taxonomía ilumina las compensaciones fundamentales entre expresividad, costo computacional, huella de memoria y simplicidad de despliegue.

\subsection{Atención Densa de Referencia: Estándar y Sigmoide}

\subsubsection{Atención Softmax Estándar}
La atención multi-cabeza estándar \citep{vaswani_attention_2017} calcula la similitud de producto punto escalado entre embeddings de tokens:

\begin{algorithmic}[1]
\State \textbf{Entrada:} Matriz de consulta $\mathbf{Q} \in \mathbb{R}^{n \times d}$, Matriz de clave $\mathbf{K} \in \mathbb{R}^{n \times d}$, Matriz de valor $\mathbf{V} \in \mathbb{R}^{n \times d}$
\State $\text{puntajes} \gets \frac{\mathbf{Q}\mathbf{K}^\top}{\sqrt{d}} \in \mathbb{R}^{n \times n}$ \Comment{calcular similitudes por pares}
\State $\text{pesos\_atencion} \gets \texttt{softmax}(\text{puntajes}, \text{eje}=1) \in \mathbb{R}^{n \times n}$ \Comment{normalizar sobre las claves}
\State \textbf{Salida:} $\mathbf{Y} = \text{pesos\_atencion} \cdot \mathbf{V} \in \mathbb{R}^{n \times d}$ \Comment{combinación ponderada de valores}
\end{algorithmic}

Para generación autorregresiva, el almacenamiento en caché KV guarda claves y valores pasados para evitar el recálculo $\mathcal{O}(n^2)$. Sin embargo, este caché de crecimiento lineal ocupa $\sim n \cdot d_{\text{oculto}}$ bytes, lo que puede consumir cientos de gigabytes para modelos de miles de millones de parámetros. La atención estándar logra una \textbf{expresividad perfecta} dentro de la ventana de contexto---cualquier token puede atender a cualquier otro token con pesos aprendidos---pero paga el precio del cómputo denso.

\begin{itemize}[leftmargin=*]
    \item \textbf{Complejidad de entrenamiento}: $\mathcal{O}(n^2 \cdot d)$ tiempo, $\mathcal{O}(n^2)$ espacio (materialización de la matriz de atención).
    \item \textbf{Complejidad de inferencia}: $\mathcal{O}(n)$ tiempo por token, $\mathcal{O}(n)$ espacio (caché KV).
    \item \textbf{Fortalezas}: Expresividad incomparable; recuperación perfecta del historial; altamente paralelizable.
    \item \textbf{Debilidades}: El cuello de botella cuadrático prohíbe contextos muy largos; el caché KV domina la memoria durante la generación.
\end{itemize}

\subsubsection{Atención Sigmoide}
La atención sigmoide reemplaza el softmax por filas con una activación sigmoide elemento a elemento \citep{ramapuram_theory_2024}:

\begin{algorithmic}[1]
\State \textbf{Entrada:} $\mathbf{Q}, \mathbf{K}, \mathbf{V}$ como arriba, sesgo aprendible $\mathbf{b} \in \mathbb{R}^{n \times n}$
\State $\text{logits} \gets \frac{\mathbf{Q}\mathbf{K}^\top}{\sqrt{d}} + \mathbf{b}$
\State $\text{pesos\_atencion} \gets \sigma(\text{logits})$ \Comment{sigmoide elemento a elemento, no softmax}
\State \textbf{Salida:} $\mathbf{Y} = \text{pesos\_atencion} \odot \mathbf{V}$
\end{algorithmic}

A diferencia de softmax, la sigmoide no impone una distribución de probabilidad (los pesos no necesitan sumar 1), lo que permite una independencia de tokens más fuerte. El análisis teórico mediante mezcla de expertos muestra que la sigmoide alcanza una complejidad muestral superior: convergencia $\mathcal{O}(n^{-0.51})$ para expertos ReLU frente a $\mathcal{O}(n^{-0.24})$ de softmax. Sin embargo, el entrenamiento empírico reveló inestabilidades de gradiente a escala. El remedio es la \textbf{norma híbrida}---agregar normalización después de la salida de atención---que estabiliza los gradientes sin sacrificar los beneficios teóricos del enrutamiento elemento a elemento.

\begin{itemize}[leftmargin=*]
    \item \textbf{Complejidad de entrenamiento}: $\mathcal{O}(n^2 \cdot d)$ (idéntico al estándar), pero las operaciones elemento a elemento permiten una aceleración de inferencia del 17\% mediante FlashSigmoid.
    \item \textbf{Complejidad de inferencia}: $\mathcal{O}(n)$ por token con caché KV (asintóticamente igual, pero con factores constantes más bajos).
    \item \textbf{Fortalezas}: Supera la competencia de suma cero; evita la sincronización por filas; implementación eficiente en hardware.
    \item \textbf{Debilidades}: Requiere estabilización cuidadosa (norma híbrida); inestabilidad de entrenamiento a grandes escalas sin pérdida auxiliar.
\end{itemize}

\subsubsection{Atención por Consultas Agrupadas (GQA)}
La atención por consultas agrupadas (GQA) \citep{ainslie_gqa_2023} tiende un puente entre la atención multi-cabeza (MHA) y la atención multi-consulta (MQA) al particionar las cabezas de consulta en grupos, cada uno compartiendo una sola cabeza clave-valor. Si $h$ es el número total de cabezas de atención, $h_k$ es el número de cabezas clave-valor, y $G = h / h_k$ es el tamaño del grupo, entonces cada grupo de $G$ cabezas de consulta comparte una cabeza KV:

\begin{algorithmic}[1]
\State \textbf{Entrada:} $\mathbf{Q} \in \mathbb{R}^{n \times h \cdot d_h}$, $\mathbf{K} \in \mathbb{R}^{n \times h_k \cdot d_h}$, $\mathbf{V} \in \mathbb{R}^{n \times h_k \cdot d_h}$
\State $\mathbf{K}_{\text{completo}} \gets \texttt{repeat\_interleave}(\mathbf{K}, G)$ \Comment{expandir de $h_k$ a $h$ cabezas}
\State $\mathbf{V}_{\text{completo}} \gets \texttt{repeat\_interleave}(\mathbf{V}, G)$
\State $\text{puntajes} \gets \frac{\mathbf{Q}\mathbf{K}_{\text{completo}}^\top}{\sqrt{d_h}} \in \mathbb{R}^{n \times n}$ \Comment{producto punto estándar tras expansión}
\State $\text{pesos\_atencion} \gets \texttt{softmax}(\text{puntajes}, \text{eje}=1)$
\State \textbf{Salida:} $\mathbf{Y} = \text{pesos\_atencion} \cdot \mathbf{V}_{\text{completo}} \in \mathbb{R}^{n \times h \cdot d_h}$
\end{algorithmic}

GQA interpola suavemente entre MHA ($h_k = h$, $G = 1$) y MQA ($h_k = 1$, $G = h$). La ventaja clave es un balance ajustable entre la huella del caché KV y la capacidad representacional. Cada cabeza KV adicional consume $2 \cdot d_h \cdot n$ bytes extra en el caché, por lo que reducir $h_k$ produce ahorros proporcionales de memoria.

\begin{itemize}[leftmargin=*]
    \item \textbf{Complejidad de entrenamiento}: $\mathcal{O}(n^2 \cdot h \cdot d_h)$ (idéntico a MHA tras la expansión de cabezas).
    \item \textbf{Complejidad de inferencia}: $\mathcal{O}(n)$ por token, espacio de caché KV $\mathcal{O}(h_k \cdot n \cdot d_h)$.
    \item \textbf{Fortalezas}: Balance flexible de cabezas; probado a escala (Llama 2 y 3 usan $h_k = 8$ cabezas KV con $h = 32$ cabezas de consulta); reduce el caché KV en factor $G \times$; soporta inferencia rápida con FlashAttention.
    \item \textbf{Debilidades}: Menos cabezas KV pueden reducir la capacidad representacional en modelos pequeños; requerir que $h_k$ divida a $h$ restringe las opciones arquitectónicas.
\end{itemize}

\subsection{Arquitecturas Recurrentes y Retentivas}

\subsubsection{Redes Retentivas (RetNet)}
RetNet \citep{sun_retentive_2023} unifica tres paradigmas de cómputo: entrenamiento paralelo, inferencia recurrente y despliegue por fragmentos. Su innovación central es el \textbf{mecanismo de retención}, que utiliza una matriz de decaimiento exponencial fija para modelar la importancia temporal:

\begin{algorithmic}[1]
\State \textbf{Forma paralela (entrenamiento):}
\State $\text{matriz\_decaimiento}[i,j] \gets \gamma^{i-j}$ para $i \geq j$, si no $0$ \Comment{decaimiento exponencial causal}
\State $\text{matriz\_decaimiento}[i,j] \gets 0$ para $i < j$ \Comment{enmascaramiento causal}
\State $\mathbf{Y}_{\text{paralelo}} \gets (\mathbf{Q}\mathbf{K}^\top \odot \text{matriz\_decaimiento}) \mathbf{V}$ \Comment{multiplicación elemento a elemento con decaimiento}
\end{algorithmic}

\begin{algorithmic}[1]
\State \textbf{Forma recurrente (inferencia):}
\For{$t = 1, \ldots, n$}
    \State $\mathbf{s}_t \gets \gamma \mathbf{s}_{t-1} + \mathbf{k}_t \mathbf{v}_t^\top$ \Comment{actualización de estado con decaimiento}
    \State $\mathbf{y}_t \gets \mathbf{q}_t \mathbf{s}_t$ \Comment{salida mediante interacción consulta-estado}
\EndFor
\end{algorithmic}

El escalar de decaimiento $\gamma \in (0,1)$ controla la ventana temporal. RetNet utiliza \textbf{retención multi-escala} con diferentes valores de $\gamma$ por cabeza (ej., $\gamma = 1 - 2^{-5}, 1 - 2^{-6}, \ldots$), permitiendo dependencias a corto y largo plazo simultáneamente. El modo recurrente por fragmentos divide las secuencias en fragmentos, procesa cada fragmento en paralelo y enhebra un estado recurrente entre fragmentos.

\begin{itemize}[leftmargin=*]
    \item \textbf{Complejidad de entrenamiento}: $\mathcal{O}(n^2 \cdot d)$ (paralela), o $\mathcal{O}(n \cdot c \cdot d)$ (recurrente por fragmentos con tamaño de fragmento $c$).
    \item \textbf{Complejidad de inferencia}: $\mathcal{O}(1)$ por token, $\mathcal{O}(d^2)$ espacio de estado (matriz fija).
    \item \textbf{Fortalezas}: Inferencia en tiempo constante; triple paradigma de cómputo; consolidación multi-escala.
    \item \textbf{Debilidades}: El decaimiento fijo impone un sesgo inductivo rígido; puede truncar patrones de largo alcance aprendidos.
\end{itemize}

\subsubsection{Mamba: Modelos de Espacio de Estado Selectivos}
Mamba \citep{gu_mamba_2023} enmarca la recurrencia como un sistema dinámico de tiempo continuo con parámetros dependientes de la entrada, logrando tanto entrenamiento lineal como inferencia en tiempo constante:

\begin{algorithmic}[1]
\State \textbf{Dinámica continua:} $h'(t) = \mathbf{A} h(t) + \mathbf{B} x(t)$, \quad $y(t) = \mathbf{C} h(t)$
\State \textbf{Discretización con tamaño de paso $\Delta_t$:}
\State $\bar{\mathbf{A}}_t \gets \exp(\Delta_t \mathbf{A})$
\State $\bar{\mathbf{B}}_t \gets (\Delta_t \mathbf{A})^{-1} (\exp(\Delta_t \mathbf{A}) - \mathbf{I}) \Delta_t \mathbf{B}_t$
\State \textbf{Actualización recurrente:}
\For{$t = 1, \ldots, n$}
    \State $\Delta_t \gets \text{softplus}(\text{Lineal}(x_t))$ \Comment{tamaño de paso dependiente de la entrada}
    \State $\mathbf{B}_t \gets \text{Lineal}(x_t)$, \quad $\mathbf{C}_t \gets \text{Lineal}(x_t)$ \Comment{matrices de proyección dependientes de la entrada}
    \State $h_t \gets \bar{\mathbf{A}}_t h_{t-1} + \bar{\mathbf{B}}_t x_t$ \Comment{transición de estado}
    \State $y_t \gets \mathbf{C}_t h_t$ \Comment{proyección de salida}
\EndFor
\end{algorithmic}

La innovación crítica es que $\mathbf{A}$ (la matriz del sistema) \textit{no} depende de la entrada, pero $\Delta_t$, $\mathbf{B}_t$ y $\mathbf{C}_t$ sí, haciendo que el sistema sea \textbf{variante en el tiempo}. Esta selectividad permite al modelo \textit{ignorar} información irrelevante estableciendo $\Delta_t$ cerca de cero, creando efectivamente una compuerta. Un algoritmo de escaneo paralelo consciente del hardware implementa la recurrencia eficientemente en GPUs fusionando el cómputo dentro de SRAM, evitando el costoso ancho de banda de HBM.

\begin{itemize}[leftmargin=*]
    \item \textbf{Complejidad de entrenamiento}: $\mathcal{O}(n \cdot d)$ mediante escaneo consciente del hardware (lineal en la longitud de la secuencia).
    \item \textbf{Complejidad de inferencia}: $\mathcal{O}(1)$ por token, $\mathcal{O}(d)$ estado (vector oculto).
    \item \textbf{Fortalezas}: Logra entrenamiento lineal e inferencia constante simultáneamente; práctico en secuencias largas (millones de tokens).
    \item \textbf{Debilidades}: La compresión del vector de estado puede debilitar la copia exacta y la recuperación asociativa densa en comparación con la atención completa.
\end{itemize}

\subsection{Transformers de Profundidad Continua: Integración EDO}

\subsubsection{Transformer EDO}
El Transformer EDO \citep{zhang_continuous_2021} interpreta la profundidad de la red como integración numérica de un sistema dinámico continuo, utilizando solucionadores Runge-Kutta de orden superior para reducir el error de truncamiento:

\begin{algorithmic}[1]
\State \textbf{Formulación continua:} $\frac{dh(t)}{dt} = f_\theta(h(t), t)$ donde $f_\theta$ es la subred del transformer
\State \textbf{Aproximación discreta Runge-Kutta-4:}
\State $k_1 \gets f_\theta(h_t, t)$
\State $k_2 \gets f_\theta(h_t + \tfrac{1}{2}k_1, t + \tfrac{1}{2}\Delta t)$
\State $k_3 \gets f_\theta(h_t + \tfrac{1}{2}k_2, t + \tfrac{1}{2}\Delta t)$
\State $k_4 \gets f_\theta(h_t + k_3, t + \Delta t)$
\State $h_{t+1} \gets h_t + \frac{\Delta t}{6}(k_1 + 2k_2 + 2k_3 + k_4)$ \Comment{combinación ponderada}
\end{algorithmic}

En lugar de las conexiones residuales simples de Euler $h_{t+1} = h_t + f_\theta(h_t)$, el bloque RK4 calcula cuatro evaluaciones intermedias y las combina con los pesos clásicos de RK4. Para evitar gradientes que se desvanecen, la arquitectura introduce \textbf{compuertas aprendidas} que interpolar entre aproximaciones intermedias:

\begin{algorithmic}[1]
\State $g \gets \sigma(\text{Lineal}([k_1, k_2, k_3, k_4]))$ \Comment{compuerta aprendible}
\State $h_{t+1} \gets h_t + g \cdot k_1 + (1-g) \cdot k_2$ \Comment{refinamiento interpolado}
\end{algorithmic}

Esta formulación reduce el número efectivo de parámetros mediante el uso compartido de pesos---la misma $f_\theta$ se evalúa múltiples veces---proporcionando un refinamiento de trayectoria más rico.

\begin{itemize}[leftmargin=*]
    \item \textbf{Complejidad de entrenamiento}: $\mathcal{O}(k \cdot n^2 \cdot d)$ donde $k$ es el orden RK (4 para RK4).
    \item \textbf{Complejidad de inferencia}: $\mathcal{O}(k \cdot n)$ por token (mayor sobrecarga constante).
    \item \textbf{Fortalezas}: Precisión significativamente mayor en tareas de generación (BLEU de última generación); eficiente en parámetros mediante uso compartido de pesos.
    \item \textbf{Debilidades}: Alto costo de cómputo por paso; la latencia de inferencia aumenta por un factor constante $k$; requiere compuertas complejas para la estabilidad.
\end{itemize}

\subsection{Memoria en Tiempo de Prueba: Titans}

La arquitectura Titans \citep{behrouz_titans_2025} introduce una dimensión ortogonal: en lugar de pesos estáticos, el modelo mantiene una memoria aprendible que se \textit{actualiza durante la inferencia} basada en una señal impulsada por sorpresa:

\begin{algorithmic}[1]
\State \textbf{Atención local a corto plazo:} Aplicar atención estándar o dispersa dentro de una ventana de contexto fija $c$
\State $y_t^{\text{local}} \gets \text{Atencion}(q_t, k_{[t-c:t]}, v_{[t-c:t]})$
\State \textbf{Señal de actualización de memoria (sorpresa):}
\State $S_t \gets \eta_t S_{t-1} - \theta_t \nabla_\ell(\mathcal{M}_{t-1}; x_t)$ \Comment{sorpresa como gradiente de la pérdida respecto a la memoria}
\State $\mathcal{M}_t \gets (1 - \alpha_t) \mathcal{M}_{t-1} + S_t$ \Comment{memoria actualizada mediante EMA}
\State \textbf{Recuperación de memoria a largo plazo:}
\State $y_t^{\text{memoria}} \gets \mathcal{M}_t^* (q_t)$ \Comment{consultar módulo de memoria para información recuperada}
\State \textbf{Combinación de salida:}
\State $y_t \gets y_t^{\text{local}} + \text{compuerta}(y_t^{\text{memoria}})$ \Comment{combinar memoria local y recuperada}
\end{algorithmic}

El módulo de memoria $\mathcal{M}$ se actualiza literalmente durante el paso hacia adelante calculando gradientes de una pérdida asociativa y aplicando pasos de SGD con momento. Las tasas de decaimiento $\eta_t, \theta_t, \alpha_t$ son en sí mismas dependientes de la entrada, permitiendo al modelo cambiar de paradigma de memoria cuando el contexto cambia.

\begin{itemize}[leftmargin=*]
    \item \textbf{Complejidad de entrenamiento}: Aproximadamente $\mathcal{O}(c^2 + n \cdot f)$ donde $c$ es la ventana local y $f$ es la sobrecarga de actualización de memoria.
    \item \textbf{Complejidad de inferencia}: $\mathcal{O}(c^2)$ atención local más $\mathcal{O}(1)$ recuperación de memoria por token.
    \item \textbf{Fortalezas}: Maneja longitudes de contexto extremas; permite recuperación asociativa real; la memoria se adapta a la entrada.
    \item \textbf{Debilidades}: Significativamente más complejo; la inferencia incluye cálculo de gradientes; mayor sobrecarga de coordinación.
\end{itemize}

\subsection{Atención Latente Multi-Cabeza (MLA)}
La Atención Latente Multi-Cabeza (MLA, por sus siglas en inglés) \citep{mehta_mla_2025} comprime el estado clave--valor por token en un único vector latente de rango bajo $c_{KV}$ de rango $r_{kv}$, el cual es la única cantidad materializada en el caché KV durante la decodificación. Un par de matrices de up-proyección reconstruye las claves y valores completos de todas las cabezas al vuelo, mientras que RoPE se aplica a los tensores de consulta y clave descomprimidos para preservar la estructura posicional. Los autores muestran que el ancho latente óptimo de Pareto se sitúa cerca de $r_{kv} = d_{\text{oculto}} / 2$, punto en el cual MLA reduce a la mitad el caché KV pero iguala la calidad de la atención multi-cabeza (MHA) en modelos pequeños.

\subsubsection{Formulación Matemática}
\[
c_{KV} = W_{DKV}\, x \in \mathbb{R}^{r_{kv}}, \qquad
k = W_{UK}\, c_{KV}, \qquad
v = W_{UV}\, c_{KV}, \qquad
q = W_{Q}\, x.
\]
\[
\tilde q = \text{RoPE}(q), \quad \tilde k = \text{RoPE}(k), \qquad
\text{atencion} = \texttt{softmax}\!\left(\frac{\tilde q\, \tilde k^\top}{\sqrt{d_h}}\right) v.
\]

\subsubsection{Pseudocódigo Algorítmico}
\begin{algorithmic}[1]
\State \textbf{Entrada:} estado oculto $x \in \mathbb{R}^{n \times d}$, rango $r_{kv}$, proyecciones $W_{DKV}, W_{UK}, W_{UV}, W_{Q}$
\State $c_{KV} \gets W_{DKV}\, x$ \Comment{down-proyección a latente de rango $r_{kv}$ (en caché en inferencia)}
\State $k \gets W_{UK}\, c_{KV}$, \quad $v \gets W_{UV}\, c_{KV}$ \Comment{up-proyección a cabezas completas}
\State $q \gets W_{Q}\, x$
\State $\tilde q \gets \text{RoPE}(q)$, \quad $\tilde k \gets \text{RoPE}(k)$ \Comment{aplicar embeddings rotatorios}
\State $\text{pesos} \gets \texttt{softmax}(\tilde q\, \tilde k^\top / \sqrt{d_h})$
\State \textbf{Salida:} $\mathbf{Y} = \text{pesos} \cdot v$
\end{algorithmic}

\subsubsection{Características Clave}
\begin{itemize}[leftmargin=*]
    \item \textbf{Complejidad de entrenamiento}: $\mathcal{O}(n^2 \cdot d)$ (atención completa sobre cabezas descomprimidas).
    \item \textbf{Complejidad de inferencia}: $\mathcal{O}(n)$ por token, espacio de caché KV $\mathcal{O}(r_{kv} \cdot n)$ (sólo latente).
    \item \textbf{Entrenable}: totalmente entrenable; $r_{kv}$ es un hiperparámetro (óptimo de Pareto $\approx d/2$).
    \item \textbf{Fortalezas}: reduce a la mitad el caché KV frente a MHA con calidad equivalente; compatible con RoPE; cuello de botella latente limpio.
    \item \textbf{Limitaciones}: la up-proyección añade FLOPs de decodificación; el rango latente es un presupuesto global fijo compartido entre cabezas.
\end{itemize}

\subsection{Atención Latente por Consultas Agrupadas (GQLA)}
La Atención Latente por Consultas Agrupadas (GQLA, por sus siglas en inglés) \citep{meng_gqla_2026} es una modificación mínima de MLA que expone dos rutas de decodificación algebraicamente equivalentes sobre los mismos pesos entrenados: una ruta MQA-absorb (que fusiona la up-proyección en la consulta, óptima en hardware tipo H100) y una ruta GQA-expanded (que materializa claves/valores completos, óptima en hardware H20 / predicción multi-token). Las dos rutas son matemáticamente idénticas porque la up-proyección es lineal, de modo que $\texttt{softmax}(q_{\text{absorbido}} \cdot c_{KV}^\top) = \texttt{softmax}((q \cdot W_{UK}^\top) \cdot c_{KV}^\top)$, permitiendo un despacho adaptativo al hardware sin reentrenamiento.

\subsubsection{Formulación Matemática}
Mismo latente que MLA:
\[
c_{KV} = W_{DKV}\, x, \qquad k = W_{UK}\, c_{KV}, \qquad v = W_{UV}\, c_{KV}, \qquad q = W_{Q}\, x.
\]
\[
\text{Ruta A (MQA-absorb): } \text{atencion} = \texttt{softmax}\!\left(\frac{(q\, W_{UK}^\top)\, c_{KV}^\top}{\sqrt{d_h}}\right) (W_{UV}\, c_{KV}).
\]
\[
\text{Ruta B (GQA-expanded): } \text{atencion} = \texttt{softmax}\!\left(\frac{q\, k^\top}{\sqrt{d_h}}\right) v, \quad k = W_{UK} c_{KV},\ v = W_{UV} c_{KV}.
\]

\subsubsection{Pseudocódigo Algorítmico}
\begin{algorithmic}[1]
\State \textbf{Entrada:} $x$, pesos $W_{DKV}, W_{UK}, W_{UV}, W_{Q}$, bandera de hardware $\text{hw} \in \{\text{H100}, \text{H20}\}$
\State $c_{KV} \gets W_{DKV}\, x$ \Comment{latente compartido, en caché}
\State $q \gets W_{Q}\, x$
\If{$\text{hw} == \text{H100}$} \Comment{ruta MQA-absorb}
    \State $q_{\text{abs}} \gets q\, W_{UK}^\top$ \Comment{absorber up-proyección en la consulta}
    \State $\text{pesos} \gets \texttt{softmax}(q_{\text{abs}}\, c_{KV}^\top / \sqrt{d_h})$
    \State $\mathbf{Y} \gets \text{pesos} \cdot (W_{UV}\, c_{KV})$
\Else \Comment{ruta GQA-expanded}
    \State $k \gets W_{UK}\, c_{KV}$, \quad $v \gets W_{UV}\, c_{KV}$
    \State $\text{pesos} \gets \texttt{softmax}(q\, k^\top / \sqrt{d_h})$
    \State $\mathbf{Y} \gets \text{pesos} \cdot v$
\EndIf
\State \textbf{Salida:} $\mathbf{Y}$
\end{algorithmic}

\subsubsection{Características Clave}
\begin{itemize}[leftmargin=*]
    \item \textbf{Complejidad de entrenamiento}: $\mathcal{O}(n^2 \cdot d)$, idéntica a MLA.
    \item \textbf{Complejidad de inferencia}: $\mathcal{O}(n)$ por token, estado latente $\mathcal{O}(r_{kv} \cdot n)$; ruta elegida por hardware.
    \item \textbf{Entrenable}: se entrena una sola vez; las dos rutas de decodificación reutilizan los mismos pesos.
    \item \textbf{Fortalezas}: decodificación adaptativa al hardware sin pérdida de calidad; la equivalencia algebraica garantiza consistencia.
    \item \textbf{Limitaciones}: mismas compensaciones de rango latente que MLA; la selección de ruta es una heurística de despliegue.
\end{itemize}

\subsection{Atención Multi-Cabeza de Rango Bajo (MLRA)}
La Atención Multi-Cabeza de Rango Bajo (MLRA, por sus siglas en inglés) \citep{liu_mlra_2026} refactoriza el latente único de MLA en $L$ sub-cabezas disjuntas cada una de rango $r/L$, de modo que el caché KV puede particionarse entre $L$ dispositivos para una decodificación tensor-paralela de 4 vías (cada dispositivo carga sólo $1/L$ del caché). Las up-proyecciones por sub-cabeza reconstruyen la porción correspondiente de las cabezas completas y se concatenan, logrando una aceleración de decodificación de 2.8$\times$ sobre MLA preservando el presupuesto latente global.

\subsubsection{Formulación Matemática}
\[
c_{KV} = [c_1, c_2, \ldots, c_L], \quad c_i \in \mathbb{R}^{r/L}, \qquad c_{KV} = W_{DKV}\, x.
\]
\[
k_i = W_{UK}^{(i)}\, c_i, \quad v_i = W_{UV}^{(i)}\, c_i, \qquad
k = [k_1; \ldots; k_L], \quad v = [v_1; \ldots; v_L].
\]
\[
\text{atencion} = \texttt{softmax}\!\left(\frac{q\, k^\top}{\sqrt{d_h}}\right) v.
\]

\subsubsection{Pseudocódigo Algorítmico}
\begin{algorithmic}[1]
\State \textbf{Entrada:} $x$, número de sub-cabezas $L$, proyecciones por sub-cabeza $W_{UK}^{(i)}, W_{UV}^{(i)}$
\State $c_{KV} \gets W_{DKV}\, x = [c_1, \ldots, c_L]$ \Comment{particionar latente en $L$ bloques}
\For{$i = 1, \ldots, L$ (en paralelo entre $L$ dispositivos)}
    \State $k_i \gets W_{UK}^{(i)}\, c_i$, \quad $v_i \gets W_{UV}^{(i)}\, c_i$
\EndFor
\State $k \gets \texttt{concat}(k_1, \ldots, k_L)$, \quad $v \gets \texttt{concat}(v_1, \ldots, v_L)$
\State $\text{pesos} \gets \texttt{softmax}(q\, k^\top / \sqrt{d_h})$
\State \textbf{Salida:} $\mathbf{Y} = \text{pesos} \cdot v$
\end{algorithmic}

\subsubsection{Características Clave}
\begin{itemize}[leftmargin=*]
    \item \textbf{Complejidad de entrenamiento}: $\mathcal{O}(n^2 \cdot d)$ (atención completa sobre cabezas concatenadas).
    \item \textbf{Complejidad de inferencia}: $\mathcal{O}(n)$ por token, estado $\mathcal{O}(r_{kv} \cdot n)$ particionado entre $L$ dispositivos.
    \item \textbf{Entrenable}: totalmente entrenable; $L$ controla la granularidad de partición.
    \item \textbf{Fortalezas}: aceleración de decodificación de 2.8$\times$; particionamiento limpio tensor-paralelo de 4 vías del caché; preserva el presupuesto latente total.
    \item \textbf{Limitaciones}: requiere que $L$ divida a $r_{kv}$; más matrices de proyección que MLA; all-reduce inter-dispositivo en la salida.
\end{itemize}

\subsection{Atención Tucker}
La Atención Tucker \citep{klein_tucker_2026} unifica MHA, GQA y MLA bajo una única factorización estilo Tucker de los tensores de pesos de consulta/clave/valor con rangos independientes $q_{\text{rank}}, k_{\text{rank}}, v_{\text{rank}}$. Se recupera MHA cuando todos los rangos igualan el tamaño oculto, GQA cuando $k_{\text{rank}} = v_{\text{rank}} < d$, y MLA cuando $k_{\text{rank}} = v_{\text{rank}}$ es igual al rango latente. Como la factorización es una repas lineal de pesos, el método es totalmente compatible con FlashAttention y RoPE, y produce un orden de magnitud menos parámetros para métricas comparables.

\subsubsection{Formulación Matemática}
\[
\mathbf{Q} = W_{Q,\text{core}}\, W_{Q,\text{factor}}\, x, \qquad
\mathbf{K} = W_{K,\text{core}}\, W_{K,\text{factor}}\, x, \qquad
\mathbf{V} = W_{V,\text{core}}\, W_{V,\text{factor}}\, x.
\]
\[
\text{atencion} = \texttt{softmax}\!\left(\frac{\mathbf{Q}\, \mathbf{K}^\top}{\sqrt{d_h}}\right) \mathbf{V}.
\]
Casos especiales: MHA $\Leftrightarrow q_{\text{rank}} = k_{\text{rank}} = v_{\text{rank}} = d$; GQA $\Leftrightarrow k_{\text{rank}} = v_{\text{rank}} < d$; MLA $\Leftrightarrow k_{\text{rank}} = v_{\text{rank}} = r_{kv}$.

\subsubsection{Pseudocódigo Algorítmico}
\begin{algorithmic}[1]
\State \textbf{Entrada:} $x$, rangos $q_{\text{rank}}, k_{\text{rank}}, v_{\text{rank}}$, matrices core/factor
\State $\mathbf{Q} \gets W_{Q,\text{core}}\, W_{Q,\text{factor}}\, x$ \Comment{consulta factorizada de rango $q_{\text{rank}}$}
\State $\mathbf{K} \gets W_{K,\text{core}}\, W_{K,\text{factor}}\, x$ \Comment{clave factorizada de rango $k_{\text{rank}}$}
\State $\mathbf{V} \gets W_{V,\text{core}}\, W_{V,\text{factor}}\, x$ \Comment{valor factorizado de rango $v_{\text{rank}}$}
\State \textbf{Aplicar RoPE a $\mathbf{Q}, \mathbf{K}$ si se usa codificación posicional}
\State $\text{pesos} \gets \texttt{softmax}(\mathbf{Q}\, \mathbf{K}^\top / \sqrt{d_h})$ \Comment{compatible con FlashAttention}
\State \textbf{Salida:} $\mathbf{Y} = \text{pesos} \cdot \mathbf{V}$
\end{algorithmic}

\subsubsection{Características Clave}
\begin{itemize}[leftmargin=*]
    \item \textbf{Complejidad de entrenamiento}: $\mathcal{O}(n^2 \cdot d)$ sobre los tensores reconstruidos.
    \item \textbf{Complejidad de inferencia}: $\mathcal{O}(n)$ por token; tamaño de caché gobernado por $k_{\text{rank}}$.
    \item \textbf{Entrenable}: los rangos son hiperparámetros; la factorización es una repas de pesos directa.
    \item \textbf{Fortalezas}: unifica MHA/GQA/MLA; reducción de parámetros de un orden de magnitud con métricas equivalentes; compatible con FlashAttention/RoPE.
    \item \textbf{Limitaciones}: dos matmuls por proyección añaden latencia; la selección de rango es un problema de ajuste por tarea.
\end{itemize}

\subsection{Atención de Cabezas Entrelazadas (IHA)}
La Atención de Cabezas Entrelazadas (IHA, por sus siglas en inglés) \citep{duvvuri_iha_2026} construye $P$ pseudo-cabezas por cabeza original (típicamente $P = H$) donde cada consulta/clave/valor pseudo es una combinación lineal aprendida de todas las $H$ cabezas originales. Esto induce hasta $P^2$ patrones de atención distintos por cabeza con una sobrecarga de $\mathcal{O}(H^2 P)$. Teóricamente, IHA necesita sólo $\Theta(\sqrt{k}\, n^2)$ parámetros frente a $\Theta(k\, n^2)$ de MHA en la tarea Polinomial; empíricamente entrega $+10$--$20\%$ en RULER de recuperación multi-clave, $+5.8\%$ en GSM8K y $+2.8\%$ en MATH-500.

\subsubsection{Formulación Matemática}
\[
q_{\text{pseudo}}[p] = \sum_{h=1}^{H} \text{mix}_q[p, h]\, q[h], \quad
k_{\text{pseudo}}[p] = \sum_{h} \text{mix}_k[p, h]\, k[h], \quad
v_{\text{pseudo}}[p] = \sum_{h} \text{mix}_v[p, h]\, v[h].
\]
\[
\text{atencion}_p = \texttt{softmax}\!\left(\frac{q_{\text{pseudo}}[p]\, k_{\text{pseudo}}[p]^\top}{\sqrt{d_h}}\right) v_{\text{pseudo}}[p],
\qquad
\mathbf{Y}[h] = \frac{1}{P} \sum_{p} \text{atencion}_p[h].
\]

\subsubsection{Pseudocódigo Algorítmico}
\begin{algorithmic}[1]
\State \textbf{Entrada:} por cabeza $q[h], k[h], v[h]$, matrices de mezcla $\text{mix}_q, \text{mix}_k, \text{mix}_v \in \mathbb{R}^{P \times H}$
\For{$p = 1, \ldots, P$}
    \State $q_p \gets \sum_h \text{mix}_q[p,h]\, q[h]$ \Comment{mezcla aprendida entre cabezas}
    \State $k_p \gets \sum_h \text{mix}_k[p,h]\, k[h]$, \quad $v_p \gets \sum_h \text{mix}_v[p,h]\, v[h]$
    \State $\text{atencion}_p \gets \texttt{softmax}(q_p\, k_p^\top / \sqrt{d_h})\, v_p$
\EndFor
\State \textbf{Salida:} $\mathbf{Y}[h] \gets \frac{1}{P} \sum_p \text{atencion}_p[h]$ \Comment{promediar salidas pseudo-cabeza por cabeza}
\end{algorithmic}

\subsubsection{Características Clave}
\begin{itemize}[leftmargin=*]
    \item \textbf{Complejidad de entrenamiento}: $\mathcal{O}(n^2 \cdot H \cdot P)$ con $P \approx H$, es decir $\mathcal{O}(n^2 \cdot H^2)$.
    \item \textbf{Complejidad de inferencia}: $\mathcal{O}(n)$ por token con caché KV; sobrecarga extra de mezcla $\mathcal{O}(H^2 P)$.
    \item \textbf{Entrenable}: las matrices de mezcla se aprenden; $P$ es un hiperparámetro.
    \item \textbf{Fortalezas}: eficiencia de parámetros $\Theta(\sqrt{k})$ en Polinomial; grandes ganancias en recuperación y razonamiento.
    \item \textbf{Limitaciones}: costo de mezcla cuadrático en cabezas; más patrones de atención a computar; latencia añadida con $H$ grande.
\end{itemize}

\subsection{Atención laTenT por Cabezas Agrupadas (GTA)}
La Atención laTenT por Cabezas Agrupadas (GTA, por sus siglas en inglés) \citep{sun_gta_2025} combina dos compresiones: (1) un \emph{mapa de atención compartido} --- un tensor de puntajes de atención por grupo de cabezas, reutilizado entre todas las cabezas del grupo, reduciendo el caché de claves; y (2) un \emph{decodificador de valores no lineal} --- el caché de valores se comprime a un latente mediante una down-proyección y se reconstruye por una up-proyección con compuerta SiLU antes de la proyección de salida. GTA reduce los FLOPs de atención hasta 62.5\% frente a GQA, encoge el caché KV hasta 70\%, y produce una aceleración de inferencia de 2$\times$ de extremo a extremo.

\subsubsection{Formulación Matemática}
\[
q_g = \texttt{mean}(q[\text{grupo } g]), \quad k_g = \texttt{mean}(k[\text{grupo } g]), \quad
\text{atencion}_g = \texttt{softmax}\!\left(\frac{q_g\, k_g^\top}{\sqrt{d_h}}\right) \text{ (reutilizado en el grupo)}.
\]
\[
v_{\text{latente}} = W_{DV}\, x, \qquad v = \text{SiLU}(W_{UV}\, v_{\text{latente}}), \qquad
\mathbf{Y} = \text{atencion}_g \cdot v.
\]

\subsubsection{Pseudocódigo Algorítmico}
\begin{algorithmic}[1]
\State \textbf{Entrada:} $q, k$ por cabeza, oculto $x$, partición de grupos $\{g\}$
\State $v_{\text{latente}} \gets W_{DV}\, x$ \Comment{comprimir valores a latente (en caché)}
\State $v \gets \text{SiLU}(W_{UV}\, v_{\text{latente}})$ \Comment{decodificación no lineal de valores}
\For{cada grupo $g$}
    \State $q_g \gets \texttt{mean}(q[g])$, \quad $k_g \gets \texttt{mean}(k[g])$ \Comment{consultas/claves compartidas por grupo}
    \State $\text{atencion}_g \gets \texttt{softmax}(q_g\, k_g^\top / \sqrt{d_h})$
    \State $\mathbf{Y}[g] \gets \text{atencion}_g \cdot v[g]$ \Comment{un mapa reutilizado para todas las cabezas en $g$}
\EndFor
\State \textbf{Salida:} $\mathbf{Y}$
\end{algorithmic}

\subsubsection{Características Clave}
\begin{itemize}[leftmargin=*]
    \item \textbf{Complejidad de entrenamiento}: $\mathcal{O}(n^2 \cdot G \cdot d_h)$ donde $G$ es el número de grupos (frente a $H$ cabezas en GQA).
    \item \textbf{Complejidad de inferencia}: $\mathcal{O}(n)$ por token; caché de valores $\mathcal{O}(r_v \cdot n)$, caché de claves reducido por compartición de grupo.
    \item \textbf{Entrenable}: la asignación de grupo y el rango latente son entrenables/hiperparámetros.
    \item \textbf{Fortalezas}: hasta 62.5\% de reducción de FLOPs de atención, 70\% de encogimiento de caché KV, 2$\times$ de aceleración de extremo a extremo.
    \item \textbf{Limitaciones}: el mapa compartido reduce la diversidad por cabeza; el decodificador no lineal añade latencia; el agrupamiento debe elegirse con cuidado.
\end{itemize}

\subsection{Atención Latente Temporal Multi-Cabeza (MTLA)}
La Atención Latente Temporal Multi-Cabeza (MTLA, por sus siglas en inglés) \citep{deng_mtla_2025} extiende MLA a lo largo del eje temporal: una hyper-red fusiona dinámicamente $m$ vectores temporalmente adyacentes del caché KV en un único slot, reduciendo la longitud temporal por un factor de $m$. Una máscara causal consciente del stride mantiene el entrenamiento paralelo consistente con la inferencia autorregresiva, produciendo una aceleración de decodificación de 5.3$\times$ y una reducción de memoria GPU de 8.3$\times$ frente a MHA en traducción de voz En--De.

\subsubsection{Formulación Matemática}
\[
c_{KV}^{(t)} = W_{DKV}\, x^{(t)} \in \mathbb{R}^{r_{kv}} \text{ (latente por token)}.
\]
\[
\bar c_{KV}^{(j)} = \sum_{i=0}^{m-1} \alpha_i\, c_{KV}^{(s j + i)}, \qquad \alpha_i = \texttt{sigmoid}(\text{HyperNet}(c_{KV}^{(s j + i)})),
\]
donde $s$ es el stride y $m$ la ventana de fusión. Las consultas atienden sobre slots fusionados $\bar c_{KV}$ con una máscara causal consciente del stride $M_{\text{stride}}$:
\[
\text{atencion} = \texttt{softmax}\!\left(\frac{q\, \bar k^\top}{\sqrt{d_h}} + M_{\text{stride}}\right) \bar v, \quad \bar k = W_{UK}\, \bar c_{KV},\ \bar v = W_{UV}\, \bar c_{KV}.
\]

\subsubsection{Pseudocódigo Algorítmico}
\begin{algorithmic}[1]
\State \textbf{Entrada:} latentes por token $c_{KV}^{(t)}$, ventana de fusión $m$, stride $s$
\State \textbf{Fase de fusión (hyper-red):}
\For{cada slot fusionado $j = 0, 1, \ldots$}
    \State $\alpha_i \gets \texttt{sigmoid}(\text{HyperNet}(c_{KV}^{(s j + i)}))$ para $i = 0, \ldots, m-1$
    \State $\bar c_{KV}^{(j)} \gets \sum_{i=0}^{m-1} \alpha_i\, c_{KV}^{(s j + i)}$ \Comment{promedio con compuerta sobre la ventana}
\EndFor
\State $\bar k \gets W_{UK}\, \bar c_{KV}$, \quad $\bar v \gets W_{UV}\, \bar c_{KV}$
\State $\text{pesos} \gets \texttt{softmax}(q\, \bar k^\top / \sqrt{d_h} + M_{\text{stride}})$ \Comment{máscara causal consciente del stride}
\State \textbf{Salida:} $\mathbf{Y} = \text{pesos} \cdot \bar v$
\end{algorithmic}

\subsubsection{Características Clave}
\begin{itemize}[leftmargin=*]
    \item \textbf{Complejidad de entrenamiento}: $\mathcal{O}(n^2 \cdot d / m)$ sobre la secuencia fusionada (longitud $\lceil n/m \rceil$) más costo de la hyper-red.
    \item \textbf{Complejidad de inferencia}: $\mathcal{O}(n/m)$ por token sobre slots fusionados; estado $\mathcal{O}(r_{kv} \cdot n / m)$.
    \item \textbf{Entrenable}: la hyper-red y la ventana de fusión $m$ se aprenden/eliges; la máscara consciente del stride garantiza consistencia.
    \item \textbf{Fortalezas}: aceleración de decodificación de 5.3$\times$, reducción de memoria GPU de 8.3$\times$ frente a MHA; compresión temporal ortogonal a la compresión de rango.
    \item \textbf{Limitaciones}: la ventana de fusión puede difuminar detalle temporal fino; la hyper-red añade cómputo; el diseño de la máscara no es trivial.
\end{itemize}

\subsection{Comparación y Síntesis Arquitectónica}

\small
\begin{longtable}{p{2.0cm}p{1.8cm}p{1.6cm}p{1.6cm}p{4.5cm}}
\toprule
\textbf{Arquitectura} & \textbf{Entrenar} & \textbf{Inferir} & \textbf{Estado} & \textbf{Característica Clave} \\
\midrule
\endhead
Atención Estándar & $O(n^2 d)$ & $O(n)$ caché & $O(nd)$ & Expresividad perfecta, enrutamiento completo de tokens, cuello de botella KV. \\
GQA & $O(n^2 d)$ & $O(n)$ caché & $O(h_k \cdot n \cdot d_h)$ & Compresión KV ajustable, probado a escala, reducción de caché en factor $G \times$. \\
Atención Sigmoide & $O(n^2 d)$ & $O(n)$ caché & $O(nd)$ & Compuerta elemento a elemento, eficiente en hardware, estable a escala con norma híbrida. \\
RetNet & $O(n^2 d)$ & $O(1)$ & $O(d^2)$ & Decaimiento multi-escala, triple modo de cómputo, patrón de olvido fijo. \\
Mamba & $O(nd)$ & $O(1)$ & $O(d)$ & Dinámica selectiva dependiente de la entrada, entrenamiento e inferencia lineales, escaneo de hardware. \\
Transformer EDO & $O(kn^2 d)$ & $O(kn)$ & $O(nd)$ & Refinamiento por integración numérica, uso compartido de pesos por etapas, mayor precisión. \\
Titans & $O(c^2 + nf)$ & $O(c^2 + 1)$ & $O(1)$ & Adaptación de memoria en tiempo de prueba, contexto extremo, actualizaciones de parámetros durante inferencia. \\
MLA & $O(n^2 d)$ & $O(n)$ caché & $O(r_{kv} n)$ & KV latente de rango $r_{kv}\!\approx\!d/2$, reduce a la mitad el caché, iguala calidad MHA. \\
GQLA & $O(n^2 d)$ & $O(n)$ caché & $O(r_{kv} n)$ & Mismos pesos MLA, dos rutas de decodificación algebraicamente equivalentes adaptativas al hardware. \\
MLRA & $O(n^2 d)$ & $O(n)$ caché & $O(r_{kv} n)$ particionado & Latente dividido en $L$ sub-cabezas para decodificación tensor-paralela de 4 vías, 2.8$\times$ de aceleración. \\
Tucker & $O(n^2 d)$ & $O(n)$ caché & $O(k_{\text{rank}} n)$ & Factorización Tucker unifica MHA/GQA/MLA, un orden de magnitud menos parámetros. \\
IHA & $O(n^2 H^2)$ & $O(n)$ caché & $O(H n d_h)$ & Pseudo-cabezas como mezclas aprendidas de cabezas, $\Theta(\sqrt{k})$ parámetros, +10--20\% RULER. \\
GTA & $O(n^2 G d_h)$ & $O(n)$ caché & $O(r_v n)$ & Mapa compartido por grupo + decodificador de valores SiLU, 62.5\% menos FLOPs, 70\% menos caché. \\
MTLA & $O(n^2 d / m)$ & $O(n/m)$ & $O(r_{kv} n/m)$ & Fusión temporal de $m$ slots KV, 5.3$\times$ decodificación, 8.3$\times$ reducción de memoria. \\
\bottomrule
\end{longtable}
\normalsize

El campo exhibe una progresión clara a lo largo de dos ejes: (i) \textbf{eficiencia computacional}, pasando de $O(n^2)$ a $O(n)$ u $O(1)$, y (ii) \textbf{adaptabilidad de memoria}, pasando de pesos estáticos a modelos dinámicos actualizados en tiempo de prueba. La atención estándar sigue siendo la línea base de expresividad; Mamba y RetNet representan la frontera práctica de eficiencia; Titans introduce un eje de innovación ortogonal (aprendizaje en tiempo de prueba). La elección de arquitectura refleja la restricción fundamental de ingeniería: expresividad versus costo de despliegue.


\section{Annex C: Comprehensive Sparse Attention Mechanisms}
\label{sec:annex-sparse}

\subsection{Resumen Ejecutivo}

Los mecanismos de atención dispersa abordan la prohibitiva complejidad computacional y de memoria $\mathcal{O}(n^2)$ de la atención de producto punto escalado estándar al restringir el conjunto de pares clave-valor atendidos. La atención dispersa moderna abarca un rico espacio de diseño distinguido por cuatro dimensiones ortogonales: (i) \textbf{patrón de dispersidad} (geométrico fijo, dependiente de datos o basado en frecuencia), (ii) \textbf{capacidad de entrenamiento} (entrenado de extremo a extremo o solo inferencia), (iii) \textbf{unidad de dispersidad} (tokens individuales, ventanas locales o bloques), y (iv) \textbf{mecanismo} (enmascaramiento, selección, filtrado o compresión). Este anexo sintetiza siete familias principales de atención dispersa---Sparse Transformer, Longformer, BigBird, SparseK, NSA, SpargeAttn y FASA---proporcionando formulaciones matemáticas, pseudocódigo algorítmico y comparaciones arquitectónicas exhaustivas.

\subsection{Sparse Transformer: Patrones Estrideados y Fijos Factorizados}

\subsubsection{Formulación Matemática}

Sparse Transformer \citep{child_sparse_transformer_2019} reduce la complejidad cuadrática a $\mathcal{O}(n\sqrt{n})$ factorizando la atención dispersa en dos cabezas dispersas complementarias. Sea $n$ la longitud de la secuencia y $l = \lfloor\sqrt{n}\rfloor$ el stride.

\paragraph{Cabeza de atención estrideada}: Cada posición $i$ atiende a cada $l$-ésima posición anterior:
\[
\mathcal{A}_i^{(\text{estrideada})} = \{j : j \leq i,\; (i - j) \bmod l = 0\}
\]

\paragraph{Cabeza de atención fija}: Cada posición $i$ atiende a posiciones dentro de su bloque local más columnas de resumen fijas:
\[
\mathcal{A}_i^{(\text{fija})} = \{j : \lfloor j/l \rfloor = \lfloor i/l \rfloor\} \cup \{j : j \bmod l \in \{l-c, \ldots, l-1\}\}
\]
donde $c$ es un hiperparámetro que controla el número de columnas de resumen (típicamente $c=1$ o $c=2$).

El cálculo de atención para cada cabeza $h$ sigue las reglas estándar de producto punto escalado restringido al conjunto disperso:
\[
\text{Attn}_h(Q, K, V)_i = \text{softmax}\left(\frac{q_i^h K_{\mathcal{A}_i}^h\top}{\sqrt{d_k}}\right) V_{\mathcal{A}_i}^h
\]

La idea clave es que con las dos cabezas factorizadas operando en paralelo a través de una profundidad de $L$ capas transformer, cada posición puede alcanzar cualquier otra posición a través de un camino de longitud máxima $L + 1$, preservando la alcanzabilidad de largo alcance con una fracción del costo de la atención densa.

\subsubsection{Pseudocódigo Algorítmico}

\begin{algorithm}[H]
\caption{Atención Sparse Transformer: Cabezas Duales Estrideada + Fija}
\begin{algorithmic}[1]
\Require Consulta $\mathbf{Q} \in \mathbb{R}^{n \times d}$, Clave $\mathbf{K} \in \mathbb{R}^{n \times d}$, Valor $\mathbf{V} \in \mathbb{R}^{n \times d}$
\Require Stride $l = \lfloor\sqrt{n}\rfloor$, columnas de resumen $c \in \{1, 2\}$
\Ensure Salida $\mathbf{Y} \in \mathbb{R}^{n \times d}$
\State \textbf{Cabeza 1: Atención Estrideada}
\For{$i = 1, \ldots, n$}
    \State $\mathcal{A}_i \gets \{j : j \leq i \text{ y } (i-j) \bmod l = 0\}$ \Comment{Cada $l$-ésima posición}
    \State $\text{puntajes}_{i} \gets \mathbf{q}_i \mathbf{K}_{\mathcal{A}_i}^\top / \sqrt{d_k}$
    \State $\text{attn}_{i} \gets \text{softmax}(\text{puntajes}_{i})$
    \State $\mathbf{y}_i^{(1)} \gets \text{attn}_{i} \mathbf{V}_{\mathcal{A}_i}$
\EndFor
\State \textbf{Cabeza 2: Atención de Bloque Fijo + Resumen}
\For{$i = 1, \ldots, n$}
    \State $\text{inicio\_bloque} \gets \lfloor i/l \rfloor \cdot l$, $\text{fin\_bloque} \gets \min(i+1, \text{inicio\_bloque}+l)$
    \State $\mathcal{A}_i \gets [\text{inicio\_bloque}, \text{fin\_bloque}) \cup \{\text{columnas de las últimas } c \text{ columnas}\}$
    \State $\text{puntajes}_{i} \gets \mathbf{q}_i \mathbf{K}_{\mathcal{A}_i}^\top / \sqrt{d_k}$
    \State $\text{attn}_{i} \gets \text{softmax}(\text{puntajes}_{i})$
    \State $\mathbf{y}_i^{(2)} \gets \text{attn}_{i} \mathbf{V}_{\mathcal{A}_i}$
\EndFor
\State Concatenar: $\mathbf{Y} = [\mathbf{y}^{(1)} \| \mathbf{y}^{(2)}]$ y proyectar mediante capa lineal de salida
\Return $\mathbf{Y}$
\end{algorithmic}
\end{algorithm}

\subsubsection{Características Clave}

\begin{itemize}[leftmargin=*]
    \item \textbf{Complejidad}: $\mathcal{O}(n\sqrt{n} \cdot d)$ durante el entrenamiento; $\mathcal{O}(n)$ por token durante la generación.
    \item \textbf{Entrenable}: Sí; retropropagación completa a través de las posiciones seleccionadas.
    \item \textbf{Patrón de dispersidad}: Fijo e independiente de los datos; los patrones no se adaptan al contenido.
    \item \textbf{Fortaleza}: Alcanzabilidad estructurada que garantiza dependencias de largo alcance; demostrado efectivo en secuencias largas (Enwik8, imágenes de texto).
    \item \textbf{Limitación}: Los patrones fijos pueden pasar por alto relaciones importantes pero no contiguas; requiere kernels CUDA personalizados para eficiencia práctica.
\end{itemize}

\subsection{Longformer: Ventana Deslizante, Dilatación y Tokens Globales}

\subsubsection{Formulación Matemática}

Longformer \citep{beltagy_longformer_2020} logra una complejidad lineal $\mathcal{O}(n \cdot w)$ combinando tres patrones de atención complementarios.

\paragraph{Atención de ventana deslizante}: Vecindario local de tamaño $w$:
\[
\mathcal{A}_i^{(\text{ventana})} = \{j : |i - j| \leq w/2\}
\]

\paragraph{Ventana deslizante dilatada}: Atención con ventana y dilatación $d$, saltando stride $d+1$:
\[
\mathcal{A}_i^{(\text{dilatada})} = \{j : |i - j| \leq (w/2)(d+1) \text{ y } (i-j) \bmod (d+1) = 0\}
\]

\paragraph{Atención global}: Tokens globales designados (ej., $[\text{CLS}]$) atienden a todas las posiciones y todos atienden a ellos:
\[
\mathcal{A}_i^{(\text{global})} = \begin{cases}
\{1, \ldots, n\} & \text{si } i \in \mathcal{G} \\
\mathcal{A}_i^{(\text{ventana})} \cup \mathcal{G} & \text{en otro caso}
\end{cases}
\]

Los tres patrones se aplican mediante proyecciones separadas. La atención local y global pueden usar las mismas proyecciones o unas separadas específicas para cada tarea.

\subsubsection{Pseudocódigo Algorítmico}

\begin{algorithm}[H]
\caption{Longformer: Ventana Deslizante + Dilatada + Global}
\begin{algorithmic}[1]
\Require Consulta $\mathbf{Q} \in \mathbb{R}^{n \times d}$, Clave $\mathbf{K}$, Valor $\mathbf{V}$, tamaño de ventana $w$, dilatación $d$, índices de tokens globales $\mathcal{G}$
\Ensure Salida $\mathbf{Y}$
\For{$i = 1, \ldots, n$}
    \If{$i \in \mathcal{G}$}
        \State $\mathcal{A}_i \gets \{1, \ldots, n\}$ \Comment{Token global atiende a todos}
    \Else
        \State Ventana local: $\mathcal{A}^{(\text{local})} \gets [i - w/2, i + w/2] \cap [1, n]$
        \State Desplazamientos dilatados: $\mathcal{A}^{(\text{dilatada})} \gets \{j : (i-j) \bmod (d+1) = 0\} \cap \mathcal{A}^{(\text{rango dilatado})}$
        \State Combinar: $\mathcal{A}_i \gets \mathcal{A}^{(\text{local})} \cup \mathcal{A}^{(\text{dilatada})} \cup \mathcal{G}$
    \EndIf
    \State $\text{puntajes}_i \gets \mathbf{q}_i \mathbf{K}_{\mathcal{A}_i}^\top / \sqrt{d_k}$
    \State $\text{attn}_i \gets \text{softmax}(\text{puntajes}_i)$
    \State $\mathbf{y}_i \gets \text{attn}_i \mathbf{V}_{\mathcal{A}_i}$
\EndFor
\Return $\mathbf{Y}$
\end{algorithmic}
\end{algorithm}

\subsubsection{Características Clave}

\begin{itemize}[leftmargin=*]
    \item \textbf{Complejidad}: $\mathcal{O}(n \cdot w)$ donde $w$ es el tamaño de ventana (lineal cuando $w \ll n$).
    \item \textbf{Entrenable}: Sí; reemplazo directo para la atención estándar.
    \item \textbf{Cobertura}: Con $L$ capas y tamaño de ventana $w$, el campo receptivo crece a $L \times w$, cubriendo toda la secuencia con redes poco profundas.
    \item \textbf{Fortaleza}: Práctico para tareas a nivel de documentos; la dilatación expande los campos receptivos locales con sobrecarga mínima; los tokens globales proporcionan correspondencia consulta-documento.
    \item \textbf{Limitación}: El tamaño de ventana sigue siendo un hiperparámetro de diseño; subóptimo para tareas que requieren patrones de atención densos.
\end{itemize}

\subsection{BigBird: Grafo Disperso Aleatorio, Local y Global}

\subsubsection{Formulación Matemática}

BigBird \citep{zaheer_bigbird_2020} preserva la aproximación universal y la completitud de Turing utilizando una combinación fundamentada de tres patrones de dispersidad.

\paragraph{Conexiones aleatorias}: Cada posición se conecta a $r$ posiciones muestreadas aleatoriamente:
\[
\mathcal{A}_i^{(\text{aleatoria})} = \text{MuestraAleatoria}(\{1, \ldots, n\}, r)
\]

\paragraph{Ventana local}: Vecindario de ventana deslizante:
\[
\mathcal{A}_i^{(\text{local})} = \{j : |i-j| \leq w/2\}
\]

\paragraph{Tokens globales}: Anclas globales específicas de la tarea:
\[
\mathcal{A}_i^{(\text{global})} = \begin{cases}
\{1, \ldots, n\} & \text{si } i \in \mathcal{G} \\
\mathcal{A}_i^{(\text{aleatoria})} \cup \mathcal{A}_i^{(\text{local})} \cup \mathcal{G} & \text{en otro caso}
\end{cases}
\]

La máscara dispersa compuesta $M$ es:
\[
M_{ij} = \mathbf{1}[j \in \mathcal{A}_i^{(\text{aleatoria})} \cup \mathcal{A}_i^{(\text{local})} \cup \mathcal{A}_i^{(\text{global})}]
\]

En la práctica, BigBird agrupa tokens en bloques de tamaño $b$ y aplica la decisión de dispersidad a nivel de bloque, permitiendo implementaciones eficientes de bloques dispersos.

\subsubsection{Garantía Teórica}

Los autores demuestran que con $O(1)$ conexiones aleatorias y tokens globales, el grafo disperso mantiene las mismas propiedades de aproximación universal y completitud de Turing que la atención densa. Formalmente, cualquier función computable puede representarse y cualquier secuencia de operaciones puede simularse dentro del grafo de conectividad de BigBird.

\subsubsection{Pseudocódigo Algorítmico}

\begin{algorithm}[H]
\caption{BigBird: Atención Dispersa por Bloques Aleatoria + Local + Global}
\begin{algorithmic}[1]
\Require Consulta $\mathbf{Q}$, Clave $\mathbf{K}$, Valor $\mathbf{V}$, tamaño de bloque $b$, enlaces aleatorios por bloque $r$, índices globales $\mathcal{G}$
\Ensure Salida $\mathbf{Y}$
\State Reformar en bloques: $n_b = \lceil n / b \rceil$ bloques
\For{bloque $i = 0, \ldots, n_b-1$}
    \For{posición $j$ en bloque $i$}
        \State \textbf{Ventana local}: $\mathcal{A}_j^{(\text{local})} = \text{posiciones cercanas en el bloque} \cup \text{posiciones frontera}$
        \State \textbf{Muestreo aleatorio de bloques}: Muestrear $r$ bloques uniformemente al azar, agregar todas las posiciones de esos bloques
        \State \textbf{Global}: Agregar todas las posiciones de tokens globales
        \State $\mathcal{A}_j \gets \mathcal{A}_j^{(\text{local})} \cup \mathcal{A}_j^{(\text{aleatoria})} \cup \mathcal{G}$
    \EndFor
\EndFor
\For{$i = 1, \ldots, n$}
    \State $\text{puntajes}_i \gets \mathbf{q}_i \mathbf{K}_{\mathcal{A}_i}^\top / \sqrt{d_k}$
    \State $\text{attn}_i \gets \text{softmax}(\text{puntajes}_i)$
    \State $\mathbf{y}_i \gets \text{attn}_i \mathbf{V}_{\mathcal{A}_i}$
\EndFor
\Return $\mathbf{Y}$
\end{algorithmic}
\end{algorithm}

\subsubsection{Características Clave}

\begin{itemize}[leftmargin=*]
    \item \textbf{Complejidad}: $\mathcal{O}(n)$ o casi lineal en implementación óptima de bloques dispersos.
    \item \textbf{Entrenable}: Sí.
    \item \textbf{Base teórica}: Mantiene demostrablemente la aproximación universal y la completitud de Turing con conectividad dispersa.
    \item \textbf{Fortaleza}: Rendimiento sólido en documentos largos para preguntas-respuestas, resúmenes y tareas de genómica.
    \item \textbf{Limitación}: El muestreo aleatorio introduce varianza y no determinismo; el ajuste de los hiperparámetros $r, w, g$ sigue dependiendo de la tarea.
\end{itemize}

\subsection{Atención SparseK: Selección Diferencial Top-$k$}

\subsubsection{Formulación Matemática}

SparseK \citep{lou_sparsek_2024} permite dispersidad aprendible mediante un \textbf{operador top-$k$ diferenciable}. Una red de puntuación $\phi_\theta$ evalúa la importancia de cada par clave-valor:

\[
u_j = \phi_\theta(\mathbf{k}_j, \mathbf{q}_i) \in \mathbb{R}
\]

El \textbf{operador SparseK} selecciona los puntajes top-$k$ manteniéndose diferenciable. Calcula un umbral $\tau(u)$ tal que la suma de los puntajes activos sea igual a $k$:

\[
\text{SparseK}(u, k)_j = \max(u_j - \tau(u), 0), \quad \text{donde} \quad \sum_j \max(u_j - \tau, 0) = k
\]

El umbral $\tau$ se encuentra mediante bisección. La atención se convierte en:

\[
m_j = \text{SparseK}(u, k)_j, \quad \text{Attn}_i = \text{softmax}\left(\frac{\mathbf{q}_i K_{\text{sel}}^\top}{\sqrt{d_k}}\right) V_{\text{sel}}
\]

donde $K_{\text{sel}}, V_{\text{sel}}$ contienen solo las entradas top-$k$ (aquellas con $m_j > 0$).

\subsubsection{Pseudocódigo Algorítmico}

\begin{algorithm}[H]
\caption{SparseK: Atención Top-$k$ Diferenciable}
\begin{algorithmic}[1]
\Require Consulta $\mathbf{q} \in \mathbb{R}^d$, Matriz de clave $\mathbf{K} \in \mathbb{R}^{n \times d}$, Matriz de valor $\mathbf{V} \in \mathbb{R}^{n \times d}$
\Require Red de puntuación $\phi_\theta$, dispersidad objetivo $k$
\Ensure Salida de atención $\mathbf{y}$
\State \textbf{Calcular puntajes de importancia:}
\For{$j = 1, \ldots, n$}
    \State $u_j \gets \phi_\theta(\mathbf{k}_j, \mathbf{q})$
\EndFor
\State \textbf{Calcular top-$k$ diferenciable mediante umbral:}
\State $u_{\text{ordenado}} \gets \text{ordenar}(u, \text{descendente}=\text{Verdadero})$
\State $\text{suma\_acum} \gets \text{suma\_acumulada}(u_{\text{ordenado}})$
\State Encontrar $\rho = \max\{i : u_{\text{ordenado}}[i] > 0 \text{ y suma\_acum}[i] \leq k\}$
\State $\tau \gets (u_{\text{ordenado}}[\rho] - k) / (\rho + 1)$ \Comment{Umbral para $k$ elementos activos}
\State $m \gets \max(u - \tau, 0)$ \Comment{Máscara de selección diferenciable}
\State \textbf{Aplicar máscara y calcular atención:}
\State $K_{\text{sel}} \gets K[m > 0]$, $V_{\text{sel}} \gets V[m > 0]$
\State $\text{puntajes} \gets \mathbf{q} K_{\text{sel}}^\top / \sqrt{d_k}$
\State $\text{attn} \gets \text{softmax}(\text{puntajes})$
\State $\mathbf{y} \gets \text{attn} \cdot V_{\text{sel}}$
\Return $\mathbf{y}$
\end{algorithmic}
\end{algorithm}

\subsubsection{Características Clave}

\begin{itemize}[leftmargin=*]
    \item \textbf{Complejidad}: $\mathcal{O}(n \cdot d)$ durante el entrenamiento (lineal en la longitud de la secuencia); $\mathcal{O}(k)$ por token durante la generación.
    \item \textbf{Entrenable}: Sí; flujo de gradiente de extremo a extremo a través del operador SparseK.
    \item \textbf{Generación incremental}: Soporta generación autorregresiva eficiente con memoria constante.
    \item \textbf{Fortaleza}: Se integra perfectamente en arquitecturas LLM existentes; mínimo ajuste fino necesario.
    \item \textbf{Limitación}: El acceso a memoria disperso de la selección top-$k$ puede limitar la eficiencia de caché en algunos hardwares; la red de puntuación añade sobrecarga.
\end{itemize}

\subsection{NSA (Native Sparse Attention): Ramas Jerárquicas Alineadas con el Hardware}

\subsubsection{Formulación Matemática}

NSA \citep{yuan_nsa_2025} descompone la atención dispersa en tres ramas paralelas que se combinan mediante compuertas aprendidas.

\paragraph{Rama de compresión}: Un MLP aprendible $\varphi$ comprime bloques de clave-valor:

\[
\tilde{K}_t^{\text{cmp}} = \left[\varphi(k_{id+1:id+l})\right]_{\substack{1 \leq i \leq \lfloor(t-l)/d\rfloor}} \quad, \quad \tilde{V}_t^{\text{cmp}} = \left[\varphi(v_{id+1:id+l})\right]
\]

\paragraph{Rama de selección}: Los bloques de alta importancia se seleccionan basándose en puntajes comprimidos:

\[
p_t^{\text{cmp}} = \text{softmax}(q_t^\top \tilde{K}_t^{\text{cmp}} / \sqrt{d}), \quad I_t = \mathrm{TopK}(p_t^{\text{cmp}}, n), \quad \tilde{K}_t^{\text{sel}} = \mathrm{Recolectar}(K, I_t)
\]

\paragraph{Rama de ventana deslizante}: Contexto local fijo:

\[
\tilde{K}_t^{\text{win}} = K_{t-w:t}, \quad \tilde{V}_t^{\text{win}} = V_{t-w:t}
\]

\paragraph{Combinación con compuertas}: Las tres salidas de atención se combinan mediante compuertas aprendidas:

\[
\alpha_t^{(\text{cmp})} = \sigma(\text{Lineal}_\text{cmp}(q_t)), \quad \alpha_t^{(\text{sel})} = \sigma(\text{Lineal}_\text{sel}(q_t)), \quad \alpha_t^{(\text{win})} = \sigma(\text{Lineal}_\text{win}(q_t))
\]

\[
y_t = \alpha_t^{(\text{cmp})} \cdot \text{Attn}(q_t, \tilde{K}^{\text{cmp}}, \tilde{V}^{\text{cmp}}) + \alpha_t^{(\text{sel})} \cdot \text{Attn}(q_t, \tilde{K}^{\text{sel}}, \tilde{V}^{\text{sel}}) + \alpha_t^{(\text{win})} \cdot \text{Attn}(q_t, \tilde{K}^{\text{win}}, \tilde{V}^{\text{win}})
\]

\subsubsection{Pseudocódigo Algorítmico}

\begin{algorithm}[H]
\caption{NSA: Ramas Comprimida + Seleccionada + Ventana con Compuertas Aprendidas}
\begin{algorithmic}[1]
\Require Consulta $\mathbf{q}_t$, secuencias de Clave/Valor almacenadas en caché, tamaño de bloque $l$, stride $d$, conteo de selección $n$, ventana $w$
\Require MLP de compresión $\varphi$, redes de compuerta $\text{Lineal}_{\text{cmp}}, \text{Lineal}_{\text{sel}}, \text{Lineal}_{\text{win}}$
\Ensure Salida $\mathbf{y}_t$
\State \textbf{Rama de compresión:}
\For{$i = 1$ to $\lfloor t/d \rfloor$}
    \State $k_i^{\text{cmp}} \gets \varphi(\mathbf{K}_{id+1:id+l})$ \Comment{Comprimir bloque mediante MLP}
    \State $v_i^{\text{cmp}} \gets \varphi(\mathbf{V}_{id+1:id+l})$
\EndFor
\State $\tilde{\mathbf{K}}^{\text{cmp}} \gets [k_1^{\text{cmp}}, \ldots, k_{\lfloor t/d \rfloor}^{\text{cmp}}]$, $\tilde{\mathbf{V}}^{\text{cmp}} \gets [v_1^{\text{cmp}}, \ldots, v_{\lfloor t/d \rfloor}^{\text{cmp}}]$
\State $\mathbf{y}_t^{(\text{cmp})} \gets \text{SDPA}(\mathbf{q}_t, \tilde{\mathbf{K}}^{\text{cmp}}, \tilde{\mathbf{V}}^{\text{cmp}})$
\State \textbf{Rama de selección:}
\State Calcular puntajes sobre comprimidos: $\mathbf{s} \gets \text{softmax}(\mathbf{q}_t \tilde{\mathbf{K}}^{\text{cmp}\top} / \sqrt{d})$
\State Seleccionar los $n$ bloques más importantes: $I \gets \mathrm{TopK}(\mathbf{s}, n)$
\State Recolectar bloques completos: $\tilde{\mathbf{K}}^{\text{sel}} \gets [\mathbf{K}_{I_i d+1:I_i d+l}]_i$, $\tilde{\mathbf{V}}^{\text{sel}} \gets [\mathbf{V}_{I_i d+1:I_i d+l}]_i$
\State $\mathbf{y}_t^{(\text{sel})} \gets \text{SDPA}(\mathbf{q}_t, \tilde{\mathbf{K}}^{\text{sel}}, \tilde{\mathbf{V}}^{\text{sel}})$
\State \textbf{Rama de ventana:}
\State $\mathbf{y}_t^{(\text{win})} \gets \text{SDPA}(\mathbf{q}_t, \mathbf{K}_{t-w:t}, \mathbf{V}_{t-w:t})$
\State \textbf{Fusión con compuertas:}
\State $\alpha^{(\text{cmp})} \gets \sigma(\text{Lineal}_{\text{cmp}}(\mathbf{q}_t))$, $\alpha^{(\text{sel})} \gets \sigma(\text{Lineal}_{\text{sel}}(\mathbf{q}_t))$, $\alpha^{(\text{win})} \gets \sigma(\text{Lineal}_{\text{win}}(\mathbf{q}_t))$
\State $\mathbf{y}_t \gets \alpha^{(\text{cmp})} \mathbf{y}_t^{(\text{cmp})} + \alpha^{(\text{sel})} \mathbf{y}_t^{(\text{sel})} + \alpha^{(\text{win})} \mathbf{y}_t^{(\text{win})}$
\Return $\mathbf{y}_t$
\end{algorithmic}
\end{algorithm}

\subsubsection{Características Clave}

\begin{itemize}[leftmargin=*]
    \item \textbf{Complejidad}: $\mathcal{O}(t/d + n \cdot l + w)$ tokens atendidos por paso (significativamente reducido frente a $\mathcal{O}(t)$ para densa).
    \item \textbf{Entrenable}: Sí; entrenamiento directo con todas las ramas diferenciables.
    \item \textbf{Alineación con hardware}: Diseñado para utilización eficiente de núcleos tensoriales; la compresión y selección reducen la presión sobre el ancho de banda de memoria.
    \item \textbf{Fortaleza}: 11.6$\times$ de aceleración en decodificación y 9$\times$ de aceleración hacia adelante en secuencias de 64k; supera a la atención completa en muchas tareas.
    \item \textbf{Limitación}: Requiere kernels Triton/CUDA personalizados; la arquitectura multi-rama compleja aumenta la sobrecarga de ingeniería.
\end{itemize}

\subsection{FASA: Atención Dispersa Consciente de la Frecuencia}

\subsubsection{Formulación Matemática}

FASA \citep{wang_fasa_2026} es un \textbf{método en tiempo de inferencia sin entrenamiento} que explota la estructura del embedding de posición rotatorio (RoPE). Bajo RoPE, el embedding del token se rota por $\theta_i = B^{-2(i-1)/d}$, descomponiendo el espacio $d$-dimensional en $d/2$ fragmentos de frecuencia (FC).

\paragraph{Idea clave}: Solo un subconjunto pequeño ($ < 1\%$) de los FC importa para la conciencia contextual; la mayoría codifica patrones posicionales.

\paragraph{Identificación de fragmento de frecuencia dominante}: Para cada capa $l$ y cabeza $h$, identificar FC dominantes mediante concordancia contextual (CA):

\[
\text{CA}^{l,h,i} = \frac{|\text{TopK}(\alpha^{l,h}) \cap \text{TopK}(\alpha^{l,h,i})|}{K}
\]

donde $\alpha^{l,h}$ es la máscara de atención completa y $\alpha^{l,h,i}$ es la máscara usando solo el FC $i$. Los FC dominantes tienen CA alta---contribuyen significativamente al patrón de atención final.

\paragraph{Etapa de predicción de importancia de tokens (TIP)}: Usando solo FC dominantes, calcular importancia ligera por token:

\[
S_t^{l,h} = \sum_{i \in \mathcal{I}_{\text{dom}}^{l,h}} \alpha^{l,h,i}(\mathbf{q}_t, \mathbf{K}_{1:t}), \quad \mathcal{T}_t = \text{TopK-Índices}(S_t, N_{\text{fac}})
\]

\paragraph{Etapa de cálculo de atención enfocada (FAC)}: Calcular atención de precisión completa en los tokens seleccionados:

\[
\hat{\alpha}_{\text{FAC}} = \text{softmax}\left(\frac{\mathbf{q}_t \mathbf{K}_{\mathcal{T}_t}^\top}{\sqrt{d}}\right), \quad \mathbf{o}_t = \hat{\alpha}_{\text{FAC}} \mathbf{V}_{\mathcal{T}_t}
\]

\subsubsection{Pseudocódigo Algorítmico}

\begin{algorithm}[H]
\caption{FASA: Atención Dispersa Consciente de la Frecuencia (en Tiempo de Inferencia)}
\begin{algorithmic}[1]
\Require Consulta actual $\mathbf{q}_t$, Clave/Valor almacenados en caché $\mathbf{K}_{1:t}, \mathbf{V}_{1:t}$, base RoPE $B$, FC dominantes $\mathcal{I}_{\text{dom}}$
\Require Presupuesto TIP $N_{\text{tip}}$, presupuesto FAC $N_{\text{fac}}$
\Ensure Salida $\mathbf{o}_t$
\State \textbf{Etapa 1: Predicción de Importancia de Tokens (TIP)}
\For{capa $l$ y cabeza $h$}
    \For{posición $i = 1$ to $t$}
        \State Calcular importancia desde FC dominantes: $s_i \gets 0$
        \For{$fc \in \mathcal{I}_{\text{dom}}^{l,h}$}
            \State Rotar consulta y clave en FC $fc$: $\mathbf{q}^{(fc)}, \mathbf{k}_i^{(fc)}$
            \State $s_i^{(fc)} \gets \text{softmax}(\mathbf{q}^{(fc)} \mathbf{k}_i^{(fc)\top} / \sqrt{d}) $
            \State $s_i \gets s_i + s_i^{(fc)}$
        \EndFor
    \EndFor
    \State $\mathcal{T}_t \gets \mathrm{TopK-Índices}([s_1, \ldots, s_t], N_{\text{fac}})$ \Comment{Seleccionar tokens principales}
\EndFor
\State \textbf{Etapa 2: Cálculo de Atención Enfocada (FAC)}
\State $\text{puntajes}_{\text{fac}} \gets \mathbf{q}_t \mathbf{K}_{\mathcal{T}_t}^\top / \sqrt{d}$ \Comment{Producto punto de precisión completa}
\State $\alpha_{\text{fac}} \gets \text{softmax}(\text{puntajes}_{\text{fac}})$ \Comment{Softmax completo}
\State $\mathbf{o}_t \gets \alpha_{\text{fac}} \mathbf{V}_{\mathcal{T}_t}$ \Comment{Suma ponderada de valores}
\Return $\mathbf{o}_t$
\end{algorithmic}
\end{algorithm}

\subsubsection{Características Clave}

\begin{itemize}[leftmargin=*]
    \item \textbf{Complejidad}: TIP es $\mathcal{O}(t \cdot N_{\text{tip}})$; FAC es $\mathcal{O}(N_{\text{fac}} \cdot d)$.
    \item \textbf{Entrenable}: No; optimización de inferencia sin entrenamiento.
    \item \textbf{Aplicabilidad}: Requiere modelos basados en RoPE; extendido a ALiBi y MLA con modificaciones.
    \item \textbf{Fortaleza}: Precisión cercana a la óptima con $\leq 256$ tokens de entre millones; hasta 2.56$\times$ de aceleración en decodificación; ortogonal a la cuantización.
    \item \textbf{Limitación}: Requiere calibración fuera de línea; la identificación de FC dominantes añade sobrecarga de preprocesamiento.
\end{itemize}

\subsection{SpargeAttn: Filtrado de Dos Etapas a Nivel de Bloques}

\subsubsection{Formulación Matemática}

SpargeAttn \citep{zhang_spargeattn_2025} es un \textbf{método universal sin entrenamiento} que predice qué bloques de la matriz de atención contendrán valores insignificantes y omite el cálculo para esos bloques.

\paragraph{Etapa 1 --- Predicción dispersa}: Calcular importancia proxy de bajo costo $\hat{s}_{ij}$ para el bloque $(i,j)$:

\[
\hat{s}_{ij} = f_{\text{pred}}(\mathbf{Q}_i, \mathbf{K}_j; \text{hiperparámetros}), \quad \text{omitir si } \hat{s}_{ij} < \epsilon_1
\]

Los predictores comunes incluyen similitud de media de bloque o estadísticas de autosimilitud.

\paragraph{Etapa 2 --- Filtrado consciente de softmax}: Después de calcular $\tilde{P}_{ij} = \text{softmax}(Q_i K_j^\top/\sqrt{d})$, verificar si la probabilidad máxima del bloque es insignificante en relación con el máximo de softmax en ejecución:

\[
\text{omitir } P_{ij} V_j \quad \text{si} \quad \max(\tilde{P}_{ij}) < e^{m_{\text{antiguo}} - m_{\text{nuevo}}} \cdot \epsilon_2
\]

donde $m_{\text{antiguo}}, m_{\text{nuevo}}$ son máximos del softmax en línea en ejecución (calculados mediante numéricos estilo FlashAttention).

\subsubsection{Pseudocódigo Algorítmico}

\begin{algorithm}[H]
\caption{SpargeAttn: Filtrado Disperso de Dos Etapas por Bloques}
\begin{algorithmic}[1]
\Require Bloques de consulta $\mathbf{Q}_i$ para $i = 1, \ldots, n_b$, Bloques de clave $\mathbf{K}_j$ para $j = 1, \ldots, n_b$, Bloques de valor $\mathbf{V}_j$
\Require Umbral de predicción $\epsilon_1$, umbral de softmax $\epsilon_2$, tamaño de bloque $b_s$
\Ensure Salida $\mathbf{Y}$
\State Inicializar: máximo de softmax en línea $m_{\text{global}} \gets -\infty$
\For{fila de bloque $i = 1$ to $n_b$}
    \State Inicializar: salida de fila de bloque $\mathbf{Y}_i \gets 0$, máximo local $m_i \gets -\infty$
    \For{columna de bloque $j = 1$ to $n_b$}
        \State \textbf{Etapa 1: Predicción Dispersa}
        \State Calcular importancia proxy: $\hat{s}_{ij} \gets f_{\text{pred}}(\mathbf{Q}_i, \mathbf{K}_j)$
        \If{$\hat{s}_{ij} < \epsilon_1$}
            \State \textbf{omitir} este bloque $[i,j]$  \Comment{Terminación temprana}
            \State \textbf{continuar} \Comment{Pasar al siguiente bloque}
        \EndIf
        \State \textbf{Etapa 2: Calcular y Filtrar}
        \State $\tilde{P}_{ij} \gets \text{softmax}(\mathbf{Q}_i \mathbf{K}_j^\top / \sqrt{d})$
        \State $m_{\text{bloque}} \gets \max(\tilde{P}_{ij})$
        \If{$m_{\text{bloque}} < e^{m_i - m_{\text{global}}} \cdot \epsilon_2$}
            \State El bloque contribuye insignificantemente; omitir  \Comment{Poda consciente de softmax}
            \State \textbf{continuar}
        \EndIf
        \State Actualizar máximo global: $m_i \gets \max(m_i, m_{\text{bloque}})$
        \State Acumular: $\mathbf{Y}_i \gets \mathbf{Y}_i + \tilde{P}_{ij} \mathbf{V}_j$
    \EndFor
    \State $m_{\text{global}} \gets \max(m_{\text{global}}, m_i)$
\EndFor
\Return $\mathbf{Y}$
\end{algorithmic}
\end{algorithm}

\subsubsection{Características Clave}

\begin{itemize}[leftmargin=*]
    \item \textbf{Complejidad}: Empírica $\mathcal{O}(n^2 \cdot s)$ donde $s$ es la fracción de bloques no omitidos (típicamente 0.2--0.5).
    \item \textbf{Entrenable}: No; aceleración plug-and-play para modelos existentes.
    \item \textbf{Universalidad}: Funciona en modelos de lenguaje, modelos de difusión de imágenes y generación de video.
    \item \textbf{Fortaleza}: 2.5--5$\times$ de aceleración comparado con métodos densos o dispersos anteriores; compatible con cuantización (integración SageAttention).
    \item \textbf{Limitación}: La aceleración depende de la dispersidad inherente; la granularidad a nivel de bloque puede pasar por alto patrones más finos; se requiere ajuste de umbral por familia de modelos.
\end{itemize}

\subsection{Resumen Comparativo}

\footnotesize
\begin{longtable}{p{1.9cm}p{1.7cm}p{1.6cm}p{1.6cm}p{1.1cm}p{3.5cm}}
\toprule
\textbf{Método} & \textbf{Complejidad} & \textbf{Entrenable} & \textbf{Unidad} & \textbf{Año} & \textbf{Contribución Principal} \\
\midrule
\endhead
Sparse Transformer & $O(n\sqrt{n}d)$ & Sí & Patrón de token & 2019 & Máscaras fijas y estrideadas factorizadas \\
Longformer & $O(nwd)$ & Sí & Local + global & 2020 & Escalado lineal con dilatación \\
BigBird & $O(n)$ & Sí & Grafo disperso & 2020 & Prueba de completitud teórica \\
SparseK & $O(nd)$ & Sí & Tokens top-$k$ & 2024 & Selección diferenciable \\
NSA & $O((t/d+nl'+w)d)$ & Sí & Multi-rama & 2025 & Diseño de 3 ramas alineado con hardware \\
FASA & $O(tN_{\text{tip}}+N_{\text{fac}}d)$ & No & Fragmentos de frecuencia & 2026 & Conocimiento de frecuencia RoPE \\
SpargeAttn & $O(n^2 s \cdot d)$ & No & Filtro de bloques & 2025 & Filtrado universal de dos etapas \\
\bottomrule
\end{longtable}
\normalsize

\subsection{Espacio de Diseño y Criterios de Selección}

Los siete métodos de atención dispersa ocupan posiciones complementarias en un espacio de diseño multidimensional:

\begin{itemize}[leftmargin=*]
    \item \textbf{Patrones geométricos/fijos} (Sparse Transformer, Longformer): Simples de implementar y analizar, pero los patrones no se adaptan al contenido.
    \item \textbf{Selección aprendida} (SparseK, NSA): Permiten adaptación mediante redes de selección entrenables, a costa de mayor complejidad de implementación.
    \item \textbf{Aceleración sin entrenamiento} (FASA, SpargeAttn): Permiten aceleración directa de modelos existentes sin reentrenamiento, adecuados para sistemas desplegados.
    \item \textbf{Garantías teóricas} (BigBird): Proporcionan demostraciones formales de completitud expresiva, importantes para comprender los márgenes de seguridad.
\end{itemize}

\textbf{Orientación de selección:}
\begin{enumerate}[leftmargin=*]
    \item \textbf{Entrenamiento de nuevos modelos} donde los recursos lo permitan: NSA o SparseK para máxima aceleración de inferencia; BigBird para garantía teórica.
    \item \textbf{Comprensión de documentos de contexto largo}: Longformer o FASA por simplicidad práctica; NSA para escala extrema.
    \item \textbf{Aceleración de modelos existentes}: FASA para LLMs basados en RoPE; SpargeAttn para cualquier arquitectura.
    \item \textbf{Entornos con recursos limitados}: Sparse Transformer por simplicidad y eficiencia probada a escala moderada.
\end{enumerate}


\section{Annex D: Gated Attention Families---Complete Literature Analysis}
\label{sec:annex-gated}

\subsection{Resumen Ejecutivo: Compuertas para el Control de Memoria}

El campo emergente de los \emph{mecanismos de atención con compuerta} aborda un desafío fundamental en el modelado de secuencias: ¿cómo debe retenerse, olvidarse y actualizarse la información a medida que el modelo procesa nuevos tokens? Los estados recurrentes aditivos tradicionales acumulan toda la información sin borrado, lo que lleva a la saturación de memoria y a la incapacidad de adaptarse a contextos cambiantes. La atención softmax proporciona expresividad completa pero usa un caché KV de $\mathcal{O}(Ld)$ durante la inferencia, limitando el tamaño de la ventana de contexto. Los mecanismos con compuerta proporcionan un punto intermedio: control selectivo sobre la supervivencia de la información.

El principio unificador entre las arquitecturas con compuerta es que las compuertas controlan \textit{qué información sobrevive}. Las compuertas operan en tres niveles distintos:
\begin{enumerate}[leftmargin=*]
    \item \textbf{Decaimiento del estado recurrente} (GLA, HGRN2): Compuertas multiplicativas sobre la matriz de memoria $S_t$ controlan cuánto del estado anterior persiste en el siguiente paso.
    \item \textbf{Intensidad de escritura en actualizaciones recurrentes} (DeltaNet, Gated DeltaNet, Gated DeltaNet-2): Las compuertas controlan la magnitud y la dirección canalizada de la nueva información escrita en la memoria.
    \item \textbf{Sesgo de logits softmax} (FoX): Las compuertas inyectan sesgo de actualidad a nivel de token en el cálculo de logits de atención.
    \item \textbf{Dispersidad post-atención} (Gated Softmax): Las compuertas escalan selectivamente los canales de salida de SDPA.
\end{enumerate}

Este apéndice sintetiza ocho arquitecturas principales de atención con compuerta publicadas entre 2023 y 2026, analizando sus fundamentos matemáticos, características de eficiencia en hardware y compensaciones empíricas.

\subsection{1. Atención Lineal con Compuerta (GLA)}

\subsubsection{Fundamento Matemático}

La Atención Lineal con Compuerta \citep{yang_gla_2023} aumenta la atención lineal estándar (que usa un estado recurrente con valor de matriz) con decaimiento multiplicativo dependiente de los datos. La atención lineal estándar reformula el mecanismo softmax tradicional como una recurrencia sobre estados ocultos:
\[
\mathbf{S}_t = \mathbf{S}_{t-1} + \mathbf{v}_t \mathbf{k}_t^\top, \quad \mathbf{o}_t = \mathbf{S}_t \mathbf{q}_t
\]
donde el estado $\mathbf{S}_t \in \mathbb{R}^{d_v \times d_k}$ acumula productos externos. Esta formulación puramente aditiva sufre de ``sobrecarga de memoria'': la información se acumula monótonamente y no puede borrarse, degradando el rendimiento en tareas que requieren selección de contexto.

GLA introduce una \textbf{matriz de compuerta diagonal} $\mathbf{G}_t \in [0,1]^{d_k \times d_k}$ dependiente de los datos:
\[
\mathbf{S}_t = \mathbf{G}_t \odot \mathbf{S}_{t-1} + \mathbf{v}_t \mathbf{k}_t^\top, \quad \mathbf{o}_t = \mathbf{S}_t \mathbf{q}_t
\]
La compuerta $\mathbf{G}_t$ se parametriza con una estructura de producto externo:
\[
\mathbf{G}_t = \mathbf{1} \boldsymbol{\alpha}_t^\top, \quad \boldsymbol{\alpha}_t = \text{logsigmoid}\left(\mathbf{W}_{gk} \mathbf{x}_t\right) / c
\]
donde $\mathbf{W}_{gk}$ es una proyección de rango bajo ($\mathbb{R}^{d \to d_k}$) y $c \approx 16$ es un normalizador. La activación logsigmoid mapea $\alpha_t \in \mathbb{R}$ a aproximadamente $[-1/2, 0]$, y la división por $c$ contrae esto aún más para que esté cerca de la identidad en la inicialización. La compuerta resultante $G_t = 1 + \alpha_t \approx 1$ cerca de la inicialización, luego aprende a decaer $(\alpha_t \to -\infty)$ información histórica.

\begin{algorithm}[H]
\caption{Atención Lineal con Compuerta (Forma Recurrente)}
\begin{algorithmic}[1]
\Require Entrada $\mathbf{x}_t$; estado anterior $\mathbf{S}_{t-1}$
\Require Proyecciones: $W_q, W_k, W_v$ (estándar); $W_{gk}$ (proyección de compuerta)
\Ensure Salida $\mathbf{o}_t$ y estado actualizado $\mathbf{S}_t$
\State $\mathbf{q}_t \gets W_q \mathbf{x}_t$
\State $\mathbf{k}_t \gets W_k \mathbf{x}_t$
\State $\mathbf{v}_t \gets W_v \mathbf{x}_t$
\State Calcular compuerta: $\boldsymbol{\alpha}_t \gets \text{logsigmoid}(W_{gk} \mathbf{x}_t) / 16$ \Comment{Decaimiento por dimensión de clave}
\State Aplicar compuerta de producto externo: $\mathbf{G}_t \gets \mathbf{1} \boldsymbol{\alpha}_t^\top$  \Comment{producto externo diagonal}
\State Decaer estado anterior: $\mathbf{S}_t^\text{decaimiento} \gets \mathbf{G}_t \odot \mathbf{S}_{t-1}$ \Comment{producto elemento a elemento}
\State Agregar nueva asociación: $\mathbf{S}_t \gets \mathbf{S}_t^\text{decaimiento} + \mathbf{v}_t \mathbf{k}_t^\top$
\State Recuperar mediante consulta: $\mathbf{o}_t^\text{crudo} \gets \mathbf{S}_t \mathbf{q}_t$
\State Aplicar compuerta de salida: $\mathbf{o}_t \gets \text{GroupNorm}(\mathbf{o}_t^\text{crudo}) \otimes \text{SiLU}(W_g \mathbf{x}_t)$
\Return $\mathbf{o}_t$, $\mathbf{S}_t$
\end{algorithmic}
\end{algorithm}

\subsubsection{Entrenamiento Eficiente en Hardware}

Para el entrenamiento paralelo, GLA usa un algoritmo por fragmentos que agrupa tokens en fragmentos $C$ y calcula transiciones recurrentes en paralelo:
\[
\mathbf{O}_{[t]} = \overleftarrow{\mathbf{Q}}_{[t]} \mathbf{S}_{[t]}^\top + \left(\mathbf{Q}_{[t]} \mathbf{K}_{[t]}^\top \odot \mathbf{\Gamma}_{[t]}\right) \mathbf{V}_{[t]}
\]
donde $\overleftarrow{\mathbf{Q}}_{[t]}$ atiende al estado en el límite del fragmento (retrospectiva), y $\mathbf{\Gamma}_{[t]}$ es la máscara causal con escalado consciente del decaimiento:
\[
\mathbf{\Gamma}_{[t],ij} = \begin{cases} \prod_{l=j}^{i-1} \alpha_l & \text{si } i > j \text{ (retrospectiva)} \\ \prod_{l=1}^{i-j} \alpha_l & \text{si } i \le j \text{ (en-fragmento)} \end{cases}
\]
Este diseño maximiza las operaciones matmul susceptibles de aceleración mediante núcleos tensoriales, logrando una complejidad total subcuadrática $\mathcal{O}(nLd^2/C)$ donde $L$ es el número de cabezas de atención y $C$ es el tamaño del fragmento.

\subsubsection{Fortalezas y Limitaciones}

\begin{table}[H]
\centering
\small
\begin{tabularx}{0.95\linewidth}{Xp{4cm}p{4cm}}
\toprule
\textbf{Aspecto} & \textbf{Fortalezas} & \textbf{Limitaciones} \\
\midrule
Eficiencia computacional & Entrenamiento subcuadrático mediante paralelismo por fragmentos. Inferencia con memoria constante $O(d^2)$. & Requiere kernels CUDA/Triton personalizados; no disponible en todos los frameworks. \\
Generalización de contexto & Extiende entrenamiento de 2K tokens a más de 20K tokens con degradación de perplejidad insignificante. & Aún rinde por debajo de softmax en benchmarks con mucha recuperación (aguja en el pajar). \\
Selectividad & Decaimiento dependiente de los datos por dimensión de clave. & La estructura de compuerta limita la selectividad por par clave-valor; no hay control directo sobre qué asociaciones específicas borrar. \\
\bottomrule
\end{tabularx}
\end{table}

\subsection{2. DeltaNet: Atención Lineal con Corrección de Errores}

\subsubsection{Fundamento Matemático}

DeltaNet \citep{yang_deltanet_2024} aplica una \textbf{regla de aprendizaje delta} clásica a las actualizaciones de estado de la atención lineal. En lugar de acumulación aditiva, DeltaNet calcula la diferencia entre el valor predicho (recuperado de la memoria) y el valor objetivo, luego realiza una actualización de corrección de errores:
\[
\mathbf{S}_t = \mathbf{S}_{t-1}(\mathbf{I} - \beta_t \mathbf{k}_t \mathbf{k}_t^\top) + \beta_t \mathbf{v}_t \mathbf{k}_t^\top
\]
donde $\beta_t \in (0,1)$ es un escalar de \textbf{intensidad de escritura} dependiente de los datos. Esta actualización puede descomponerse como una operación de borrado-escritura:
\[
\mathbf{v}_t^\text{anterior} = \mathbf{S}_{t-1} \mathbf{k}_t \quad \text{(predicción actual)}, \quad
\mathbf{v}_t^\text{actualizado} = \beta_t \mathbf{v}_t + (1-\beta_t)\mathbf{v}_t^\text{anterior} \quad \text{(mezcla anterior/nuevo)}
\]
de modo que:
\[
\mathbf{S}_t = \mathbf{S}_{t-1} - \mathbf{v}_t^\text{anterior}\mathbf{k}_t^\top + \mathbf{v}_t^\text{actualizado}\mathbf{k}_t^\top
\]

La idea clave es que esta regla de actualización minimiza una pérdida de ECM en línea en cada paso temporal:
\[
\mathcal{L}_t(\mathbf{S}) = \tfrac{1}{2}\|\mathbf{S}\mathbf{k}_t - \mathbf{v}_t\|^2 \Rightarrow \nabla_{\mathbf{S}} \mathcal{L}_t = (\mathbf{S}\mathbf{k}_t - \mathbf{v}_t)\mathbf{k}_t^\top
\]

Un paso de SGD con tasa de aprendizaje $\beta_t$ produce la actualización de la regla delta. Esta conexión con el entrenamiento en tiempo de prueba (TTT) proporciona fundamentación teórica: DeltaNet optimiza directamente la predicción de valores en tiempo de inferencia.

\subsubsection{Entrenamiento Paralelo Eficiente}

La matriz de transición de DeltaNet $\mathbf{I} - \beta_t \mathbf{k}_t \mathbf{k}_t^\top$ es una matriz de Householder generalizada (actualización de rango 1 a la identidad). Para el entrenamiento por fragmentos, DeltaNet usa la \textbf{representación WY}, que representa el producto de $m$ de estas actualizaciones de rango 1 de forma compacta:
\[
\prod_{i=1}^{m} (\mathbf{I} - \beta_i \mathbf{k}_i \mathbf{k}_i^\top) = \mathbf{I} - \mathbf{W}\mathbf{Y}^\top
\]
donde $\mathbf{W}, \mathbf{Y} \in \mathbb{R}^{d \times m}$ se construyen recursivamente. Esta factorización permite paralelismo rico en matmul sin materializar el producto completo, manteniendo el entrenamiento por fragmentos eficiente.

\begin{algorithm}[H]
\caption{Actualización de Estado de DeltaNet con Representación WY}
\begin{algorithmic}[1]
\Require Fragmento de entrada $\mathbf{X}_{[t]}$; estado anterior $\mathbf{S}_{t-1}$
\Require Consultas normalizadas L2 $\hat{\mathbf{Q}}_t$, claves $\hat{\mathbf{K}}_t$, valores $\mathbf{V}_t$
\Ensure Salida $\mathbf{O}_{[t]}$ y estado actualizado $\mathbf{S}_t$
\State \textbf{Fase por fragmentos:}
\For{posición $i = 1$ to tamaño\_fragmento}
    \State Normalizar: $\hat{\mathbf{k}}_i \gets \hat{\mathbf{K}}_t[i] / \|\hat{\mathbf{K}}_t[i]\|$, $\hat{\mathbf{q}}_i \gets \hat{\mathbf{Q}}_t[i] / \|\hat{\mathbf{Q}}_t[i]\|$
    \State Calcular intensidad de escritura: $\beta_i \gets \text{sigmoid}(W_\beta \mathbf{x}_i)$
    \State Predicción: $\hat{\mathbf{v}}^\text{pred} \gets \mathbf{S}_{i-1}^\text{en-fragmento} \hat{\mathbf{k}}_i$
    \State Mezcla consciente del error: $\mathbf{v}_i^\text{actualización} \gets \beta_i(\mathbf{v}_i - \hat{\mathbf{v}}^\text{pred})$
    \State Borrar-escribir: $\mathbf{S}_i^\text{en-fragmento} \gets \mathbf{S}_{i-1}^\text{en-fragmento} - \hat{\mathbf{v}}^\text{pred}\hat{\mathbf{k}}_i^\top + (\beta_i \mathbf{v}_i + (1-\beta_i)\hat{\mathbf{v}}^\text{pred})\hat{\mathbf{k}}_i^\top$
    \State Salida: $\mathbf{o}_i \gets \mathbf{S}_i^\text{en-fragmento} \hat{\mathbf{q}}_i$
\EndFor
\State \textbf{Fase entre fragmentos:} Usar representación WY de los productos de transición
\Return $\mathbf{O}_{[t]}$, $\mathbf{S}_{t}$
\end{algorithmic}
\end{algorithm}

\subsubsection{Fortalezas y Limitaciones}

\begin{table}[H]
\centering
\small
\begin{tabularx}{0.95\linewidth}{Xp{4cm}p{4cm}}
\toprule
\textbf{Aspecto} & \textbf{Fortalezas} & \textbf{Limitaciones} \\
\midrule
Recuperación asociativa & Recuperación perfecta en el benchmark MQAR (Recuperación Asociativa Multi-Consulta); la semántica de corrección de errores se ajusta naturalmente a patrones de recuperación. & Sin olvido global, la memoria aún se satura en longitudes de secuencia extremas; requiere mecanismos de reinicio periódico para contexto ilimitado. \\
Teoría & Fundamentada en optimización de ECM en línea; conexión clara con el paradigma de entrenamiento en tiempo de prueba. & Requiere claves normalizadas L2 para estabilidad numérica; la normalización de doble ancho añade sobrecarga. \\
Rendimiento & La representación WY permite entrenamiento eficiente por fragmentos. & Ligeramente más lento que Mamba2 por token debido a matrices de transición más ricas (rango 1 en lugar de diagonales). \\
\bottomrule
\end{tabularx}
\end{table}

\subsection{3. Gated DeltaNet: Síntesis de Compuerta y Corrección de Errores}

\subsubsection{Fundamento Matemático}

Gated DeltaNet \citep{yang_gated_deltanet_2024} sintetiza las fortalezas de GLA y DeltaNet combinando una \textbf{compuerta de decaimiento} $\alpha_t$ (de GLA) con la \textbf{compuerta de intensidad de escritura} $\beta_t$ (de DeltaNet):
\[
\mathbf{S}_t = \mathbf{S}_{t-1}\left(\alpha_t (\mathbf{I} - \beta_t \mathbf{k}_t \mathbf{k}_t^\top)\right) + \beta_t \mathbf{v}_t \mathbf{k}_t^\top
\]
Reescribiendo para mayor claridad:
\[
\mathbf{S}_t = \alpha_t \mathbf{S}_{t-1}(\mathbf{I} - \beta_t \mathbf{k}_t \mathbf{k}_t^\top) + \beta_t \mathbf{v}_t \mathbf{k}_t^\top
\]
Las dos compuertas son complementarias:
\begin{itemize}[leftmargin=*]
    \item $\alpha_t \to 0$ borra rápidamente el estado histórico: $\mathbf{S}_t \approx \beta_t \mathbf{v}_t \mathbf{k}_t^\top$ (reinicio de contexto).
    \item $\alpha_t \to 1$ recupera el comportamiento puro de regla delta: $\mathbf{S}_t \approx \mathbf{S}_{t-1}(\mathbf{I} - \beta_t \mathbf{k}_t \mathbf{k}_t^\top) + \beta_t \mathbf{v}_t \mathbf{k}_t^\top$ (actualizaciones dirigidas).
    \item $\beta_t \to 0$ omite la escritura pero decae: $\mathbf{S}_t \approx \alpha_t \mathbf{S}_{t-1}$ (olvido puro).
    \item $\beta_t \to 1$ reemplaza la memoria agresivamente: $\mathbf{S}_t \approx \alpha_t \mathbf{S}_{t-1} + \mathbf{v}_t \mathbf{k}_t^\top$ (reemplazo).
\end{itemize}

Desde una perspectiva de aprendizaje en línea, Gated DeltaNet corresponde a:
\[
\min_{\mathbf{S}_t} \|\mathbf{S}_t - \alpha_t \mathbf{S}_{t-1}\|_F^2 - 2\langle \mathbf{S}_t \mathbf{k}_t, \beta_t(\mathbf{v}_t - \alpha_t \mathbf{S}_{t-1} \mathbf{k}_t)\rangle
\]
que introduce un decaimiento de pesos adaptativo $\alpha_t$ en una actualización tipo SGD---análogo al decaimiento de pesos desacoplado en la optimización de aprendizaje profundo.

\begin{algorithm}[H]
\caption{Paso Adelante Paralelo por Fragmentos de Gated DeltaNet}
\begin{algorithmic}[1]
\Require Fragmento de entrada $\mathbf{X}_{[i]}$; estado anterior $\mathbf{S}_{i-1}$; tamaño de fragmento $C$
\Require Proyecciones: $W_q, W_k, W_v, W_\alpha, W_\beta$
\Ensure Salida $\mathbf{O}_{[i]}$ y estado $\mathbf{S}_i$
\State \textbf{Cálculo en-fragmento:}
\For{$j = 1$ to $C$}
    \State $\mathbf{q}_j \gets W_q \mathbf{x}_j$, $\mathbf{k}_j \gets W_k \mathbf{x}_j$, $\mathbf{v}_j \gets W_v \mathbf{x}_j$
    \State $\alpha_j \gets \text{sigmoid}(W_\alpha \mathbf{x}_j)$, $\beta_j \gets \text{sigmoid}(W_\beta \mathbf{x}_j)$
    \State Aplicar compuerta combinada: $\mathbf{S}_j \gets \alpha_j \mathbf{S}_{j-1}(\mathbf{I} - \beta_j \mathbf{k}_j \mathbf{k}_j^\top) + \beta_j \mathbf{v}_j \mathbf{k}_j^\top$
    \State $\mathbf{o}_j \gets \mathbf{S}_j \mathbf{q}_j$
\EndFor
\State \textbf{Recurrencia entre fragmentos:}
\State $\mathbf{S}_i \gets \mathbf{S}_C$ del en-fragmento; propagar entre fragmentos mediante máscaras $\prod \alpha$ acumulativas
\State \textbf{Compuerta de salida:} $\mathbf{O}_{[i]} \gets \text{GroupNorm}(\mathbf{O}^\text{crudo}) \otimes \text{SiLU}(W_g \mathbf{X}_{[i]})$
\Return $\mathbf{O}_{[i]}$, $\mathbf{S}_i$
\end{algorithmic}
\end{algorithm}

\subsubsection{Rendimiento Empírico y Compensaciones}

Gated DeltaNet se ha integrado en los modelos de producción Qwen3-Next y Qwen3.5 de Alibaba. Empíricamente:
\[
\begin{array}{l|cccc}
\textbf{Tarea} & \textbf{Gated DeltaNet} & \textbf{Mamba2} & \textbf{DeltaNet} & \textbf{GLA} \\
\hline
\text{Modelado de Lenguaje} & \mathbf{1.0\times} & 1.08\times & 1.05\times & 1.12\times \\
\text{Recuperación Asociativa} & \mathbf{1.0\times} & --(sin perfecto) & \mathbf{1.0\times} & 0.75 \\
\text{Extrapolación de Longitud} & \mathbf{1.0\times} & 0.98\times & 0.96\times & 0.94\times \\
\text{Rendimiento (tokens/seg)} & 0.85\times & \mathbf{1.0\times} & 0.82\times & 0.88\times
\end{array}
\]
Gated DeltaNet logra el mejor equilibrio entre diversas tareas a costa de un rendimiento ligeramente reducido debido a matrices de transición más ricas.

\subsection{4. HGRN2: Compuerta Jerárquica con Expansión de Producto Externo}

\subsubsection{Fundamento Matemático}

HGRN2 \citep{qin_hgrn2_2024} usa una expansión de estado basada en producto externo con \textbf{compuertas de olvido acotadas inferiormente de forma jerárquica}. La actualización de estado es:
\[
\mathbf{S}_t = \text{diag}(\mathbf{g}_t) \cdot \mathbf{S}_{t-1} + \mathbf{v}_t \mathbf{k}_t^\top, \quad \mathbf{o}_t = \mathbf{S}_t \mathbf{q}_t
\]
donde $\mathbf{g}_t \in [b_\ell, 1]^d$ es una compuerta de olvido con cota inferior para la capa $\ell$. Las cotas satisfacen:
\[
0 \le b_{\ell_{\text{poco profunda}}} < b_{\ell_{\text{media}}} < b_{\ell_{\text{profunda}}} \le 1
\]
es decir, las cotas aumentan (se vuelven menos restrictivas) en las capas más profundas. La compuerta en sí se calcula como:
\[
\mathbf{g}_t = b_\ell + (1 - b_\ell) \cdot \sigma(W_g \mathbf{x}_t)
\]
donde $\sigma$ es sigmoide. Cuando $W_g$ produce logits débilmente negativos, $\mathbf{g}_t \approx b_\ell$ (retención forzada en capas superficiales). Cuando las salidas son fuertemente positivas, $\mathbf{g}_t \approx 1$ (refresco de memoria en cualquier capa).

La estructura jerárquica anima a diferentes capas a especializarse en diferentes escalas temporales: las capas superficiales modelan dependencias locales (alta retención debido a la cota $b_\ell$), mientras que las capas profundas capturan estructura de largo alcance (compuerta flexible con cota superior alta).

\begin{algorithm}[H]
\caption{Paso Adelante de HGRN2 con Cotas Jerárquicas}
\begin{algorithmic}[1]
\Require Entrada $\mathbf{x}_t$; estado anterior $\mathbf{S}_{t-1}$; índice de capa $\ell \in [0, L-1]$
\Require Cotas no decrecientes: $0 \le b_0 < b_1 < \cdots < b_{L-1} \le 1$
\Ensure Salida $\mathbf{o}_t$ y estado $\mathbf{S}_t$
\State Cota de retención específica de capa: $b \gets b_\ell$
\State Proyectar: $\mathbf{q}_t \gets W_q \mathbf{x}_t$, $\mathbf{k}_t \gets W_k \mathbf{x}_t$, $\mathbf{v}_t \gets W_v \mathbf{x}_t$
\State Calcular compuerta de olvido con cota: $\mathbf{g}_t \gets b + (1-b) \cdot \sigma(W_g \mathbf{x}_t)$ \Comment{Confinado a $[b, 1]$}
\State Multiplicación diagonal (vectorizada): $\mathbf{S}_t \gets \mathbf{S}_{t-1} \circ \text{diag}(\mathbf{g}_t) + \mathbf{v}_t \mathbf{k}_t^\top$
\State Recuperar: $\mathbf{o}_t \gets \mathbf{S}_t \mathbf{q}_t$
\State Normalización opcional: $\mathbf{o}_t \gets \text{RMSNorm}(\mathbf{o}_t)$
\Return $\mathbf{o}_t$, $\mathbf{S}_t$
\end{algorithmic}
\end{algorithm}

\subsubsection{Escalado y Resultados Empíricos}

A escala de 3B en 100B tokens, HGRN2 supera ligeramente a Mamba2 y a los Transformers con arquitectura LLaMA en modelado de lenguaje. La compuerta jerárquica proporciona modelado temporal multi-escala sin variantes arquitectónicas explícitas por capa. Sin embargo, la compuerta diagonal con valor vectorial (en oposición a la selección elemento a elemento) proporciona menos flexibilidad por elemento que las variantes de regla delta.

\subsection{5. Forgetting Transformer (FoX): Compuerta en el Espacio de Logits Softmax}

\subsubsection{Fundamento Matemático}

Forgetting Transformer (FoX), propuesto por Lin et al.~\citep{lin_forgetting_transformer_2025}, incrusta una compuerta de olvido directamente en los logits de atención softmax. En lugar de reemplazar softmax con recurrencia lineal, FoX preserva la expresividad completa de softmax mientras añade control de actualidad mediante un sesgo de logit dependiente de los datos.

Para cada posición de token $t$ y todas las claves $j \le t$, calcular un escalar de compuerta de olvido:
\[
f_t = \sigma(w_f^\top \mathbf{x}_t + b_f)
\]
El sesgo de logit en la posición $(i,j)$ es el producto acumulativo log-olvido:
\[
d_{ij} = \sum_{l=j+1}^{i} \log f_l = \log \prod_{l=j+1}^{i} f_l
\]
La salida de atención completa se convierte en:
\[
\mathbf{O} = \text{softmax}\left(\frac{\mathbf{Q}\mathbf{K}^\top}{\sqrt{d_k}} + \mathbf{D}\right) \mathbf{V}
\]
donde $\mathbf{D}_{ij} = d_{ij}$. Equivalentemente:
\[
\text{Atencion}(i,j) = \frac{\exp(\mathbf{q}_i^\top \mathbf{k}_j / \sqrt{d_k} + d_{ij})}{\sum_{j'=1}^{i} \exp(\mathbf{q}_i^\top \mathbf{k}_{j'} / \sqrt{d_k} + d_{ij'})}
\]

Esta ponderación reduce el peso de los tokens más antiguos multiplicativamente: un token de hace 10 pasos se escala por $\prod_{l=j+1}^{i} f_l$, que es típicamente mucho menor que 1 si $f_t < 1$ en promedio.

El mecanismo es matemáticamente equivalente a una \textbf{variante aprendible dependiente de los datos de ALiBi (Atención con Sesgos Lineales)}, donde trabajos anteriores usaban $d_{ij} = -\infty \cdot |i-j|$ fijo o $d_{ij} = -(i-j)$.

\begin{algorithm}[H]
\caption{FoX: Atención con Olvido y Sesgo de Actualidad}
\begin{algorithmic}[1]
\Require Consultas $\mathbf{Q}$, claves $\mathbf{K}$, valores $\mathbf{V}$; entrada $\mathbf{X}$
\Require Parámetros de compuerta de olvido: $w_f \in \mathbb{R}^d$, $b_f \in \mathbb{R}$
\Ensure Salida de atención $\mathbf{O}$
\State \textbf{Calcular compuertas de olvido:}
\For{$t = 1$ to $T$}
    \State $f_t \gets \sigma(w_f^\top \mathbf{x}_t + b_f)$ \Comment{Probabilidad de olvido por token}
\EndFor
\State \textbf{Calcular sesgos acumulativos de log-olvido:}
\For{$i = 1$ to $T$}
    \For{$j = 1$ to $i$}
        \State $d_{ij} \gets \sum_{l=j+1}^{i} \log f_l$ \Comment{Producto acumulativo en espacio logarítmico}
    \EndFor
\EndFor
\State \textbf{SDPA estándar con sesgo:}
\State Calcular puntajes: $\mathbf{S} \gets \mathbf{Q}\mathbf{K}^\top / \sqrt{d_k}$
\State Agregar sesgo: $\mathbf{S} \gets \mathbf{S} + \mathbf{D}$
\State Aplicar softmax: $\mathbf{A} \gets \text{softmax}(\mathbf{S}, \text{dim}=-1)$
\State Ponderar valores: $\mathbf{O} \gets \mathbf{A} \mathbf{V}$
\Return $\mathbf{O}$
\end{algorithmic}
\end{algorithm}

\subsubsection{Integración con FlashAttention}

La ventaja clave de FoX es la compatibilidad con FlashAttention y las implementaciones optimizadas existentes. Los sesgos de olvido se añaden en el cálculo de logits softmax, que ya es una operación central en el algoritmo por bloques de FlashAttention. La sobrecarga es:
\[
\text{Sobrecarga} \approx 0.5--2\% \text{ del tiempo de pared}
\]
ya que la adición de sesgo se amortiza entre las operaciones matmul.

\subsubsection{Fortalezas y Limitaciones}

FoX logra resultados de última generación en extrapolación de longitud, recuperación casi perfecta de aguja en el pajar y comprensión superior de contexto largo en comparación con los Transformers. Sin embargo, sigue siendo $\mathcal{O}(L^2 d)$ en entrenamiento y $\mathcal{O}(L d)$ por paso en inferencia con caché KV, limitando las longitudes de secuencia extremas.

\subsection{6. Atención con Compuerta (Post-SDPA Sigmoide)}

\subsubsection{Fundamento Matemático}

El Mejor Artículo de NeurIPS 2025 del equipo de Qwen propone aplicar una compuerta sigmoide \textbf{después} de la atención de producto punto escalado (SDPA):
\[
\mathbf{Y} = \text{SDPA}(\mathbf{Q}, \mathbf{K}, \mathbf{V}) = \text{softmax}\left(\frac{\mathbf{Q}\mathbf{K}^\top}{\sqrt{d_k}}\right) \mathbf{V}
\]
\[
\mathbf{Y}' = \mathbf{Y} \odot \sigma\left(\mathbf{X} \mathbf{W}_\theta\right)
\]
donde el puntaje de compuerta $\sigma(\mathbf{X} \mathbf{W}_\theta)$ puede tener forma $\mathbb{R}^{n \times q \times d_k}$ (compuerta elemento a elemento por cabeza de consulta) o $\mathbb{R}^{n \times q}$ (compuerta fija por cabeza). En más de 30 variantes y modelos MoE de 15B + densos de 1.7B entrenados en 3.5T tokens, el equipo encontró:

\begin{itemize}[leftmargin=*]
    \item La compuerta por cabeza añade solo $\sim 1.6M$ parámetros a un modelo de 15B (sobrecarga insignificante).
    \item Las compuertas efectivas son \textbf{dispersas}: activación media $\approx 0.116$ (es decir, el 88\% son cero o casi cero).
    \item Las compuertas actúan como filtros dependientes de la consulta, suprimiendo canales de bajo valor mientras preservan señales de alto valor.
    \item Las compuertas eliminan demostrablemente el fenómeno de \textbf{sumidero de atención}, donde el patrón de atención se vuelve dominado por un único token temprano.
\end{itemize}

\begin{algorithm}[H]
\caption{Atención Softmax con Compuerta (Post-SDPA)}
\begin{algorithmic}[1]
\Require Consultas $\mathbf{Q}$, Claves $\mathbf{K}$, Valores $\mathbf{V}$; entrada $\mathbf{X}$
\Require Matriz de proyección de compuerta $W_g \in \mathbb{R}^{d \to d_{\text{compuerta}}}$ donde $d_{\text{compuerta}} \in \{1, d_k\}$
\Ensure Salida con compuerta $\mathbf{Y}'$
\State \textbf{SDPA estándar:}
\State Calcular atención: $\mathbf{A} \gets \text{softmax}(\mathbf{Q}\mathbf{K}^\top / \sqrt{d_k})$ 
\State Ponderar valores: $\mathbf{Y} \gets \mathbf{A} \mathbf{V}$
\State \textbf{Calcular compuertas:}
\State Proyectar entrada: $\mathbf{G} \gets \mathbf{X} W_g$  \Comment{forma: $[n, q, d_{\text{compuerta}}]$}
\State Aplicar sigmoide: $\mathbf{G} \gets \sigma(\mathbf{G})$ \Comment{Compuerta elemento a elemento}
\State \textbf{Aplicar compuerta:}
\If{$d_{\text{compuerta}} = 1$}
    \State Transmitir y escalar: $\mathbf{Y}' \gets \mathbf{Y} \cdot \mathbf{G}$ \Comment{Escalado por cabeza}
\Else
    \Comment{$d_{\text{compuerta}} = d_k$}
    \State Modulación elemento a elemento: $\mathbf{Y}' \gets \mathbf{Y} \odot \mathbf{G}$ \Comment{Compuerta por dimensión}
\EndIf
\Return $\mathbf{Y}'$
\end{algorithmic}
\end{algorithm}

\subsubsection{Hallazgos Clave}

\begin{itemize}[leftmargin=*]
    \item \textbf{Supresión del sumidero de atención}: La compuerta rompe el ciclo de retroalimentación donde softmax se concentra en el primer token, permitiendo una atención más equilibrada.
    \item \textbf{Estabilidad de entrenamiento}: Los modelos con compuerta toleran tasas de aprendizaje más grandes ($+30\%$ mayor $\eta_{max}$ estable).
    \item \textbf{Sobrecarga mínima}: El impacto en latencia es solo del $1.6\%$ en GPUs H100; el rendimiento cae como máximo $2-3\%$.
    \item \textbf{Mejora en la ley de escalado}: Mejora la pérdida en todas las escalas de modelo, desde 800M hasta 110B parámetros de forma consistente.
\end{itemize}

\normalsize

\subsection{7. Gated DeltaNet-2: Borrado y Escritura Canalizados Decouplados}

\subsubsection{Fundamento Matemático}

Gated DeltaNet-2 \citep{hatamizadeh_gated_deltanet2_2026} aborda una limitación compartida de Gated DeltaNet y Kimi Delta Attention (KDA): ambos usan una única \emph{compuerta delta escalar} $\beta_t$ para controlar simultáneamente (i) cuánto contenido antiguo se borra en la dirección de lectura del eje de claves y (ii) cuánto valor nuevo se escribe en el eje de valores. Gated DeltaNet-2 \textbf{desacopla} estos dos roles en una compuerta de borrado canalizada $\mathbf{b}_t$ (eje de claves) y una compuerta de escritura canalizada $\mathbf{w}_t$ (eje de valores), heredando además el decaimiento canalizado $\boldsymbol{\alpha}_t$ de KDA. Con estado $\mathbf{S}_t \in \mathbb{R}^{d_k \times d_v}$, decaimiento $\mathbf{D}_t = \mathrm{Diag}(\boldsymbol{\alpha}_t)$, dirección de borrado $\mathbf{e}_t = \mathbf{b}_t \odot \mathbf{k}_t$ y objetivo de escritura $\mathbf{z}_t = \mathbf{w}_t \odot \mathbf{v}_t$, la actualización de estado es:
\[
\mathbf{S}_t = \left(\mathbf{I} - \mathbf{k}_t \mathbf{e}_t^\top\right) \mathbf{D}_t \, \mathbf{S}_{t-1} + \mathbf{k}_t \mathbf{z}_t^\top, \qquad \mathbf{o}_t = \mathbf{S}_t^\top \mathbf{q}_t
\]
o, expandiendo los productos de las compuertas:
\[
\mathbf{S}_t = \left(\mathbf{I} - \mathbf{k}_t (\mathbf{b}_t \odot \mathbf{k}_t)^\top\right) \mathbf{D}_t \, \mathbf{S}_{t-1} + \mathbf{k}_t (\mathbf{w}_t \odot \mathbf{v}_t)^\top
\]
Las dos compuertas se producen mediante proyecciones lineales independientes seguidas de sigmoide:
\[
\mathbf{b}_t = \sigma(\mathbf{W}_b \mathbf{x}_t), \qquad \mathbf{w}_t = \sigma(\mathbf{W}_w \mathbf{x}_t)
\]
y el decaimiento canalizado sigue una parametrización de log-decaimiento (calculada en fp32 para evitar pérdida de precisión en productos acumulativos largos):
\[
\mathbf{g}_t = \mathbf{a} - \mathrm{softplus}(\boldsymbol{\delta}), \qquad \boldsymbol{\alpha}_t = \exp(\mathbf{g}_t)
\]
Los roles estructurales de las dos compuertas son complementarios:
\begin{itemize}[leftmargin=*]
    \item $\mathbf{b}_t$ (eje de claves) hace que la \textbf{dirección de lectura} sea selectiva por canal: determina qué canales de clave se eliminan antes de escribir.
    \item $\mathbf{w}_t$ (eje de valores) hace que el \textbf{objetivo de escritura} sea selectivo por canal: determina qué canales de valor se depositan en memoria.
    \item $\boldsymbol{\alpha}_t$ (por canal de clave) proporciona decaimiento global, como en KDA.
\end{itemize}
\textbf{Reducciones.} Fijar $\mathbf{b}_t = \beta_t \mathbf{1}_{d_k}$ y $\mathbf{w}_t = \beta_t \mathbf{1}_{d_v}$ recupera KDA; colapsar además $\boldsymbol{\alpha}_t$ a un escalar recupera Gated DeltaNet. Desde la perspectiva de aprendizaje en línea/pesos rápidos, la actualización es el minimizador de:
\[
\mathcal{L}_t(\mathbf{S}) = \tfrac{1}{2}\|\mathbf{S} - \mathbf{D}_t \mathbf{S}_{t-1}\|_F^2 - 2\left\langle \mathbf{S}^\top \mathbf{k}_t,\; \mathbf{z}_t - (\mathbf{D}_t \mathbf{S}_{t-1})^\top \mathbf{e}_t \right\rangle
\]

\begin{algorithm}[H]
\caption{Avance Recurrente de Gated DeltaNet-2 (Forma de Referencia)}
\begin{algorithmic}[1]
\Require Entrada $\mathbf{x}_t$; estado previo $\mathbf{S}_{t-1} \in \mathbb{R}^{d_k \times d_v}$
\Require Proyecciones: $W_q, W_k, W_v, W_b, W_w$; parámetros de log-decaimiento $\mathbf{a}, \boldsymbol{\delta}$
\Ensure Salida $\mathbf{o}_t$ y estado $\mathbf{S}_t$
\State $\mathbf{q}_t \gets \mathrm{normalize}(W_q \mathbf{x}_t)$, $\mathbf{k}_t \gets \mathrm{normalize}(W_k \mathbf{x}_t)$, $\mathbf{v}_t \gets \mathrm{silu}(W_v \mathbf{x}_t)$
\State $\mathbf{b}_t \gets \sigma(W_b \mathbf{x}_t)$ \Comment{compuerta de borrado eje-clave, $\mathbb{R}^{d_k}$}
\State $\mathbf{w}_t \gets \sigma(W_w \mathbf{x}_t)$ \Comment{compuerta de escritura eje-valor, $\mathbb{R}^{d_v}$}
\State $\boldsymbol{\alpha}_t \gets \exp(\mathbf{a} - \mathrm{softplus}(\boldsymbol{\delta}))$ \Comment{decaimiento canalizado, fp32}
\State Dirección de borrado: $\mathbf{e}_t \gets \mathbf{b}_t \odot \mathbf{k}_t$; objetivo de escritura: $\mathbf{z}_t \gets \mathbf{w}_t \odot \mathbf{v}_t$
\State Decaer estado: $\mathbf{S}_t \gets \mathrm{Diag}(\boldsymbol{\alpha}_t)\, \mathbf{S}_{t-1}$
\State Lectura: $\mathbf{r}_t \gets \mathbf{S}_t^\top \mathbf{e}_t$ \Comment{suma sobre eje de claves}
\State Actualización delta: $\mathbf{S}_t \gets \mathbf{S}_t - \mathbf{k}_t \mathbf{r}_t^\top + \mathbf{k}_t \mathbf{z}_t^\top$
\State Recuperar: $\mathbf{o}_t \gets \mathbf{S}_t^\top \mathbf{q}_t$
\State Compuerta de salida: $\mathbf{o}_t \gets \mathrm{GroupNorm}(\mathbf{o}_t) \odot \mathrm{silu}(W_g \mathbf{x}_t)$
\State \Return $\mathbf{o}_t$, $\mathbf{S}_t$
\end{algorithmic}
\end{algorithm}

El algoritmo de entrenamiento por fragmentos extiende la representación WY de KDA absorbiendo el decaimiento canalizado en una recurrencia delta asimétrica sobre un estado normalizado por decaimiento, permitiendo paralelismo rico en matmuls en tensor cores. Como $\mathbf{b}_t$ y $\mathbf{w}_t$ son operadores diagonales por canal que varían por fila, el atajo de post-escalado escalar de KDA se rompe: el paso hacia atrás requiere una \textbf{inversión WY consciente de las compuertas} que incorpora las compuertas en cada producto escalar que acumula el gradiente.

\subsubsection{Rendimiento Empírico y Compromisos}

A 1.3B parámetros entrenados con 100B tokens FineWeb-Edu, Gated DeltaNet-2 supera a Mamba-2, Mamba-3, Gated DeltaNet y KDA en modelado de lenguaje, razonamiento de sentido común y---más notablemente---recuperación de contexto largo del mundo real y RULER (especialmente multi-key needle-in-a-haystack). Su rendimiento en H100 se mantiene casi plano ($38.0 \to 36.1$ Kt/s) al crecer la longitud de secuencia de 2K a 16K, donde el de un Transformer cae bruscamente. Los estudios de ablación muestran que la compuerta de borrado $\mathbf{b}_t$ explica la mayor parte de la ganancia; la compuerta de escritura $\mathbf{w}_t$ aporta una mejora adicional menor. La variante híbrida (Gated DeltaNet-2 recurrente + atención de ventana deslizante) reduce aún más la brecha con los modelos de atención pura en tareas que requieren agregación de evidencia local.

\begin{table}[H]
\centering
\footnotesize
\begin{tabularx}{\linewidth}{lXX}
\toprule
\textbf{Aspecto} & \textbf{Fortaleza} & \textbf{Limitación} \\
\midrule
Recuperación & La mejor entre modelos delta recurrentes en multi-key NIAH. & La variante recurrente aún supera al Transformer en NQ/DROP. \\
Rendimiento & Escalado casi plano 2K$\to$16K; pequeña sobrecarga vs KDA. & Requiere un nuevo kernel Triton WY consciente de las compuertas. \\
Memoria & Estado de inferencia constante $\mathcal{O}(d^2)$. & Rango de borrado de autovalores negativos $[0,2]$ no da ganancia consistente a 1.3B. \\
\bottomrule
\end{tabularx}
\caption{Compromisos de Gated DeltaNet-2 a escala 1.3B / 100B tokens \citep{hatamizadeh_gated_deltanet2_2026}.}
\label{tab:gated-deltanet2-tradeoffs-es}
\end{table}

\normalsize

\subsection{8. Kimi Delta Attention (KDA): Compuerta de Decaimiento Canalizada para la Regla Delta}

\subsubsection{Fundamento Matemático}

Kimi Delta Attention \citep{kimi_kda_2025} es el módulo central de Kimi Linear y extiende Gated DeltaNet con una \textbf{compuerta de decaimiento canalizada} de grano más fino. Mientras que Gated DeltaNet usa un único decaimiento escalar por cabeza $\alpha_t$ junto con la compuerta escalar de escritura delta $\beta_t$, KDA reemplaza el decaimiento escalar con un decaimiento \textbf{por canal de clave} $\boldsymbol{\alpha}_t \in \mathbb{R}^{d_k}$, parametrizado mediante una formulación de log-decaimiento y aplicado como una matriz diagonal $\mathbf{D}_t = \mathrm{diag}(\boldsymbol{\alpha}_t)$ antes de la actualización de la regla delta. La regla delta en sí retiene una compuerta de escritura escalar por cabeza $\beta_t$:
\[
\mathbf{S}_t = (\mathbf{I} - \beta_t \mathbf{k}_t \mathbf{k}_t^\top)\, \mathbf{D}_t\, \mathbf{S}_{t-1} + \beta_t \mathbf{v}_t \mathbf{k}_t^\top, \qquad \mathbf{o}_t = \mathbf{S}_t^\top \mathbf{q}_t
\]
El decaimiento canalizado sigue una parametrización de log-decaimiento (calculada en fp32 por estabilidad numérica en productos acumulativos largos):
\[
\mathbf{g} = \mathbf{a} - \mathrm{softplus}(\boldsymbol{\delta}), \qquad \boldsymbol{\alpha}_t = \exp(\mathbf{g})
\]
donde $\mathbf{a}$ y $\boldsymbol{\delta}$ son parámetros aprendidos por canal. Esto da a cada canal de clave una escala temporal de olvido independiente, permitiendo al modelo retener algunos canales más tiempo que otros dentro de la misma cabeza.

\paragraph{Reducción a Gated DeltaNet}: Cuando $\boldsymbol{\alpha}_t$ colapsa a un escalar (es decir, todos los canales comparten el mismo decaimiento), KDA se reduce exactamente a Gated DeltaNet. El decaimiento canalizado es por tanto una generalización estricta que añade selectividad por canal sobre el decaimiento/escritura escalares de Gated DeltaNet.

\paragraph{Entrenamiento por fragmentos}: Un algoritmo por fragmentos ad hoc usa una variante \textbf{Diagonal-Plus-Low-Rank (DPLR)} especializada de la matriz de transición para mantener la eficiencia en hardware. El producto de las transiciones por paso $(\mathbf{I} - \beta_t \mathbf{k}_t \mathbf{k}_t^\top)\mathbf{D}_t admite una descomposición DPLR que permite paralelismo rico en matmuls en tensor cores preservando la estructura de decaimiento canalizado.

\begin{algorithm}[H]
\caption{Avance Recurrente de Kimi Delta Attention (Forma de Referencia)}
\begin{algorithmic}[1]
\Require Entrada $\mathbf{x}_t$; estado previo $\mathbf{S}_{t-1} \in \mathbb{R}^{d_k \times d_v}$
\Require Proyecciones: $W_q, W_k, W_v, W_\beta$; parámetros de log-decaimiento $\mathbf{a}, \boldsymbol{\delta}$
\Ensure Salida $\mathbf{o}_t$ y estado $\mathbf{S}_t$
\State $\mathbf{q}_t \gets W_q \mathbf{x}_t$, $\mathbf{k}_t \gets W_k \mathbf{x}_t$, $\mathbf{v}_t \gets W_v \mathbf{x}_t$
\State $\beta_t \gets \mathrm{sigmoid}(W_\beta \mathbf{x}_t)$ \Comment{compuerta de escritura escalar por cabeza}
\State $\boldsymbol{\alpha}_t \gets \exp(\mathbf{a} - \mathrm{softplus}(\boldsymbol{\delta}))$ \Comment{decaimiento canalizado, fp32}
\State $\mathbf{D}_t \gets \mathrm{diag}(\boldsymbol{\alpha}_t)$ \Comment{matriz diagonal de decaimiento}
\State Decaer estado: $\mathbf{S}_t \gets \mathbf{D}_t \mathbf{S}_{t-1}$
\State Lectura: $\mathbf{r}_t \gets \mathbf{S}_t^\top \mathbf{k}_t$ \Comment{predicción actual a lo largo de la clave}
\State Actualización delta: $\mathbf{S}_t \gets \mathbf{S}_t - \beta_t \mathbf{k}_t \mathbf{r}_t^\top + \beta_t \mathbf{v}_t \mathbf{k}_t^\top$
\State Recuperar: $\mathbf{o}_t \gets \mathbf{S}_t^\top \mathbf{q}_t$
\State Compuerta de salida: $\mathbf{o}_t \gets \mathrm{GroupNorm}(\mathbf{o}_t) \odot \mathrm{silu}(W_g \mathbf{x}_t)$
\State \Return $\mathbf{o}_t$, $\mathbf{S}_t$
\end{algorithmic}
\end{algorithm}

\subsubsection{Rendimiento Empírico y Compromisos}

KDA está desplegado en Kimi Linear, un modelo híbrido (3B activados / 48B parámetros totales) que combina capas KDA con Multi-head Latent Attention (MLA). Kimi Linear supera a un baseline de MLA completa en todas las tareas evaluadas, reduce el caché KV hasta un $75\%$, y logra hasta $6\times$ de throughput de decodificación para contexto de 1M de tokens. El decaimiento canalizado es el principal impulsor de la ganancia: permite control fino sobre qué canales de clave persisten, mejorando el recuerdo de contexto largo frente a Gated DeltaNet con decaimiento escalar.

\begin{table}[H]
\centering
\footnotesize
\begin{tabularx}{\linewidth}{lXX}
\toprule
\textbf{Aspecto} & \textbf{Fortaleza} & \textbf{Limitación} \\
\midrule
Recuperación & El decaimiento canalizado mejora el recuerdo de contexto largo frente a baselines de decaimiento escalar; desplegado en Kimi Linear (híbrido 3B/48B). & El estado recurrente matricial $\mathcal{O}(d^2)$ es más pesado que la atención lineal con compuerta diagonal. \\
Rendimiento & Hasta $6\times$ de throughput de decodificación a 1M de contexto; $75\%$ de reducción de caché KV frente a MLA completa. & Requiere kernels por fragmentos conscientes de DPLR para eficiencia en hardware. \\
Memoria & Estado de inferencia constante $\mathcal{O}(d^2)$ (recurrente matricial). & El cálculo log-decaimiento en fp32 añade sobrecarga; la reducción a Gated DeltaNet pierde selectividad por canal. \\
\bottomrule
\end{tabularx}
\caption{Compromisos de KDA en Kimi Linear (3B activados / 48B total) \citep{kimi_kda_2025}.}
\label{tab:kda-tradeoffs-es}
\end{table}

\normalsize

\subsection{Análisis Comparativo: Taxonomía de Mecanismos con Compuerta}

\footnotesize
\begin{longtable}{p{1.8cm}p{1.7cm}p{1.8cm}p{1.6cm}p{2.0cm}p{2.4cm}p{1.3cm}}
\toprule
\textbf{Arquitectura} & \textbf{Año de Publicación} & \textbf{Representación de Estado} & \textbf{Tipo de Compuerta} & \textbf{Complejidad de Entrenamiento} & \textbf{Complejidad de Inferencia} & \textbf{Ref} \\
\midrule
\endhead
GLA & 2023 & Matriz (recurrente) & Decaimiento diagonal & $\mathcal{O}(Ld^2)$ & $\mathcal{O}(d^2)$ memoria & \citep{yang_gla_2023} \\
DeltaNet & 2024 & Matriz (recurrente) & Escritura escalar ($\beta$) & $\mathcal{O}(Ld^2)$ & $\mathcal{O}(d^2)$ memoria & \citep{yang_deltanet_2024} \\
Gated DeltaNet & 2024 & Matriz (recurrente) & Decaimiento + escritura & $\mathcal{O}(Ld^2)$ & $\mathcal{O}(d^2)$ memoria & \citep{yang_gated_deltanet_2024} \\
RetNet & 2023 & Matriz (recurrente) & Exponencial fija & $\mathcal{O}(Ld^2)$ & $\mathcal{O}(d)$ por paso & \citep{sun_retentive_2023} \\
HGRN2 & 2024 & Matriz (recurrente) & Diagonal (acotada) & $\mathcal{O}(Ld^2)$ & $\mathcal{O}(d^2)$ memoria & \citep{qin_hgrn2_2024} \\
FoX & 2025 & Atención completa & Sesgo en espacio de logits & $\mathcal{O}(L^2d)$ & $\mathcal{O}(L)$ caché KV & \citep{lin_forgetting_transformer_2025} \\
Gated Softmax & 2025 & Atención completa & Sigmoide post-SDPA & $\mathcal{O}(L^2d)$ & $\mathcal{O}(L)$ caché KV & \citep{qiu_gated_attention_2025} \\
Gated DeltaNet-2 & 2026 & Matriz (recurrente) & Borrado + escritura decouplados (canalizado) & $\mathcal{O}(Ld^2)$ & $\mathcal{O}(d^2)$ memoria & \citep{hatamizadeh_gated_deltanet2_2026} \\
KDA & 2025 & Matriz (recurrente) & Decaimiento canalizado + escritura escalar & $\mathcal{O}(Ld^2)$ & $\mathcal{O}(d^2)$ memoria & \citep{kimi_kda_2025} \\
\bottomrule
\end{longtable}
\normalsize

\medskip

\footnotesize
\begin{longtable}{p{1.8cm}p{2.0cm}p{2.2cm}p{2.2cm}p{2.2cm}p{1.8cm}}
\toprule
\textbf{Arquitectura} & \textbf{Entrenable} & \textbf{Recuperación en Contexto} & \textbf{Recuperación Asociativa} & \textbf{Extrapolación de Longitud} & \textbf{Ventaja Clave} \\
\midrule
\endhead
GLA & Sí & Moderada & Débil & Buena (20K) & Eficiencia en hardware; paralelismo por fragmentos \\
DeltaNet & Sí & Fuerte & Perfecta (MQAR) & Buena & Semántica de corrección de errores; recuperación sólida \\
Gated DeltaNet & Sí & Muy fuerte & Perfecta & Excelente & Equilibrada: olvido + actualizaciones dirigidas \\
RetNet & Sí & Débil & Débil & Buena & Simplicidad; no necesita caché KV \\
HGRN2 & Sí & Moderada & Débil & Moderada & Modelado temporal multi-escala \\
FoX & Sí & Muy fuerte & Muy fuerte (casi perfecta) & Excelente & Preserva softmax; sobrecarga mínima \\
Gated Softmax & Sí & Fuerte & Buena & Buena & Simplicidad; no reemplaza SDPA \\
Gated DeltaNet-2 & Sí & Muy fuerte & Muy fuerte & Excelente & Borrado/escritura canalizados decouplados; mejor RULER multi-key \\
KDA & Sí & Muy fuerte & Muy fuerte & Excelente & Decaimiento canalizado; desplegado en Kimi Linear (75\% reducción KV, $6\times$ decodificación) \\
\bottomrule
\end{longtable}
\normalsize

\medskip

\footnotesize
\begin{longtable}{p{2.0cm}p{2.4cm}p{2.8cm}p{3.4cm}}
\toprule
\textbf{Arquitectura} & \textbf{Rendimiento Típico} & \textbf{Memoria por Inferencia} & \textbf{Caso de Uso Principal} \\
\midrule
\endhead
GLA & Alto (CUDA por fragmentos) & $O(d^2)$ constante & Contexto largo con restricciones de memoria \\
DeltaNet & Medio-Alto & $O(d^2)$ constante & Recuperación y razonamiento asociativo \\
Gated DeltaNet & Medio & $O(d^2)$ constante & Producción: Qwen3-Next, rendimiento equilibrado \\
RetNet & Alto & $O(d)$ por paso & Inferencia extremadamente rápida; modelos simples \\
HGRN2 & Medio-Alto & $O(d^2)$ constante & Patrones temporales multi-escala \\
FoX & Moderado (cuadrático) & $O(Ld)$ caché KV & Extrapolación de longitud; generación de contexto largo \\
Gated Softmax & Alto (sobrecarga mínima) & $O(Ld)$ caché KV & Mejora directa para modelos existentes \\
Gated DeltaNet-2 & Alto (escalado casi plano) & $O(d^2)$ constante & Recuperación de contexto largo; multi-key NIAH \\
KDA & Alto (hasta $6\times$ decodificación a 1M) & $O(d^2)$ constante & Producción híbrida (Kimi Linear); decodificación de contexto largo \\
\bottomrule
\end{longtable}
\normalsize

\subsection{Principio Unificado de Compuerta}

Las ocho arquitecturas comparten un principio común: \textbf{la compuerta es un mecanismo para controlar el flujo de información y la retención selectiva de memoria}. Las decisiones de diseño específicas divergen a lo largo de varias dimensiones:

\begin{enumerate}[leftmargin=*]
    \item \textbf{Ubicación de la compuerta}: Actualizaciones de estado recurrente (GLA, DeltaNet, Gated DeltaNet, Gated DeltaNet-2, HGRN2) versus espacio de logits/salida softmax (FoX, Gated Softmax, RetNet).
    \item \textbf{Granularidad de control}: Por dimensión (GLA, HGRN2 diagonal), por clave (DeltaNet), por token (FoX), por cabeza (Gated Softmax) o borrado/escritura canalizados decouplados (Gated DeltaNet-2).
    \item \textbf{Dependencia de datos}: Completamente dependiente de datos (GLA, DeltaNet, Gated DeltaNet, Gated DeltaNet-2, FoX, Gated Softmax) versus esquemas fijos (RetNet con decaimiento exponencial).
    \item \textbf{Compensación de expresividad}: Eficiencia recurrente lineal (GLA, DeltaNet, Gated DeltaNet, Gated DeltaNet-2, HGRN2, RetNet) versus expresividad completa de softmax (FoX, Gated Softmax).
\end{enumerate}

\subsection{Implementación y Orientación Práctica}

\subsubsection{Cuándo Usar Cada Arquitectura}

\begin{enumerate}[leftmargin=*]
    \item \textbf{GLA}: Inferencia rápida con contextos moderados a largos (2K--20K tokens), cuando los kernels CUDA personalizados son aceptables y el rendimiento de recuperación no es crítico.
    \item \textbf{DeltaNet}: Fuerte recuperación asociativa y tareas de aprendizaje en contexto; bueno para tareas sintéticas (MQAR) y pruebas de asociación clave-valor.
    \item \textbf{Gated DeltaNet}: Modelos de producción que requieren el mejor equilibrio de recuperación, olvido y rendimiento (probado en Qwen3-Next; ICLR 2025).
    \item \textbf{Gated DeltaNet-2}: Cargas de recuperación de contexto largo (multi-key NIAH) donde el borrado/escritura canalizado decouplado da el recuerdo más fuerte; cuando se dispone de un kernel WY consciente de las compuertas.
    \item \textbf{RetNet}: Inferencia extremadamente eficiente sin almacenamiento en caché KV; adecuado para despliegues en dispositivos móviles/de borde o modelos de juguete.
    \item \textbf{HGRN2}: Jerarquías temporales multi-escala; cuando los límites de olvido específicos por capa son deseables.
    \item \textbf{FoX}: Extrapolación de longitud y comprensión de contexto largo; cuando el entrenamiento cuadrático es aceptable y no se desea reemplazar softmax.
    \item \textbf{Gated Softmax}: Mejora rápida para modelos softmax existentes con cambios mínimos de código; mejor para profesionales que desean ganancias inmediatas sin una reestructuración importante.
\end{enumerate}

\subsubsection{Consideraciones de Hardware}

\begin{table}[H]
\centering
\small
\begin{tabularx}{0.95\linewidth}{Xp{3.5cm}p{3.5cm}}
\toprule
\textbf{Objetivo de Hardware} & \textbf{Arquitecturas Recomendadas} & \textbf{Notas} \\
\midrule
GPU (H100/A100) & Gated Softmax, FoX, Gated DeltaNet, Gated DeltaNet-2 & Implementaciones maduras de softmax/lineal. GLA/DeltaNet requieren kernels personalizados. \\
TPU (alta memoria) & GLA, DeltaNet, Gated DeltaNet, Gated DeltaNet-2, HGRN2 & Hardware soporta matmuls eficientes; los métodos recurrentes lineales brillan. \\
Móvil/Borde & RetNet, Gated Softmax ligero & Inferencia con memoria constante es crítica. \\
\bottomrule
\end{tabularx}
\end{table}


\input{appendices/annex-e-tutorial}

\end{document}
